\section{Des documents en informatique appliquée}

\paragraphe{P003}{De cette revue de littérature, il ressort que les objectifs principaux des documents en informatique appliquée sont d'aider à réaliser et à établir des entités informatiques, et plus couramment des logiciels. Pendant que la documentation est créée et utilisée à ces fins, plusieurs écueils se dressent. La majorité de ceux-ci proviennent de la façon d'œuvrer en informatique. Pour pallier ces écueils, il existe des approches constituées de méthodes et de techniques, mais ces dernières ne résolvent pas tout.}{}{}

\paragraphe{P004}{Au préalable, il est nécessaire de donner ces définitions pour éviter des ambiguïtés~: 
\begin{itemize}
    \item logiciel~: ensemble de suites d’instructions interprétables par un ordinateur;
    \item composant logiciel~: logiciel ou partie d’un logiciel\footnote{Le logiciel a de particulier qu'il peut comprendre lui-même d'autres logiciels.};
    \item entité informatique~: système, sous-système, composant logiciel ou composant matériel en lien avec l'informatique;
    \item document~: écrit servant d’information;
    \item documentation informatique~: ensemble de documents associé à un ensemble d’entités informatiques.
\end{itemize}
\vspace{-1\baselineskip}
}{D01, D02, D03, D04, D05}{D01}

\enonce{D01}{Définition}{Logiciel (entité)}{
Ensemble de suites d'instructions interprétable par un ordinateur.
}{P004}
\explication{D01}{de la définition}{Logiciel (entité)}{
La définition formulée est~: ensemble de suites d’instructions interprétables par un ordinateur.

La littérature fournit ces informations qui ont servi à construire cette définition~:
\begin{itemize}
    \item définition de programme~: «~Suite d'instructions écrites sous une forme que l'ordinateur peut comprendre pour traiter un problème ou pour effectuer une tâche.~»
    \cite{Gdt_programme_0};
    \item définition 1 de logiciel~: «~Ensemble des programmes constituant une unité destinée à effectuer un traitement particulier sur un ordinateur.~»
    \cite{Gdt_logiciel_0};
    \item définition 2 de logiciel~: \enfr
    {all or part of the programs, procedures, rules, and associated documentation of an information processing system}
    {tout ou partie des programmes, procédures, règles et de la documentation associée d'un système de traitement de l'information}
    \cite{iso_24765_2017}.
\end{itemize}
Un logiciel est composé de programmes et compose une ou plusieurs entités informatiques. Lorsqu'il est employé seul, il désigne une entité logicielle.
}{}

\enonce{D02}{Définition}{Composant logiciel}{
Logiciel ou partie d'un logiciel.
}{P004}
\explication{D02}{de la définition}{Composant logiciel}{
La définition formulée est~: logiciel ou partie d'un logiciel.

La littérature fournit ces informations qui ont servi à construire cette définition~:
\begin{itemize}
    \item définition de composant~: «~Objet, substance, élément qui entre dans la composition de quelque chose \elide.~»
    \cite{Larousse_composant_0};
    \item définition de composant logiciel~: 
    «~Logiciel, programme ou élément d'un logiciel ou d'un programme qui constitue un module indépendant utilisé comme élément d'un système plus complexe et qui est spécialement conçu pour fonctionner sans problèmes avec d'autres logiciels ou programmes.~»
    \cite{Gdt_composantlogiciel_0};
    \item description de composant~:  \enfr
    {A component can be hardware or software and can be subdivided into other components. Component refers to a part of a whole, such as a component of a software product or a component of a software identification tag. \elide}
    {Un composant peut être matériel ou logiciel et peut être subdivisé en d’autres composants. Un composant désigne une partie d’un tout, comme un composant d’un produit logiciel ou un composant d’une étiquette d’identification logicielle.} 
    \ccite{iso_24765_2017}[p.~82].
\end{itemize}
}{}

\enonce{D03}{Définition}{Entité informatique}{
Système, sous-système, composant logiciel ou composant matériel en lien avec l'informatique.
}{P004}
\explication{D03}{de la définition}{Entité informatique}{
La définition formulée est~: Système, sous-système, composant logiciel ou composant matériel en lien avec l'informatique.

La littérature fournit ces informations qui ont servi à construire cette définition~:
\begin{itemize}
    \item définition d'entité~: «~Chose considérée comme un être ayant son individualité.~»
    \cite{Gdt_entite_0};
    \item définition d'informatique~: «~Discipline qui s'intéresse à tous les aspects, tant théoriques que pratiques, reliés au traitement automatique de l'information, à la conception, à la programmation, au fonctionnement et à l'utilisation des ordinateurs.~» \cite{Gdt_informatique_0};
    \item définition de système~: \enfr
    {combination of interacting elements organized to achieve one or more stated purposes}
    {combinaison d'éléments interagissant organisés pour atteindre un ou plusieurs objectifs énoncés}
    \cite{iso_24765_2017};
    \item définition 1 de produit~: \enfr
    {an artifact that is produced, is quantifiable, and can be either an end item in itself or a component item}
    {un artefact produit, quantifiable et qui peut être soit un produit final en soi, soit un composant}
    \cite{iso_24765_2017};
    \item définition 2 de produit~: \enfr
    {complete set of computer programs, procedures, and associated documentation designed for delivery to a user}
    {ensemble complet de programmes informatiques, de procédures et de documentation associée conçu pour être livré à un utilisateur}
    \cite{iso_24765_2017};
    \item définition 3 de produit~: \enfr
    {result of a process}
    {résultat d'un processus}
    \cite{iso_24765_2017}.
\end{itemize}
Puisque les multiples définitions de \emphase{produit} peuvent porter à confusion, l'utilisation du terme \emphase{produit} est réservée comme étant le résultat d'un processus, et, dans ce mémoire, ce résultat est précisément désigné par les termes  \emphase{entité informatique}, \emphase{composant logiciel} ou \emphase{logiciel}. Plus particulièrement, le produit qui importe le plus dans ce mémoire est le logiciel.
}{}

\enonce{D04}{Définition}{Document}{
Écrit servant d'information.
}{P004}
\explication{D04}{de la définition}{Document}{
La définition formulée est~: écrit servant d'information.

La littérature fournit ces informations qui ont servi à construire cette définition~:
\begin{itemize}
\item définition 1~: «~pièce écrite servant d'information, de preuve~»
\cite{Larousse_document_0};
\item définition 2~: «~objet quelconque servant de preuve, de témoignage~»
\cite{Larousse_document_0};
\item définition 3~: \enfr
{uniquely identified unit of information for human use, such as a report, specification, manual or book, in printed  or electronic form}
{unité d'information identifiée de manière unique et destinée à l'usage humain, telle qu'un rapport, un cahier des charges, un manuel ou un livre, sous forme imprimée ou électronique}
\cite{iso_24765_2017};
\item définition 4~: \enfr
{medium, and the information recorded on it, that generally has permanence and can be read by a person or a machine.}
{support ainsi que l'information qui y est enregistrée, qui a généralement une permanence et peut être lue par une personne ou une machine.}
\cite{iso_24765_2017}.
\end{itemize}
}{}

\enonce{D05}{Définition}{Documentation informatique}{
Ensemble de documents associé à un ensemble d'entités informatiques.
}{P004}
\explication{D05}{de la définition}{Documentation informatique}{
La définition formulée est~: ensemble de documents associé à un ensemble d'entités informatiques.

La littérature fournit ces informations qui ont servi à construire cette définition~:
\begin{itemize}
    \item définition 1 de documentation informatique~: «~Ensemble de documents accompagnant et décrivant un produit informatique.~» \cite{Gdt_documentation_0} ;
    \item note jointe à cette définition~: «~La documentation sert à faciliter l'installation, la compréhension et la maintenance d'un produit informatique.~» \cite{Gdt_documentation_0} ;
    \item définition 2 de documentation informatique~: \enfr 
    {written or pictorial information describing, defining, specifying, reporting, or certifying activities, requirements, procedures, or results}
    {informations écrites ou illustrées décrivant, définissant, spécifiant, rapportant ou certifiant des activités, des exigences, des procédures ou des résultats}
    \cite{iso_24765_2017};
    \item avis de Paul Clements~: \enfr
    {Documentation can be formal or not, as appropriate, and may contain models or not, as appropriate. Documents may describe, or they may specify. Hence, the term is appropriately general.}
    {La documentation peut être formelle ou informelle, selon le contexte, et peut contenir des modèles ou non, selon le contexte. Les documents peuvent décrire ou spécifier. Par conséquent, le terme est suffisamment général.}    
    \ccite{clements_documenting_2014}[p.~12].
\end{itemize}
La documentation logicielle est indissociable du produit informatique, c'est pour cette raison que le terme \emphase{documentation informatique} est aussi répandu. 
}{}

\subsection{Des objectifs}

\paragraphe{P005}{Les objectifs de la documentation sont déterminés en considérant l'aspect fondamental et l'aspect pratique. L'aspect fondamental est principalement déterminé par l'association entre le document et l'entité informatique en elle-même. L'aspect pratique est déterminé principalement par l'usage des documents et leurs catégorisations. Ce faisant, il est également montré que les documents dans leur ensemble représentent fréquemment le seul véhicule pour objectiver l'entité informatique. Ainsi, ces documents se révèlent nécessaires à la réalisation et à l'établissement du logiciel.}{}{}

\paragraphe{P006}{\fl{L’association fondamentale est que le document sert\footnote{Il peut être utilisé, cité ou requis. Aussi, il est possible qu'un seul document puisse servir à plusieurs entités informatiques.} à une entité informatique.}  Les documents et les entités informatiques font partie des artefacts du procédé de développement. Voici des définitions~:
\begin{itemize}
    \item procédé~: ensemble structuré d'activités visant un but logiquement déterminé par des antécédents et des conséquents où des intrants donnent lieu à des extrants;
    \item procédé de développement~: procédé dont le but est de créer des entités informatiques;
    \item processus~: instance d'un procédé;
    \item artefact~: intrant tangible ou extrant tangible à un processus;
    \item phase d'un projet~: regroupement d’activités logiquement interreliées découlant de la gestion d'un projet;
    \item phase d'un procédé de développement~: regroupement d’activités logiquement interreliées découlant d'un procédé de développement.
\end{itemize}
Le processus va créer un ou plusieurs artefacts. Cependant, le logiciel demeure le but du procédé de développement, les autres artefacts, comme les documents, servent assurément ce but, et il sera décrit la manière. Dans le cadre de ce mémoire, l'utilisation du terme \emphase{phase} porte sur les procédés de développement. 
}{OD01, OD02, OD03, OD04, OD05, OD06}{OD01}

\enonce{OD01}{Définition}{Procédé}{
Ensemble structuré d'activités visant un but logiquement déterminé par des antécédents et des conséquents où des intrants donnent lieu à des extrants.
}{P006}
\explication{OD01}{de la définition}{Procédé}{
La définition formulée est~: ensemble structuré d'activités visant un but logiquement déterminé par des antécédents et des conséquents où des intrants donnent lieu à des extrants. La littérature fournit ces informations qui ont servi à construire cette définition~:
\begin{itemize}
\item définition 1 de procédé~: «~une suite continue de faits, de phénomènes présentant une certaine unité ou une certaine régularité dans leur déroulement~» \cite{Bdl_procede_0};
\item définition 2 de procédé~: «~une méthode utilisée en vue d’obtenir un résultat déterminé.~» \cite{Bdl_procede_0};
\item définition 3 de procédé~: «~\elide un procédé est une activité humaine ayant des éléments d'entrées (matières premières ou personnes) et des éléments de sorties (produits finis ou personnes). Il y a donc bien transformation d'objets ayant certaines caractéristiques en objets en possédant d'autres.~» \footnote{Wikipedia, «~Procédé~», \url{https://fr.wikipedia.org/wiki/Proc\%C3\%A9d\%C3\%A9} (consulté le 2026-05-15).};
\item définition de méthode~: «~Ensemble de démarches que suit l'esprit pour découvrir et démontrer la vérité.~»  \cite{Robert_methode_0};
\item définition 1 de processus~: «~Suite continue de faits ou de phénomènes qui présentent une certaine unité ou une certaine régularité dans leur déroulement.~» \cite{Gdt_processus_1} ;
\item définition 2 de processus~: «~Ensemble d'activités logiquement interreliées qui produisent un résultat déterminé.~» \cite{Gdt_processus_0}.
\end{itemize}
Le procédé et le processus peuvent être confondus. Il est retenu que le procédé a un objectif à remplir tandis que le processus se déroule. 
}{}

\enonce{OD02}{Définition}{Procédé de développement}{
Procédé dont le but est de créer des entités informatiques.
}{P006}
\explication{OD02}{de la définition}{Procédé de développement}{
La définition formulée est~: procédé dont le but est de créer des entités informatiques.

La littérature fournit ces informations qui ont servi à construire cette définition~:
\begin{itemize}
\item définition de développement~: «~Mise au point d'un appareil, d'un produit en vue de sa vente; période précédant la commercialisation.~» \cite{Larousse_developpement_0};
\item définition de procédé~: suite d’activité visant un but;
\item définition de cycle de vie du développement logiciel~: «~Ensemble des étapes permettant le développement d'un logiciel de manière systématique.~» \cite{Gdt_sdlc_0};
\item définition de cycle de développement logiciel~: \enfr
{\elide period of time that begins with the decision to develop a software product and ends when the software is delivered \elide.}
{\elide période de temps qui commence avec la décision de développer un produit logiciel et se termine avec la livraison du logiciel \elide.}
\cite{iso_24765_2017};
\item définition de cycle de vie du produit~: «~Série de phases qui représentent l’évolution d'un produit, du concept à la livraison, la croissance, la maturité et jusqu’à son retrait du marché.~» \cite{pmi_guidefr_2017}.
\end{itemize}

Les procédés de développement logiciel et système peuvent se confondre, c'est pour cela que le terme \emphase{procédé de développement} est utilisé \footnote{Il existe une très grande variété de procédés de développement. Ils se construisent grâce des modèles et des méthodes. Il existe notamment le \textit{Scrum}, l'\textit{eXtreme Programming} (XP), le \textit{Rational Unified Process} (RUP), le \textit{Personal Software Process} (PSP), la \textit{Cleanroom}, le modèle en V, le modèle en spirale ainsi que le \textit{Capability Maturity Model Integration} (CMMI).}. Aussi, le procédé de développement s'inscrit dans un procédé de gestion de projet.
}{}

\enonce{OD03}{Définition}{Processus}{
Instance de procédé.
}{P006}
\explication{OD03}{de la définition}{Processus}{
La définition formulée est~: instance de procédé.

La littérature fournit ces informations qui ont servi à construire cette définition~:
\begin{itemize}
\item définition 1 de processus~: «~Suite continue de faits ou de phénomènes qui présentent une certaine unité ou une certaine régularité dans leur déroulement.~» \cite{Gdt_processus_1} ;
\item définition 2 de processus~: «~Ensemble d'activités logiquement interreliées qui produisent un résultat déterminé.~» \cite{Gdt_processus_0};
\item définition 1 de cycle~: «~Suite de phénomènes se renouvelant sans arrêt dans un ordre immuable.~»  \cite{Robert_cycle_0};
\item définition 2 de cycle~: «~Ensemble d'opérations qui se répètent de façon constante à chaque intervalle de temps du déroulement d'un événement périodique.~» \cite{Gdt_cycle_0};
\item définition de cycle de vie du produit~: «~Série de phases qui représentent l’évolution d'un produit, du concept à la livraison, la croissance, la maturité et jusqu’à son retrait du marché.~» \cite{pmi_guidefr_2017};
\item qualificatif de cycle de vie~: incrémentiel, itératif, prédictif \cite{pmi_guidefr_2017}.
\end{itemize}
Les qualificatifs de cycle de vie en gestion de projet entremêlent les concepts de processus et de procédé. Néanmoins, il est possible de restreindre le cycle de vie au processus. Puisque le processus peut avoir une notion de régularité, le cycle de vie est un cas particulier de processus. Le processus se déroule et le but est facultatif.
}{}

\enonce{OD04}{Définition}{Artefact}{
Intrant tangible ou extrant tangible à un processus.
}{P006}
\explication{OD04}{de la définition}{Artefact}{
La définition formulée est~: intrant tangible ou extrant tangible à un processus.

La littérature fournit ces informations qui ont servi à construire cette définition~:
\begin{itemize}
\item définition d'artefact~: «~Module d'information utilisé ou produit lors de la conception d'un logiciel.~» \cite{Gdt_artefact_0}; 
\item note jointe à cette définition~: «~L'artefact peut prendre la forme d'un modèle d'architecture, d'une description, d'une documentation, etc.~» \cite{Gdt_artefact_0};
\item définition d'artefact logiciel (\textit{software artifact}, en anglais)~: \enfr
{The software artifact is a tangible machine-readable document created during software development. Examples are requirement specification documents, design documents, source code and executables.}
{Un artefact logiciel est un document tangible lisible par machine, créé lors du développement d'un logiciel. Les documents de spécification des exigences, les documents de conception, le code source et les exécutables en sont des exemples.} \ccite{iso_19506_2012}[p.~8].
\end{itemize} 
Dans le cadre d'un processus informatique, les extrants sont des artefacts informatiques et certains sont des entités informatiques.
}{}

\enonce{OD05}{Définition}{Phase d'un projet}{
Regroupement d’activités avec un début et une fin logiquement interreliés et découlant de la gestion d'un projet.
}{P006}
\explication{OD05}{de la définition}{Phase d'un projet}{
La définition formulée est~: regroupement d’activités avec un début et une fin logiquement interreliés et découlant de la gestion d'un projet.

La littérature fournit ces informations qui ont servi à construire cette définition~:
\begin{itemize}
\item définition de phase~: «~Chacun des états successifs d'un phénomène en évolution.~» \cite{Larousse_phase_0};
\item définition de phase d'un projet~: \enfr{\elide a collection of logically related project activities that culminates in the completion of one or more deliverables \elide.}{\elide un ensemble d'activités de projet logiquement liées qui aboutissent à la réalisation d'un ou plusieurs livrables \elide.}
\cite{iso_24765_2017};
\item Définition phase du projet~: «~Ensemble d'activités conjointes du projet qui aboutit à la finalisation d'un ou de plusieurs livrables.~». \ccite{pmi_guidefr_2017}[p.~718]
\end{itemize}
Les phases sont d'abord une relation entre deux points, puis vient des activités connexes ou logiquement reliées.
}{}

\enonce{OD06}{Définition}{Phase d'un procédé de développement}{
Regroupement d’activités avec un début et une fin logiquement interreliés et découlant d'un procédé de développement.
}{P006}
\explication{OD06}{de la définition}{Phases d'un procédé de développement}{
La définition formulée est~: regroupement d’activités avec un début et une fin logiquement interreliés et découlant d'un procédé de développement.

La littérature fournit ces informations qui ont servi à construire cette définition~:
\begin{itemize}
\item définition de phase~: «~Chacun des états successifs d'un phénomène en évolution.~» \cite{Larousse_phase_0};
\item phase de développement~: \enfr{Development phase A set of related tasks that produce one or more work products. Development phases can be interleaved, overlapped, and iterated as specified by the development process being used.}{Ensemble de tâches liées entre elles qui produisent un ou plusieurs livrables. Les phases de développement peuvent être imbriquées, se chevaucher et être itérées selon le processus de développement utilisé.}
\ccite{fairley_managing_2009}[p.~473];
\item Note sur \textit{software development cycle}~: \enfr{This cycle typically includes a requirements phase, design phase, implementation phase, test phase, and sometimes, installation and checkout phase.}{Ce cycle comprend généralement une phase d'analyse des besoins, une phase de conception, une phase de mise en œuvre, une phase de tests et parfois une phase d'installation et de vérification.}
\cite{iso_24765_2017}. 
\end{itemize}
}{}

\paragraphe{P007}{\fl{L'usage des catégories de documents, ainsi que le contenu des catégories de documents varie.} D'abord, les catégories de documents utilisés vont varier selon la taille du projet, les lectorats, les procédés utilisés et la réglementation applicable \ccite{nbs_guide_1987}[p.~8][rierson_do178c_2017][p.~287]. 
\sep
Ensuite, les catégories de documents elles-mêmes vont varier, puisqu'elles ne sont pas décrites formellement. 
\begin{itemize}
    \item le standard MIL-STD-498 (1994) fournit des gabarits \ccite{dod_std498_1994}[][codur_dod_2012][];
    \item l'ISO/IEC/IEEE 15289 (2011) fournit une liste de points à couvrir pour certains documents \cite{iso_15289_2011};
    \item l'ISO/IEC/IEEE 24765 (2017) fournit des descriptions sommaires \cite{iso_24765_2017}.
\end{itemize}
Ce qui demeure invariable est qu'un document est écrit par un autorat et vise un lectorat.
}{OD07, ODO8, ODO9, OD10}{OD07}

\enonce{OD07}{Assertion}{Le nombre de catégories de documents à employer pour un projet de développement dépend minimalement de ces facteurs~: la taille du projet, les lectorats, les procédés utilisés et la réglementation applicable.}{}{P007}
\explication{OD07}{sur l'assertion}{Le nombre de catégories de documents à employer pour un projet de développement dépend minimalement de ces facteurs~: la taille du projet, les lectorats, les procédés utilisés et la réglementation applicable.}{
La taille du projet, les lectorats et les procédés utilisés se retrouvent dans le 
\titleen{Management Guide for Software Documentation}{} de 1987 publiés par le National Bureau of Standards \footnote {En 1988, le nom a changé pour le National Institute of Standards and Technology (NIST).}. Voici ce qui y est inscrit~:
\quoteenfr
{The number of document types produced for any one project varies and depends on the size of the project, the audiences addressed, the number of phases identified for management control, and other factors.
}
{Le nombre de types de documents produits pour un projet donné varie et dépend de la taille du projet, des publics cibles, du nombre de phases identifiées pour le contrôle de gestion et d'autres facteurs.}
{\ccite{nbs_guide_1987}[p.~8]}

La réglementation applicable exige des catégories de documents, par exemple, dans l'avionique. Lenna Rierson l'explicite pour le Software Configuration Index (SCI) et le Software Accomplishment Summary (SAS)~:
\quoteenfr
{The SCI and SAS are the two life cycle data items that are typically submitted to the certification authority to demonstrate compliance with the DO-178C and the regulations.}
{Le SCI et le SAS sont les deux éléments de données du cycle de vie qui sont généralement soumis à l'autorité de certification pour démontrer la conformité à la norme DO-178C et aux réglementations.}
{\ccite{rierson_do178c_2017}[p.~287]}
}{}

\enonce{ODO8}{Assertion}{L'attribution d'un document à une catégorie est souvent vague.}{}{P007}
\explication{ODO8}{sur l'assertion}{L'attribution d'un document à une catégorie est souvent vague.}{
Les gabarits peuvent disparaître des standards. Un article résume l'évolution de standards dans le secteur militaire (tableau B.1) \cite{codur_dod_2012}. Il y a dans ce tableau les standards suivants ~: MIL-STD-1679 (1978), DOD-STD-2167 (1985), MIL-STD-498 (1994), IEEE/EIA 12207 (1996). Il n'y a aucun gabarit dans l'ISO/IEC/IEEE 12207\hc2017 \cite{iso_12207_2017}. 

\taben
{\textit{Coverage of some software engineering concepts by standards}}
{figure/chap1/documentation/stdfeature.png}
{(source~: \textit{Table} 1 \cite{codur_dod_2012})}{}{}{tb}{Couverture des standards sur des concepts de génie logiciel}

Les catégories de documents peuvent être définies sommairement à l'intérieur de standards. Voici des exemples provenant de l'ISO/IEC/IEEE 24765\hc2017 ayant pour titre \titlefr{Ingénierie des systèmes et du logiciel — Vocabulaire} \cite{iso_24765_2017}~:
\begin{itemize}
    \item définition de manuel d'installation~: \enfr {\elide document that provides the information necessary to install a system or component, set initial parameters, and prepare the system or component for operational use cf. diagnostic manual, installation manual, maintenance manual, operator manual, programmer manual, user manual \elide.}
    {\elide document fournissant les informations nécessaires à l'installation d'un système ou d'un composant, au paramétrage initial et à la préparation du système ou du composant pour son utilisation opérationnelle (cf. manuel de diagnostic, manuel d'installation, manuel de maintenance, manuel d'utilisation, manuel du programmeur, manuel utilisateur) \elide.} \cite{iso_24765_2017};
    \item définition 1 des spécifications des exigences logicielles (SEL)~: \enfr
    {\elide documentation of the essential requirements (functions, performance, design constraints, and attributes) of the  software and its external interfaces \elide.}
    {\elide documentation des exigences essentielles (fonctions, performances, contraintes de conception et attributs) du logiciel et de ses interfaces externes \elide.} \cite{iso_24765_2017};
    \item définition 2 de SEL~: \enfr
    {\elide structured collection of the requirements (functions, performance, design constraints, and  attributes) of the software and its external interfaces \elide.}
    {\elide recueil structuré des exigences (fonctions, performances, contraintes de conception et attributs) du logiciel et de ses interfaces externes \elide.} \cite{iso_24765_2017}.
\end{itemize}

Des catégories de documents peuvent être décrites avec des items les composant. Il y a au moins une norme qui le fait, l'ISO/IEC/IEEE 15289\hc2011 intitulé \textit{Ingénierie des systèmes et du logiciel — Contenu des systèmes et produits d'information du processus de cycle de vie du logiciel (documentation)} \cite{iso_15289_2011}. Les éléments nécessaires pour la spécification des exigences logicielles sont énumérés dans la figure~B.1.
}{}

\enoncewithoutexp{ODO9}{Définition}{Autorat}{
Ensemble des auteurs.
}{P007}

\enoncewithoutexp{OD10}{Définition}{Lectorat}{
Ensemble des lecteurs (d'un quotidien, d'une revue, etc.). \cite{GdtLectorat_0}
}{P007}

\paragraphe{P008}{\fl{Les documents ont une utilité selon des organisations de standardisation.} D'abord, plusieurs standards prescrivent des catégories de documents. Certains standards définissent ces catégories, il y a~: la MIL-STD-498, le DO-178C, l'IEC 61513 et l'ISO 15289. Les trois premiers standards s'appliquent à des secteurs critiques comprenant des entités informatiques et le dernier concerne les entités informatiques, peu importe le secteur. Ensemble, ils prescrivent au moins 76~catégories de documents, certaines partageant plusieurs éléments. 
\sep
Ensuite, les organismes de certification peuvent s'assurer du respect d'un standard pour un procédé par les artefacts générés incluant les documents. En effet, les artefacts sont les traces laissées par les processus, ils constituent la preuve de leur exécution. Ils peuvent être examinés pour comprendre les processus et pour s'assurer de la bonne mise en place d'un modèle de procédé, comme le modèle de maturité des capacités d’intégration (CMMI, \textit{Capability Maturity Model Integration} en anglais) \footnote{Le CMMI est un ensemble de bonnes pratiques pour le développement de systèmes et de logiciels.} \ccite{basque_cmmi_2009}[203].
}{OD11, OD12}{OD11}

\figannexe{
\figen
{The software requirements specification includes}
{figure/chap1/documentation/iso15289.png}
{(source~: \ccite{iso_15289_2011}[p.~71])}{}{}{tb}{Liste des éléments d'une spécification des exigences logicielles}

}

\enonce{OD11}{Assertion}{Plusieurs standards définissent des catégories de documents.}{}{P008}
\explication{OD11}{de l'assertion}{Plusieurs standards définissent des catégories de documents.}{
Voici des standards ou des normes qui définissent suffisamment les catégories de documents~: 
\begin{itemize}
    \item la MIL-STD-498 \textit{software development and documentation} est utilisée pour son aspect historique. Elle a influencé, entre autres, la DO-178 et l'ISO~12207 \ccite{dod_std498_1994}[];
    \item la DO-178C, \textit{Software Considerations in Airborne Systems and Equipment Certification} est la dernière version de la norme servant à certifier les logiciels critiques en avionique \cite{rierson_do178c_2017,hilderman_aviation_2021, gorschek_requirements_2022};
    \item l'IEC 61513 \textit{Nuclear power plants Instrumentation and control important to safety — General requirements for systems} sert pour assurer la qualité du logiciel et du matériel utilisé dans les centrales nucléaires \cite{smith_safety_2016, iec_61513_2013};
    \item l'ISO 15289 \textit{Systems and software engineering Content of life-cycle information items (documentation)} explicite les documents provenant des normes sur le cycle de vie système ou logiciel \cite{iso_15289_2011}. Quant au cycle de vie système, il est décrit dans l'ISO 15288, tandis que le cycle de vie logiciel est décrit dans l'ISO 12207.
\end{itemize}
}{}

\enonce{OD12}{Théorème}{Certains artefacts sont nécessaires pour vérifier le respect d'un procédé de développement.}{}{P008}
\explication{OD12}{du théorème}{Certains artefacts sont nécessaires pour vérifier le respect d'un procédé de développement.}{
Ce n'est pas explicité, mais plusieurs l'insinuent. Richard Basque fournit une description d'artefact dans son ouvrage sur le CMMI~: 
\quotefr 
{\elide lorsqu'on réalise une activité de développement à l'aide d'un processus, on laisse des traces \elide À la façon d'un archéologue qui examine des artefacts comme preuve de l'activité d'une civilisation \elide.}
{\ccite{basque_cmmi_2009}[p.~203]}

Leanna Rierson indique ces artefacts pour le DO-178C~:
\quoteenfr
{DO-178C section 11 identifies the life cycle data generated while planning, developing, and verifying the software. In many situations the term software life cycle data is also referred to as data items or artifacts.}
{La section 11 de la norme DO-178C identifie les données du cycle de vie générées lors de la planification, du développement et de la vérification du logiciel. Dans de nombreux cas, ces données sont également appelées éléments de données ou artefacts.}
{\ccite{rierson_do178c_2017}[p.~56]}
\vspace{-1\baselineskip}
}{}

\paragraphe{P009}{\fl{Les documents doivent être présents tout au long du processus de développement.} D'abord, plusieurs experts donnent leurs avis sur l'importance et la place de la documentation~:
\begin{itemize}
    \item Frederick P. Brooks Jr. indique que la documentation complète est nécessaire pour terminer le processus \ccite{brooks_mythical_1995}[p.~6];
    \item C. A. R. Hoare et David L. Parnas avertissent qu'il s'agit d'une erreur de faire les documents à la fin du processus \ccite{hoare_hints_1983}[][parnas_software_2001][p.~563];
    \item Ian F. Sommerville indique que certains documents servent à planifier et se retrouvent donc au début du processus \cite{sommerville_software_2001}.
\end{itemize}  
Cependant, il n'est pas dit que la documentation couvre l’entièreté du processus. Les activités sont regroupées en phases à l'intérieur du processus. Si toutes les phases sont couvertes, alors le processus le sera également. Mais, il n'y a pas de consensus sur les phases d'un processus de développement \ccite{iso_12207_2017}[p.~17]. En s'inspirant, des phases qui existent dans la littérature, les phases suivantes ont été sélectionnées~: l'analyse, la conception, la concrétisation, l'évaluation et la transition \footnote{Pour un problème donné, l'analyse est son internalisation. La conception est celle d'au moins une solution. La concrétisation passe par la création d'une machine. L'évaluation est celle de cette machine par rapport au problème et à la solution. La transition vise à transmettre cette machine.}. Ainsi, il est possible de distribuer les catégories de documents sur les différentes phases pour observer qu'ils peuvent en effet couvrir l'entièreté du processus (tableau~1.1.).
}{OD13, OD14, OD15}{OD13}

\enonce{OD13}{Axiome}{Les documents servent à faire du logiciel et à le compléter.}{}{P009}
\explication{OD13}{de l'axiome}{Les documents servent à faire du logiciel et à le compléter.}{
Plusieurs experts en informatique soulignent la place et l'importance de la documentation et expliquent où les documents interviennent dans le processus de développement~: 
\begin{itemize}
    \item avis de Hoare~: \enfr{The view that documentation is something that is added to a program after it has been commissioned seems to be wrong in principle and counterproductive in practice. Instead,  documentation must be regarded as an integral part of the process of design and coding.}{L'idée que la documentation soit un élément ajouté à un programme après sa mise en service semble erronée par principe et contreproductive en pratique. Au contraire, la documentation doit être considérée comme une partie intégrante du processus de conception et de développement.}
    \cite{hoare_hints_1983}
    \item avis de Frederick P. Brooks Jr.~:  \enfr{promotion of a program to a programming product requires its thorough documentation, so that anyone may use it, fix it, and extend it.}{promotion d'un programme en tant que produit informatique qui nécessite une documentation complète, afin que chacun puisse l'utiliser, le corriger et l'étendre.}
    \ccite{brooks_mythical_1995}[p.~6]
    \item avis de Parnas~: \enfr{If a product is not documented as it is designed, using documentation as a design medium, we will save a little today, but pay far more in the future.}{Si un produit n'est pas documenté dès sa conception, en utilisant la documentation comme outil de conception, nous économiserons un peu aujourd'hui, mais paierons beaucoup plus cher à l'avenir.}
    \ccite{parnas_software_2001}[p.~563];
    \item avis de Ian F. Sommerville~:  \enfr{1. They should act as a communication medium between members of the development team.  2. They should be a system information repository to be used by maintenance engineers.  3. They should provide information for management to help them plan, budget and schedule the software development process.  4. Some of the documents should tell users how to use and administer the system.}{1. Ils doivent servir de support de communication entre les membres de l'équipe de développement. 2. Ils doivent constituer un référentiel d'informations système à l'usage des ingénieurs de maintenance. 3. Ils doivent fournir à la direction les informations nécessaires à la planification, au budget et à l'ordonnancement du processus de développement logiciel. 4. Certains documents doivent expliquer aux utilisateurs comment utiliser et administrer le système.}
    \cite{sommerville_software_2001}.
\end{itemize}
\vspace{-1\baselineskip}
}{}

\enonce{OD14}{Axiome}{Il n'y a pas consensus sur les phases de développement.}{}{P009}
\explication{OD14}{sur l'axiome}{Il n'y a pas consensus sur les phases de développement.}{
La littérature fournit ces informations qui ont servi à construire cet axiome~:
\begin{itemize}
    \item Phases provenant d'un ouvrage de Barry W. Boehm et Richard Turner~: \enfr{concept development, requirements, design, development, maintenance}
    {développement de concepts, exigences, conception, développement, maintenance}
    \ccite{boehm_balancing_2008}[p.~194];
    \item Exemple de phases dans l'ISO~: \enfr{concept, development, production, utilization, support, and retirement}
    {conception, développement, production, utilisation, assistance et mise hors service}
    \ccite{iso_12207_2017}[p.~17];
    \item Note provenant de l'ISO qui suit l'exemple~: \enfr{Use of these terms to define stages is not normative.}
    {L'utilisation de ces termes pour définir les étapes n'est pas normative.}
    \ccite{iso_12207_2017}[p.~17].
\end{itemize}
\vspace{-1\baselineskip}
}{}

\enonce{OD15}{Assertion}{Les phases peuvent être, entre autres, l'analyse, la conception, la concrétisation, l'évaluation et la transition.}{}{P009}
\explication{OD15}{sur l'assertion}{Des phases peuvent être, entre autres, l'analyse, la conception, la concrétisation, l'évaluation et la transition.}{
Des noms de phases ont été sélectionnés par rapport à d'autres termes déjà employés par des procédés décrits à l'intérieur de normes ou d'ouvrages~:
\begin{itemize}
    \item les normes~: 
    \begin{itemize}
        \item l'ISO 15288 \cite{incose_systems_2023, iso_15288_2015} ;
        \item l'ISO 12207 \cite{iso_12207_2017} ;
        \item le MIL-STD-498 \cite{dod_std498_1994} ;
    \end{itemize}
    \item les ouvrages~: 
        \begin{itemize}
            \item \textit{Balancing Agility and Discipline: A Guide for the Perplexed} de Barry W. Boehm et et Richard Turner \ccite{boehm_balancing_2008}[p.~194];
            \item \textit{Software Project Management: A Unified Framework} de Walker Royce \ccite{royce_software_2004}[p.~120];
            \item \textit{Managing and Leading Software Projects} de Richard E. Fairley \ccite{fairley_managing_2009}[p.~56];
        \end{itemize}
\end{itemize}

La phase d'analyse~:
\begin{itemize}
    \item définition d'analyse~:  «~Opération intellectuelle consistant à décomposer un tout en ses éléments constituants et d'en établir les relations.~» \cite{Robert_analyse_0};
    \item description~: l'analyse de problèmes est en fait son internalisation et cela mène à des artefacts, tels que la spécification des exigences;
    \item connexe aux phases~: \textit{Concept Development}, \textit{Requirements}, \textit{Requirements analysis}, \textit{Business analysis}, \textit{System analysis};
\end{itemize}

La phase de conception~:
\begin{itemize}
    \item définition de conception~: «~Action de concevoir.~» \cite{Robert_conception_0} ;
    \item définition de concevoir~: «~Créer par l'imagination.~» \cite{Robert_concevoir_0};
    \item description~: la conception des solutions mène à des artefacts comme l'architecture;
    \item connexe aux phases~: \textit{Design}, \textit{Architecture};
\end{itemize}

La phase de concrétisation~:
\begin{itemize}
    \item définition de concrétisation~: «~Fait de se concrétiser, de devenir concret ou réel.~» \cite{Robert_concretisation_0} ;
    \item description~: la concrétisation des solutions mène à des artefacts comme le code source ;
    \item connexe aux phases~: \textit{Development}, \textit{Implementation};
\end{itemize}

La phase d’évaluation~:
\begin{itemize}
    \item définition d'évaluation~: «~Action d'évaluer~» \cite{Robert_evaluation_0} ;
    \item définition d'évaluer~: «~Porter un jugement sur la valeur, le prix de~» \cite{ Robert_evaluer_0};
    \item description~: l'évaluation porte sur les artefacts;
    \item connexe aux phases~: \textit{Assessment}, \textit{Verification and Validation}, \textit{Test};
\end{itemize}

La phase de transition~:
\begin{itemize}
    \item définition de transition~:  «~Passage d'un état à un autre, en général lent et graduel; état intermédiaire.~» \cite{Robert_transition_0} ;
    \item description~: la transition consiste à l'externalisation d'artefacts pour en nécessiter d'autres - il existe des transitions pour en faire l'exploitation ou la migration;
    \item connexe aux phases~: \textit{Maintenance}, \textit{Deployment}, \textit{Utilization}, \textit{Support}, \textit{Retirement}, \textit{Operation}.
\end{itemize}
\vspace{-1\baselineskip}
}{}

\paragraphe{P010}{\fl{Les documents en informatique appliquée ont comme objectifs d'aider à réaliser et à établir des logiciels.} Le premier objectif, réaliser, consiste à rendre tangible ce qui est abstrait. Un logiciel peut être considéré comme une chose immatérielle qui nécessite, pour exister, l'emploi d'un support matériel, comme des disques, des imprimés, etc. Le logiciel peut également être perceptible par les interactions rendues possibles avec l'aide d'un écran, d'un clavier, etc. Dans les deux cas, le logiciel devient tangible. 
\sep
Le second objectif, établir, vise à créer les conditions pour soutenir la pérennité. Pour le logiciel, cela englobe ce qui facilitera la validation, l'exploitation, la maintenance, et ce, pour les différents milieux. Les différentes catégories de documents présentées au tableau 1.1 peuvent servir l’un, l'autre ou même les deux objectifs.
}{OD16, OD17}{OD16}

%Maxime : une espace là :( footnote entretien utilise dans le cas de logiciel seul.

\paragraphe{}{%
\tablandscape
{Classement des catégories de documents par phase}
{figure/chap1/documentation/docpro.png}
{}{}{
\begin{tablenotes}
\footnotesize
\item[a] {Le MIL-STD-498 aborde le code source, les tests unitaires, les tests d’intégration sans les considérer comme des \emphase{documents connexes}, donc n'ont pas de \textit{Data Item Description} (DID).}
%\item[a] {Dans la MIL-STD-498, il n'y a pas de documents pour la concrétisation. Le code source, les tests unitaires, les tests d’intégration sont abordés, mais n'ont pas de \textit{Data Item Description} (DID) puisqu'ils sont considérés comme des produits et non comme des \emphase{documents connexes}.}
\item[b] {Le DO-178C considère la transition à l'extérieur du produit, donc n'a pas à être certifié.}
\item[b] {Dans le DO-178C, il n'y a pas de documents de transition, puisqu'ils sont considérés à l'extérieur du produit et qu'ils n'ont pas à être certifiés.}
\end{tablenotes}
}{H}
}{}{}

\enonce{OD16}{Définition}{Réaliser}{
Rendre tangible ce qui est abstrait.
}{P010}
\explication{OD16}{de la définition}{Réaliser}{
La définition formulée est~: rendre tangible ce qui est abstrait.

La littérature fournit ces informations qui ont servi à construire cette définition~:
\begin{itemize}
\item définition de réaliser~: «~Faire passer à l'état de réalité concrète (ce qui n'existait que dans l'esprit).~» \cite{Robert_realiser_0};
\item définition 1 de concret~: «~Qui peut être perçu par les sens ou imaginé.~» \cite{Robert-concret_0};
\item définition 2 de concret~: «~Par opposition à abstrait, qui est directement perceptible par les sens ; palpable, tangible, matériel.~» \cite{Larousse_concret_0};
\item définition de tangible~: «~Qu'on connaît par le toucher ; matériel, sensible~» \cite{Larousse_tangible_0}.
\end{itemize}
\vspace{-1\baselineskip}
}{}

\enonce{OD17}{Définition}{Établir}{
Créer les conditions pour soutenir la pérennité.
}{P010}
\explication{OD17}{de la définition}{Établir}{
La définition formulée est~: créer les conditions pour soutenir la pérennité.

La littérature fournit ces informations qui ont servi à construire cette définition~:
\begin{itemize}
\item définition 1 d'établir~: «~Mettre, faire tenir (une chose) dans un lieu et d'une manière stable~» \cite{Robert_etablir_0};
\item définition 2 d'établir~: Construire, fonder quelque chose sur quelque chose de solide~» \cite{Larousse_etablir_0};
\item définition d'être établi~: S'instaurer, exister de façon durable~» \cite{Larousse_etablir_0}.
\end{itemize}
\vspace{-1\baselineskip}
}{}



\subsection{Des écueils}

\paragraphe{P011}{Les écueils aux documents en informatique appliquée arrivent tout au long du processus de développement. Le processus de développement étant rarement linéaire, il devient difficile d'énumérer les écueils dans un ordre totalement logique. Néanmoins, étant donné le nombre important d'écueils, ils sont tout de même divisés entre ceux arrivant pendant un projet de réalisation et ceux arrivant après la réalisation.}{}{}

\subsubsection{Pendant la réalisation}

\paragraphe{P012}{Les premiers écueils présentés sont ceux tirés de l'avis des professionnels informatiques et du contexte de développement. Par la suite sont énumérés les écueils touchant à la qualité de la documentation, et plus précisément à ses caractéristiques, à ses mesures et à son contrôle.
}{}{}

\paragraphe{P013}{\fl{Plusieurs professionnels informatiques indiquent que les compétences informatiques et le temps sont essentiels pour surmonter les problèmes de documentations.} D'abord, plusieurs professionnels informatiques répondent à un sondage $(n=323)$. Ils nomment comme problèmes de documentation principaux~: les ambiguïtés, l'inexactitude et l'incomplétude. L'article conclut que, pour surmonter ces problèmes, une bonne expertise informatique est requise \cite{uddin_how_2015}. 
\sep
Pour des professionnels informatiques ayant répondu à un autre sondage~$(n=146)$, la principale cause des problèmes entourant la documentation est le manque de temps qui lui est consacré \cite{aghajani_doc_2020}. Lors de la planification, il faudrait allouer plus de temps à la documentation et la commencer au début du processus.
}{ED01}{ED01}

\enonce{ED01}{Assertion}{Le professionnel informatique (temps, compétences) est nécessaire pour la documentation.}{}{P013}
\explication{ED01}{de l'assertion}{Le professionnel informatique (temps, compétences) est nécessaire pour la documentation.}{
Un article de 2015 décrit deux sondages réalisés auprès de 323~professionnels informatiques d'IBM \cite{uddin_how_2015}. Le sujet porte sur les problèmes de documentation de l'interface du programme d’application (API, \textit{Application Programme Interface} en anglais) et les résultats apparaissent à la figure~B.2. L'article conclut que les compétences techniques sont importantes pour écrire la documentation~:

\figen
{\textit{The results of our analysis of the problems’ frequency, their severity, and the necessity to solve them. A circle’s size indicates the percentage of respondents who gave that problem top priority.}}%«~Tangled information~» and «~Excessive structural information~» are absent because no one selected them as the top problems.
{figure/chap1/documentation/ibm.png}
{(source~: \textit{Figure} 2 \cite{uddin_how_2015})}
{}{}{tb}{Causes des problèmes de documentation selon des professionnels}

\quoteenfr
{Perhaps unsurprisingly, the biggest problems with API documentation were also the ones requiring the most technical expertise to solve. Completing, clarifying, and correcting documentation require deep, authoritative knowledge of the API’s implementation. This makes accomplishing these tasks difficult for nondevelopers or recent contributors to a project.}
{Sans surprise, les problèmes les plus importants liés à la documentation des API étaient aussi ceux qui exigeaient le plus d'expertise technique. Compléter, clarifier et corriger la documentation requiert une connaissance approfondie et experte de l'implémentation de l'API. De ce fait, ces tâches s'avèrent difficiles à accomplir pour les non-développeurs ou les nouveaux contributeurs au projet.}
{\cite{uddin_how_2015}}

Un sondage effectué en 2020 interroge 146~professionnels en informatique, entre autres, sur les causes des problèmes de documentation \cite{aghajani_doc_2020}~:
\quoteenfr
{Increasing the budget dedicated to documentation was a recurring solution often mentioned by participants. This suggests that software documentation does not receive the attention it deserves when planning and allocating software resources.
\elide
Lack of time to write documentation is the only issue in this category that was indicated as important by the majority (65\%) of participants. Also, it is a frequently encountered issue. The proposed solutions boil down to: (i) explicitly allocating time/effort/resources to documentation in the project planning; and (ii) starting documentation activities in the early stage of the software lifecycle, \elide.}
{L'augmentation du budget alloué à la documentation a été une solution récurrente souvent mentionnée par les participants. Cela suggère que la documentation logicielle ne reçoit pas l'attention qu'elle mérite lors de la planification et de l'allocation des ressources logicielles.
\elide Le manque de temps pour rédiger la documentation est le seul problème de cette catégorie qui a été jugé important par la majorité (65~\%) des participants. De plus, il s'agit d'un problème fréquent. Les solutions proposées se résument à~: (i) allouer explicitement du temps, des efforts et des ressources à la documentation dans la planification du projet ; et (ii) démarrer les activités de documentation dès les premières étapes du cycle de vie du logiciel, \elide.}
{\cite{aghajani_doc_2020}}
\vspace{-1\baselineskip}
}{}

\paragraphe{P014}{\fl{Plusieurs professionnels informatiques ne désirent par mettre l'effort dans la documentation.} D'abord, parmi les discours encore fréquents à ce jour, il y a~:
\begin{itemize}
    \item \enfr{The nature of programmers is such that interesting work gets done at the expense of dull work, and documentation is dull work.}{La nature des programmeurs est telle que le travail intéressant est réalisé au détriment du travail ennuyeux, et la documentation est un travail ennuyeux.}
    \ccite{wolverton_cost_1974}[p.~615];
    \item \enfr{Overall, the key message reported by some professionals is: «~Only the code explains the code.~»}{Globalement, le message clé rapporté par certains professionnels est le suivant~: «~Seul le code explique le code.~»}
    \cite{nielebock_commenting_2019}.
\end{itemize}
Ensuite, cette aversion des professionnels informatiques à l'égard de la documentation a conduit à la réalisation d'une documentation minimaliste. L'objectif était de rendre la documentation plus attrayante et efficace. Un article publié en 1990 relève un certain nombre de problèmes au sujet de la couverture et de la compréhension de la documentation minimaliste \cite{brockmann_min_1990}.
}{ED02, ED03}{ED02}

\enonce{ED02}{Assertion}{Plusieurs professionnels informatiques considèrent le temps consacré à la documentation comme étant inutile.}{}{P014}
\explication{ED02}{de l'assertion}{Plusieurs professionnels informatiques considèrent le temps consacré à la documentation comme étant inutile.}{
En 1974, une opinion devient répandue, elle provient d'un article écrit par Ray W. Wolverton~:
\quoteenfr
{The nature of programmers is such that interesting work gets done at the expense of dull work, and documentation is dull work.}
{La nature des programmeurs est telle que le travail intéressant est réalisé au détriment du travail ennuyeux, et la documentation est un travail ennuyeux.}
{\ccite{wolverton_cost_1974}[p.~615]}

En 2010, un sondage contenant 260~questions est réalisé auprès de 83~professionnels en informatique \cite{causevic_industrial_2010}. L'objectif du sondage est d'en apprendre davantage sur les tests lors d'un procédé de développement (tableau B.2). Plusieurs ont exprimé leur désaccord avec le fait qu'une documentation compréhensible devrait être essentielle.

\taben
{\textit{Questions with the highest degree of dissatisfaction}}
{figure/chap1/documentation/dissatisfaction.png}
{(source~: \cite{causevic_industrial_2010})}
{}
{
\begin{tablenotes}
\footnotesize
\item[a] \parbox{0.95\linewidth}{\textit{the dissatisfaction $d_{q,R} = \frac{1}{|R|}\sum_{r\in R} |c_{q,r} - p_{q,r}|$ where $q$ is the question on the  practice of interest, $R$ set of respondents in a particular category, $c_{q,r}$ and $p_{q,r}$ refer to the current and the preferred practice of question $q$ as reported by the respondent $r$}}
\end{tablenotes}
}{tb}{Désaccord des professionnels envers des affirmations}

En 2015, un sondage est réalisé sur le comportement de 178~professionnels en informatique pendant la maintenance logicielle~: 
\quoteenfr
{We asked the respondents about what type of documentation they use during software maintenance. The majority use artifacts that are based on textual documentation (66\%) and on the source code by itself (70\%), but 40\% of those surveyed use a graphical notation (26\% use only UML, 9\% another notation)
\elide
14 \% do not use any complementary documentation (graphical or textual) to maintain source code, i.e., they only employ the source code as documentation.}
{Nous avons interrogé les répondants sur le type de documentation qu'ils utilisent lors de la maintenance logicielle. La majorité utilise des artefacts basés sur la documentation textuelle (66~\%) et sur le code source lui-même (70~\%), mais 40~\% des personnes interrogées utilisent une notation graphique (26~\% utilisent uniquement UML, 9~\% une autre notation).
\elide 14~\% n'utilisent aucune documentation complémentaire (graphique ou textuelle) pour la maintenance du code source, c'est-à-dire qu'ils utilisent uniquement le code source comme documentation.}
{\cite{fernandez_use_2015}}
En 2018, un article présentait les résultats d’entretiens menés auprès de 17~professionnels en informatique et d’un questionnaire rempli par un second groupe  composé de 114 répondants. Deux groupes se distinguent~: 
\quoteenfr
{Some developers seem to follow a top-down approach that is concepts oriented. These developers spend some effort to build a deeper and more thorough understanding of the API before they attempt to implement particular use cases.
\elide Developers in this group pointed out that they systematically search and regularly consult documentation provided by the API supplier.
\newpage
\elide On the other hand, there are developers that take a bottom-up approach that is code oriented.
\elide Typically, developers from this group want to start coding right away. Participants mentioned that they start learning by trying examples which they extend step-by-step. They check documentation only if they cannot find a solution on their own.}
{Certains développeurs semblent privilégier une approche descendante, axée sur les concepts. Ils consacrent du temps à acquérir une compréhension approfondie de l'API avant de tenter d'implémenter des cas d'utilisation spécifiques.
\elide Les développeurs de ce groupe ont souligné qu'ils recherchent systématiquement et consultent régulièrement la documentation fournie par le fournisseur de l'API.
\elide D'autres développeurs, en revanche, adoptent une approche ascendante, axée sur le code.
\elide Généralement, les développeurs de ce groupe souhaitent commencer à coder immédiatement. Les participants ont indiqué qu'ils apprennent en essayant des exemples qu'ils enrichissent progressivement. Ils ne consultent la documentation que s'ils ne trouvent pas de solution par eux-mêmes.}
{\cite{meng_application_2018}}
Une opinion revient fréquemment dans une étude menée en 2019 auprès de 227~professionnels informatiques sur l'efficacité des commentaires sur le code source~:
\quoteenfr
{Overall, the key message reported by some professionals is: «~Only the code explains the code.~»}
{Globalement, le message clé rapporté par certains professionnels est le suivant~: «~Seul le code explique le code.~»}
{\cite{nielebock_commenting_2019}}

En 2022, une revue de littérature de Theo Theunissen \textit{et al.} sur 63~articles publiés entre  2001 et 2018 dans le contexte des méthodes agiles. La conclusion met en évidence un constat important~:
\quoteenfr
{The challenges include: informal documentation is hard to understand, documentation is considered waste when it does not contribute to the end product, productivity is limited to the measured amount of working software, documentation is easily out-of-sync with the actual code, and there is short term focus. The practices include: a significant amount of communications happens verbally and informally; there is a positive usage of development artifacts for documentation purposes, such as TDD; and the use of architecture frameworks might positively influence documentation quality.}
{Les défis sont les suivants~: la documentation informelle est difficile à comprendre, elle est perçue comme un gaspillage lorsqu’elle ne contribue pas au produit final, la productivité est limitée à la quantité de logiciels fonctionnels mesurée, la documentation se désynchronise facilement avec le code et la vision est souvent à court terme. Les bonnes pratiques incluent~: une communication importante, verbale et informelle ; une utilisation pertinente des artefacts de développement à des fins de documentation, comme le TDD ; et l’utilisation de cadres d’architecture peut avoir un impact positif sur la qualité de la documentation.}
{\cite{theunissen_mapping_2022}}
\vspace{-1\baselineskip}
}{}

\enonce{ED03}{Assertion}{L'aversion des professionnels informatiques pour la documentation a mené à la réalisation d'une documentation minimaliste.}{}{P014}
\explication{ED03}{de l'assertion}{L'aversion des professionnels informatiques pour la documentation a mené à la réalisation d'une documentation minimaliste.}{
En 1990, un article repasse en revue la documentation minimaliste \cite{brockmann_min_1990}. L'article cite les différentes expérimentations menées par John Carroll chez IBM, HP, Xerox et Bell Northern Research. Plusieurs comportements sont énumérés, en voici quatre~: 
\quoteenfrlist
{
\item Are impatient learners and want to get started quickly on something productive
\item Skip around in manuals and on-line documents and rarely read them fully
\item Make mistakes but learn most often from correcting such mistakes
\item Are best motivated by self-initiated exploration
\item Are discouraged, not empowered, by large manuals with each task decomposed into its subtask minutiae.
}
{
\item Ce sont des apprenants impatients qui veulent se lancer rapidement dans un projet productif
\item Ils parcourent les manuels et les documents en ligne sans les lire attentivement
\item Ils font des erreurs, mais apprennent surtout en les corrigeant
\item L'exploration autonome les motive au mieux
\item Les manuels volumineux, où chaque tâche est décomposée en sous-tâches insignifiantes, les découragent plutôt que de les stimuler.
}
{\cite{brockmann_min_1990}}

L'article soulève des problèmes avec la documentation minimaliste~:
\quoteenfrlist
{
\item Learning by self-discovery is less predictable and gaps or lack of depth in learning may appear \elide
\item This design assumes a motivated audience, However, all the other document designs such as  playscript and structured writing assume just the opposite. Which is more accurate?
\item In allowing the free choosing of goals, some users picked unattainable goals or ineffective goals.
\item It’s quite different from the ways documenters have traditionally been encouraged to write \elide
\item When does one know when concision becomes cryptic? Will minimalist designers be more open to liability challenges than systematic comprehensive documenters ?
\item Will designers miss the all-important testing element-the hard work-and just employ the easiest  step, i.e., cut out words.
}
{
\item L’apprentissage par la découverte personnelle est moins prévisible et des lacunes ou un manque de profondeur dans l’apprentissage peuvent apparaître \elide
\item Cette conception suppose un public motivé. Or, toutes les autres conceptions de documents, comme les scénarios et l’écriture structurée, supposent exactement le contraire. Laquelle est la plus juste ?
\item En permettant le libre choix des objectifs, certains utilisateurs ont opté pour des objectifs inatteignables ou inefficaces.
\item C’est très différent des méthodes traditionnelles d’écriture utilisées par les rédacteurs de documents \elide
\item À quel moment la concision devient-elle obscure ? Les concepteurs minimalistes seront-ils plus ouverts aux questions de responsabilité que les rédacteurs de documents systématiques et exhaustifs ?
\item Les concepteurs négligeront-ils l’élément essentiel des tests – le travail de fond – et se contenteront-ils de la solution de facilité, c’est-à-dire de supprimer des mots ?
}
{\cite{brockmann_min_1990}}
\vspace{-1\baselineskip}
}{}

\paragraphe{P015}{\fl{La tendance actuelle privilégie les professionnels informatiques qui veulent éviter la documentation à ceux qui la préconisent.} Après les années 2000, les méthodologies agiles se sont imposées comme un courant dominant \cite{dima_agile_2018}, et ce, jusqu'au secteur critique \cite{rosa_empirical_2022}. Au fondement de ce courant se trouve le manifeste agile qui regroupe des principes. Les quatre premiers principes priorisent certaines valeurs (dont la communication) par rapport à d'autres (les éléments tangibles et vérifiables), tandis que le cinquième (et dernier) principe justifie les quatre premiers sur la base de leurs préférences personnelles \footnote{Kent Beck \textit{et al.}, site web, «~Manifesto for Agile Software Development~», \url{https://agilemanifesto.org/iso/fr/manifesto.html} (consulté le 2025-05-12).}. Voici les principes numérotés~; 
\begin{enumerate}
    \item «~Les individus et leurs interactions plus que les processus et les outils.~»
    \item «~Des logiciels opérationnels plus qu’une documentation exhaustive.~»
    \item «~La collaboration avec les clients plus que la négociation contractuelle.~»
    \item «~L’adaptation au changement plus que le suivi d’un plan.~»
    \item «~Nous reconnaissons la valeur des seconds éléments, mais privilégions les premiers.~»
\end{enumerate}
Cela a amené les professionnels informatiques à privilégier la communication à la documentation \cite{raglianti_rise_2023}.
}{ED04, ED05}{ED04}


\enonce{ED04}{Assertion}{Le courant dominant en développement de logiciel est les méthodologies agiles.}{}{P015}
\explication{ED04}{de l'assertion}{Le courant dominant en développement de logiciel est les méthodologies agiles.}{
Des interviews menés en 2018 sont réalisés avec les experts de 19 entreprises internationales en technologie de l'information~:
\quoteenfr
{most implemented model was the Agile model with its three forms, respectively: 68\% of the experts mentioned Agile as the model used within their department, while the rest mentioned the Waterfall model,
\elide
while half of these mentioned the Agile model was implemented after 2011 in their department.}
{Le modèle le plus fréquemment mis en œuvre était le modèle Agile, sous ses trois formes~: 68~\% des experts ont mentionné Agile comme le modèle utilisé au sein de leur département, tandis que les autres ont mentionné le modèle en cascade. \elide La moitié d’entre eux ont indiqué que le modèle Agile avait été mis en œuvre après 2011 dans leur département.}
{\cite{dima_agile_2018}}

\tabmediumen
{}
{figure/chap1/documentation/cico.png}
{(modifié les \textit{Tables} 7 \cite{cico_exploring_2021})}
{}
{
\begin{tablenotes}
\footnotesize
\item[a] {Retrait du filigrane \textit{Journal Pre-proof}.}
\end{tablenotes}
}{tb}{Catégorisation formée depuis les articles}

Une exploration sur les articles portant sur l'enseignement du génie logiciel conduit à obtenir 147~articles. Les critères de sélection de ces articles portaient essentiellement sur les mots-clés (génie logiciel, apprentissage, étudiant, professeur, projet et, etc.) et s'étendaient de 2008 à 2018. Les tendances formées sont décrites par le tableau B.3. Le tableau B.4 met en évidence  l'évolution du nombre d'articles par tendance, démontrant la popularité croissante des méthodes agiles dans le secteur de la recherche en enseignement du génie logiciel.
\taben
{\textit{The evolution of SE trends as topics in software engineering education}}
{figure/chap1/documentation/agileevo.png}
{(modifié la \textit{Figure}~8 \cite{cico_exploring_2021})}
{}
{
\begin{tablenotes}
\footnotesize
\item[a] {Découpage de la figure pour ne garder que le tableau et retrait du filigrane \textit{Journal Pre-proof}.}
\end{tablenotes}
}{tb}{Tendances dans l'enseignement du génie logiciel}

Un article indique la tendance de plus en plus grande à utiliser des méthodes agiles lors des exigences logicielles. Le titre de l'article est \titleen{The Current and Evolving Landscape of Requirements Engineering in Practice}{Le paysage actuel et évolutif de l'ingénierie des exigences en pratique}~:
\quoteenfr
{When asked which software development life cycle (SDLC) best describes the ones used in the project, 51\% of the projects reported using more than one approach and agile methodologies (for example, Scrum, Extreme Programming, or feature-driven development) were dominating the landscape as they were selected in 70\% of the surveyed projects [Figure~1(a)]. This is an increase of 24\% since the 2013 Survey. The Waterfall maintained its presence at the same level since 2013 (21\% of the projects).}
{Interrogés sur le cycle de vie du développement logiciel (SDLC) le plus approprié à leur projet, 51~\% des projets ont déclaré utiliser plusieurs approches. Les méthodologies agiles (par exemple, Scrum, Extreme Programming ou le développement piloté par les fonctionnalités) dominent largement, puisqu'elles ont été choisies dans 70~\% des projets analysés [Figure~1(a)]. Cela représente une augmentation de 24 \% par rapport à l'enquête de 2013. Le modèle en cascade conserve sa part de marché habituelle (21~\% des projets).}
{\cite{kassab_current_2022}}

Un article décrit la récente utilisation des méthodes agiles dans le secteur critique~:
\quoteenfr
{The development of a safety-critical system following an agile process has little empirical evidence of its effectiveness. Nevertheless, some case studies [24][25] are available in the literature, and they are discussed in the following sections.
A pharmaceutical company is the case study in which a Scrum process was applied to combine agile software development in a traditional safetycritical project in [24]. The product should be compliant with several safety standards, including those associated with the US FDA, and started with more than 100 managers, engineers, and developers involved.
\elide
Storer and Islam [25] present a case study where semi-structured interviews were performed with software engineers employed in a large avionics company in the United Kingdom. Their goal was to investigate the company’s experiences in the adoption of agile software development to SCS and the difficulties experienced.}
{Le développement d'un système critique pour la sécurité selon une méthodologie agile ne repose que sur peu de preuves empiriques de son efficacité. Néanmoins, quelques études de cas [24][25] sont disponibles dans la littérature et sont abordées dans les sections suivantes.
L'étude de cas [24] porte sur une entreprise pharmaceutique où la méthodologie Scrum a été appliquée pour combiner le développement logiciel agile à un projet critique pour la sécurité traditionnelle. Le produit devait être conforme à plusieurs normes de sécurité, notamment celles de la FDA américaine, et plus de 100 responsables, ingénieurs et développeurs y ont participé.
\elide
Storer et Islam [25] présentent une étude de cas basée sur des entretiens semi-directifs menés auprès d'ingénieurs logiciels d'une grande entreprise d'avionique au Royaume-Uni. Leur objectif était d'étudier l'expérience de l'entreprise en matière d'adoption du développement logiciel agile pour les systèmes critiques pour la sécurité et les difficultés rencontrées.}
{\ccite{gorschek_requirements_2022}[p.37-54]}

Une analyse décrit l'adoption des méthodes agiles dans le secteur de la défense américaine~:
\quoteenfr
{In 2009, the U.S. Congress enacted the National Defense Authorization Act for Fiscal Year 2010 (2010 NDAA) requiring the DoD to implement a new acquisition process for IT systems [51]. This new process included principles of agile development such as early and continual involvement of the user, multiple rapidly executed increments or releases, early successive prototyping to support an evolutionary approach, and a modular open-systems approach. Since the enactment of the 2010 NDAA, an increasing number of nonIT DoD acquisition programs have turned to agile software development as a method for delivering new and enhanced capabilities to the warfighters on a rapid and repeatable basis, avoiding delays and cost overruns associated with traditional methods such as waterfall [51].}
{En 2009, le Congrès américain a promulgué la loi d'autorisation de la défense nationale pour l'exercice 2010 (NDAA 2010), imposant au département de la Défense (DoD) la mise en œuvre d'un nouveau processus d'acquisition des systèmes informatiques [51]. Ce nouveau processus intégrait les principes du développement agile, tels que l'implication précoce et continue de l'utilisateur, le déploiement rapide de multiples incréments ou versions, le prototypage successif précoce pour soutenir une approche évolutive, et une approche modulaire de systèmes ouverts. Depuis la promulgation de la NDAA 2010, un nombre croissant de programmes d'acquisition du DoD, hors technologies de l'information, ont adopté le développement agile de logiciels comme méthode pour fournir aux combattants des capacités nouvelles et améliorées de manière rapide et reproductible, évitant ainsi les retards et les dépassements de coûts associés aux méthodes traditionnelles, telles que la méthode en cascade [51].}
{\cite{rosa_empirical_2022}}
Cette même analyse conclut que~:
\quoteenfr
{The caveat is the uneven sample of 36 for agile vs. 176 for traditional. \elide
Agile software projects appear to perform slightly better than traditional in term of productivity gains. \elide Whether agile development is more effective than traditional in general remains unanswered as differences are not significant.}
{Il convient toutefois de noter la disparité de l'échantillon~: 36 projets agiles contre 176 pour les projets traditionnels.
\elide
Les projets logiciels agiles semblent légèrement plus performants que les projets traditionnels en termes de gains de productivité. \elide La question de savoir si le développement agile est globalement plus efficace que le développement traditionnel reste ouverte, car les différences observées ne sont pas significatives.}
{\cite{rosa_empirical_2022}}
\vspace{-1\baselineskip}
}{}

\enonce{ED05}{Assertion}{La communication est privilégiée à la documentation.}{}{P015}
\explication{ED05}{de l'assertion}{La communication est privilégiée à la documentation.}{
En 2002, un article publie les résultats d'un sondage mené auprès de 48~professionnels informatiques sur la documentation. Dans la conclusion, il est mentionné~:
\quoteenfrlist
{
\item Document content can be relevant even if it is not up to date. (However, keeping it up to date is still a good objective). As such, technologies should strive for easy to use, lightweight and disposable solutions for documentation.
\item Documentation is an important tool for communication and should always serve a purpose. Technologies should enable quick and efficient means to communicate ideas as opposed to providing strict validation and verification rules for building facts.
}
{
\item Le contenu d'un document peut être pertinent même s'il n'est pas à jour. (Toutefois, le maintenir à jour reste un objectif louable). De ce fait, les technologies devraient privilégier des solutions de documentation simples d'utilisation, légères et éphémères.
\item La documentation est un outil de communication essentiel et doit toujours avoir une utilité. Les technologies devraient permettre de communiquer des idées rapidement et efficacement, plutôt que d'imposer des règles strictes de validation et de vérification pour établir des faits.
}
{\cite{forward_relevance_2002}}

Une analyse de 2023 porte sur les fichiers de documentation \textit{README} (fichiers lisez-moi en français) à l'intérieur de projets $(n>10000)$ en code source ouvert. L'analyse met en évidence que les instruments de communication prennent de plus en plus de place au détriment de la documentation~:
\quoteenfr
{Classical software documentation is being replaced by «~communication~». At least in open source software on GitHub, it is supported by a plethora of platforms characterized by high throughput, volatility, and heterogeneity. The original vision of on-demand developer documentation [55] advocated for a paradigm shift. A shift did happen, but it was not in the direction foreseen by Robillard et al. five years ago. The new communication platforms bring new challenges and opportunities for modern software documentation. It is time to shed light on new forms of documentation.}
{La documentation logicielle classique cède la place à la communication. Du moins, dans le domaine des logiciels libres sur GitHub, cette dernière s'appuie sur une multitude de plateformes caractérisées par un débit élevé, une grande volatilité et une forte hétérogénéité. La vision initiale d'une documentation développeur à la demande [55] préconisait un changement de paradigme. Ce changement a bien eu lieu, mais pas dans le sens envisagé par Robillard \textit{et al}. il y a cinq ans. Les nouvelles plateformes de communication soulèvent de nouveaux défis et offrent de nouvelles opportunités pour la documentation logicielle moderne. Il est temps d'explorer ces nouvelles formes de documentation.}
{\cite{raglianti_rise_2023}}

Voici un résumé du tableau \emphase{Documentation tools in CSD retrieved from the studies for RQ2. (Table 11)} présenté dans la revue de littérature de 2022 de Theo Theunissen \textit{et al.}, qui présente les outils et le nombre d'études en lien \cite{theunissen_mapping_2022}: Word, Excel, Powerpoint, etc. (40); \textit{Wiki} (39); \textit{Source control} (5); \textit{Scripts} (9); \textit{Markdown editors} (2); \textit{Verbal communication} (7).
}{}

\paragraphe{P016}{\fl{En général, la qualité de la documentation est considérée secondaire.} D'abord, il existe un intérêt sur les modèles de qualité, même si cet intérêt baisse depuis 2010 pour les six caractéristiques les plus notables que sont~: \enfr{functionality, maintainability, reliability, efficiency, usability, portability}{fonctionnalité, maintenabilité, fiabilité, efficacité, facilité d'utilisation, portabilité} \cite{ndukwe_how_2023}. Progressivement, le développement de caractéristiques de qualité et de métriques se porte davantage sur la qualité du code source plutôt que sur celle du produit \cite{colakoglu_metrics_2021, ndukwe_how_2023}.
%\sep
Pourtant, la qualité du logiciel est intrinsèquement liée à celle de la documentation. Plusieurs experts provenant de secteurs critiques nomment des métriques importantes en lien avec la documentation \cite{ming_ranking_2003}~:
\begin{itemize}
    \item le nombre de changements dans la spécification des exigences ;
    \item la densité de défauts à l'intérieur de la conception.
\end{itemize}
Un sondage de 2014 mené auprès de professionnels informatiques $(n=88)$ révèle que 60~\% d'entre eux préfèrent s'abstenir d'utiliser des standards ou des modèles de qualité pour leur documentation, tandis qu'au moins 30~\%  utilisent d'anciens\footnote{Cela peut être pour de bonnes raisons.} standards ou modèles de qualité \cite{plosch_doc_2014}.
Ce manque d'intérêt se produit aussi en recherche sur la documentation logicielle. Une citation de David L. Parnas explique ce peu d’intérêt~: 
\quoteenfr
{Most Computer Science researchers do not see software documentation as a topic within computer science. They can see no mathematics, no algorithms, and no models of computation; consequently, they show no interest in it.}
{La plupart des chercheurs en informatique ne considèrent pas la documentation logicielle comme un sujet relevant de l'informatique. Ils ne voient ni mathématiques, ni algorithmes, ni modèles de calcul ; par conséquent, ils ne s'y intéressent pas.}
{\cite{parnas_precise_2011}}

}{ED06, ED07, ED08, ED09}{ED06}

\enonce{ED06}{Axiome}{Il y a une dépendance entre la qualité du logiciel et la qualité de la documentation logicielle.}{}{P016}
\explication{ED06}{de l'axiome}{Il y a une dépendance entre la qualité du logiciel et la qualité de la documentation.}{
Un article publié en 2003 analyse l'avis d'experts au sujet de différentes métriques \cite{ming_ranking_2003}. Ces experts sont~: Alain Abran, David Card, William Everett, Jon Hagar, Hebert Hecht, Watts Humphrey, Michael Lyu, Jean-Claude Laprie, William Petrick et Allen Nikora. Les trois mesures les plus importantes par phase de développement sont présentées au tableau B.5. Plusieurs d'entre elles concernent directement la documentation ou des artefacts contenus dans la documentation (\textit{Requirements specification change request, Cyclomatic complexity}, etc.). 
\taben
{\textit{Top Three Measures per Phase}}
{figure/chap1/documentation/fault.png}
{(source~: \textit{Table} 9 \cite{ming_ranking_2003})}
{}{}{tb}{Experts indiquant les trois mesures les plus importantes}
}{}

\enonce{ED07}{Citation}{Il y a un manque d’intérêt scientifique sur l'étude de la documentation.}{}{P016}
\explication{ED07}{de la citation}{Il y a un manque d’intérêt scientifique sur l'étude de la documentation.}{
David L. Parnas explique le peu d’intérêt envers la documentation de la part de la recherche \cite{parnas_precise_2011}. Selon lui, l'absence de modèle y est pour beaucoup. (voir P11 p.\href{P016}{\pageref{paragraphe:P016}} la citation complète) 
}{}

\enonce{ED08}{Assertion}{Depuis 2010, l'attention se concentre désormais sur la qualité du code source plutôt que sur celle du produit.}{}{P016}
\explication{ED08}{de l'assertion}{Depuis 2010, l'attention se concentre désormais sur la qualité du code source plutôt que sur celle du produit.}{
En 2023, un article a été publié qui regroupait deux revues de littérature et un sondage sur le thème de la qualité logiciel à travers le temps \cite{ndukwe_how_2023}. Les caractéristiques de qualité logiciel ont une plus grande fréquence en 2010. Les principales caractéristiques de qualité d'un produit en 2010 avec la fréquence d'apparition sont ; \textit{functionality (19), maintainability (17), reliability (14), efficiency (13), usability (12), portability (6)}. Cela coïncide avec la popularité des modèles de qualité provenant de l'ISO 9126 puis de l'ISO 25010. Après, un déclin s'amorce et arrive des caractéristiques de qualité sur le code source. En 2023, les caractéristiques de qualité observée pour le code source avec la fréquence d'apparition sont ; \textit{reusability (10), security (9), readability (6), completeness (3), reliability (3), conciseness (2), correctness (2)}.

Des métriques sur les documents s'appliqueraient sur toutes les phases. Une revue de littérature de 2021 porte sur les métriques se retrouvant dans des articles de 2009 à 2019 \cite{colakoglu_metrics_2021}. Ces articles peuvent s'adresser à toutes les phases ou à certaines, voici la distribution des articles par phase~: \textit{all (14); code (43); planification (1); verification and validation (5); requirement (3); design (10)}.
}{}

\figannexe{
\figen
{\textit{Quality models/standards used for ensuring software  documentation quality}}
{figure/chap1/documentation/qualitymodel.png}
{(source~: \textit{Figure} 8 \cite{plosch_doc_2014})}
{}{}{tb}{Standards utilisés pour la documentation selon des professionnels}
}

\enonce{ED09}{Assertion}{Plusieurs affirment ne pas utiliser de modèles de qualité pour la documentation ou utiliser d'anciens modèles de qualité.}{}{P016}
\explication{ED09}{de l'assertion}{Plusieurs affirment ne pas utiliser de modèles de qualité pour la documentation ou utiliser d'anciens modèles de qualité.}{
En 2014, un sondage sur la documentation interroge 88~professionnels informatiques \cite{plosch_doc_2014}. Certains résultats se retrouvent dans la figure~B.3~: près de 60~\% des professionnels informatiques n’utilisent aucun standard ou modèle de qualité. Seuls 9~\% déclarent utiliser un standard ou un modèle de qualité actuel, soit  l'ISO 25000 et l'ISO 26514. Voici une description de quelques standards~:
\begin{itemize}
    \item l’IEEE~830-1998 est remplacé en 2011 par l’ISO~ 29148, la dernière version date de 2018. Il s'agit d'un guide pour écrire les différents types d'exigences. Les types d'exigences sont : \enfr{\elide business requirements specification, stakeholder requirements specification, system requirements specification, software requirements specification \elide.}{\elide spécification des exigences métier, spécification des exigences des parties prenantes, spécification des exigences système, spécification des exigences logicielles \elide.} \cite{iso_29148_2018};
    \item l'ISO 9126 est remplacée par l’ISO~25010 qui est comprise dans la série ISO~25000 qui date de 2005. À l'époque, elle fournit quelques métriques sur la documentation pour mesurer~: l'efficacité de la documentation d'utilisateur, la complétude de la documentation d'utilisateur, la traçabilité entre les différents documents \cite{iso_9126part1_2001, iso_9126part2_2003, iso_9126part3_2003, iso_9126part4_2004};
    \item l'ISO~25000 décrit le \titleen{Systems and software Quality Requirements and Evaluation}{Exigences et évaluation de la qualité des systèmes et des logiciels} (SQuaRE). L'ISO~25010 décrit quant à elle les qualités d'un produit.
\end{itemize}
\vspace{-1\baselineskip}
}{}

\paragraphe{}{%
\tab%
{Problèmes sur la doc. comprenant des caractéristiques de qualité}%
{figure/chap1/documentation/quality.png}%
{}{}%
{}{tb}%

}{}{}

\paragraphe{P017}{\fl{Plusieurs caractéristiques de qualité portant sur la documentation sont relatives, donc plus difficilement objectivables.}
Deux analyses fournissent ou dissimulent plusieurs caractéristiques de qualité sur la documentation. La première s'applique sur des articles traitant des problèmes de documentation sélectionnés depuis une revue de littérature \mbox{$(n=69)$  \cite{zhi_cost_2015}}. La deuxième s'applique sur différents artefacts $(n=878)$ concernant des problèmes avec la documentation \cite{aghajani_software_2019}. Ces problèmes et les caractéristiques de qualité qu'ils contiennent sont résumés dans le tableau 1.2. Dans ce tableau, il y a des critères relatifs aux autres artefacts, au lectorat ou à l'autorat~: la complétude, l'accessibilité, la compréhension, la mise à jour, la maintenabilité. Lorsqu'il est demandé à des professionnels informatiques leurs méthodes d'évaluations de la documentation, la plupart d’entre eux affirment qu’ils n’utilisent que des revues informelles pour évaluer ces différentes caractéristiques de qualité \cite{plosch_doc_2014}.
}{ED10, ED11}{ED10}

\figannexe{
\figen
{\textit{Documentation quality attributes}}
{figure/chap1/documentation/zhi.png}
{(source~: \textit{Figure}~18 \cite{zhi_cost_2015})}
{}
{
\begin{tablenotes}
\footnotesize
\item[a] {Degré un~: articles traitant la caractéristique de manière qualitative.}
\item[b] {Degré deux~: articles traitant la caractéristique de manière quantitative.}
\end{tablenotes}
}{tb}{Nombre d'articles abordant la qualité de la documentation}
}

\enonce{ED10}{Assertion}{Les caractéristiques de qualité de documentation peuvent être déduites de la littérature ou des artefacts.}{}{P017}
\explication{ED10}{de l'assertion}{Les caractéristiques de qualité de documentation peuvent être déduites de la littérature ou des artefacts.}{
Une revue de littérature portant sur le coût, le bénéfice et la qualité de la documentation conduit à considérer 69~ articles publiés sur une période s'étalant entre 1971 et 2011  \cite{zhi_cost_2015}. Cette revue a permis de déduire de ces différents articles, des caractéristiques de qualité pour la documentation qui se trouvent illustrées dans la figure~B.4. 

En 2019, une étude a été menée sur les problèmes reliés à la documentation \cite{aghajani_software_2019}. Au total, 878~artefacts ont été examinés, incluant des courriels, des discussions sur des forums et diverses demandes, tous reliés à la documentation. Voici le classement de ceux-ci sous différentes caractéristiques dans le tableau B.6. Il y a la complétude, l'usage et la lisibilité qui reviennent fréquemment.
\tabtabensmall
{}
{
\bgroup
\def\arraystretch{0.5}
\begin{tabular}{p{10cm}p{6.4cm}}
    \begin{itemize}[leftmargin=-0cm]
        \item [] \textit{Information content - What} (485)~:
        \begin{itemize}
        \item \textit{correcteness} (72);
        \item \textit{completeness} (268);
        \item \textit{up-to-dateness} (190);
        \end{itemize}
        \item [] \textit{Process related - How} (81)~:
        \begin{itemize}
        \item \textit{internationalization} (20);
        \item \textit{contributing to documentation} (50);
        \item \textit{doc-generator configuration} (4);
        \item \textit{development issues caused by documentation} (8);
        \item \textit{traceability} (5);
        \end{itemize}
    \end{itemize}
    &
    \begin{itemize}[leftmargin=-0cm]
        \item [] \textit{Information content - How} (255)~:
        \begin{itemize}
        \item \textit{usability} (138);
        \item \textit{maintainability} (21);
        \item \textit{readability} (107);
        \item \textit{usefulness} (20);
        \end{itemize}
        \item [] \textit{Tool related - How} (134)~:
        \begin{itemize}
        \item \textit{bug/issue} (17);
        \item \textit{support/expectations} (34);
        \item \textit{help required} (90);
        \item \textit{tool migration} (4).
        \end{itemize}
    \end{itemize}
\end{tabular}
\egroup
}
{(modifié de~: \ccite{aghajani_software_2019})}
{}
{
}{H}{Problèmes reliés à la documentation}

Plusieurs de ces caractéristiques de qualité reviennent encore lors d'une revue de littérature comportant 14~articles portant sur les problèmes de documentation (tableau~B.7) \cite{santos_review_2020}.
}{}

\figannexe{
\taben
{\textit{Quality attributes mentioned on each publication.}}
{figure/chap1/documentation/santos.png}
{(source~: \textit{Table} 1 \cite{santos_review_2020})}
{}{}{tb}{Caractéristiques de qualité de la doc. abordées dans des articles}
}

\enonce{ED11}{Citation}{Plusieurs professionnels informatiques affirment évaluer la documentation à travers des revues informelles.}{}{P017}
\explication{ED11}{de la citation}{Plusieurs professionnels informatiques affirment évaluer la documentation à travers des revues informelles.}{
En 2014, un sondage interroge 88~professionnels informatiques sur la documentation, sur les techniques de vérification et la validation employées (figure~B.5) la revue informelle est la technique la plus utilisée \cite{plosch_doc_2014}. D'ailleurs, dans le même article, plusieurs affirment ne pas utiliser de modèles de qualité (\hyperref[explication:ED09]{vers ED09}).  
}{}

\figannexe{
\figens
{\textit{Software documentation analysis techniques per quality attribute}}
{figure/chap1/documentation/tools.png}
{(source~: \textit{Figure}~13 \cite{plosch_doc_2014})}
{}{}{H}{Techniques d'évaluation de la qualité de la documentation utilisées par des professionnels}
\vspace{-1\baselineskip}
}

\paragraphe{P018}{\fl{Les mesures sur les documents sont de moins en moins mises de l'avant.} En 1990, un sondage mené dans l'industrie $(n=800)$ indique que certaines mesures touchent directement la documentation \ccite{hetzel_making_1993}[p.~6]~:
\begin{itemize}
    \item le niveau de complétude et d'exactitude de la documentation $(42~\%)$;
    \item la taille et la complexité de la documentation $(20~\%)$.
\end{itemize}
En 2018, un sondage auprès de 129~professionnels informatiques indique qu'au plus deux individus connaissent certaines mesures applicables aux documents, sans pour autant préciser lesquelles \cite{eisty_survey_2018}. 
\sep
Sans mesure, il devient impossible d'établir des comparatifs et, ainsi, d'effectuer des contrôles. Horst Zuse indique l'importance d'intégrer les mesures dès la première phase du développement, afin d'épargner l'effort et d'éviter des problèmes, du moins lors de la maintenance \ccite{zuse_framework_1998}[p.~491].
}{ED12}{ED12}

\enonce{ED12}{Assertion}{Il y a une tendance à moins mesurer la documentation.}{}{P018}
\explication{ED12}{de l'assertion}{Il y a une tendance à moins mesurer la documentation.}{
\tabsmallens
{\textit{List of the mostly used measures in industry}}
{figure/chap1/documentation/zuse.png}
{(source~: \textit{Figure} 9.3 \ccite{zuse_framework_1998}[p.~491])}
{}{}{tb}{Mesures les plus utilisées en industrie en 1990}
L'ouvrage \textit{A framework of Software Measurement} de Horst Zuse cite de Bill Hetzel, les résultats sont dans la tableau~B.8 et il en déduit ceci~: 
\quoteenfr
{It is interesting to observe that most of the listed measures do not characterize the product quality. Many of them can only be used in the maintenance phase. In order to avoid maintenance effort and problems, software measures also should be used in the very early phases of the software life-cycle. It is our view, that software measurement is necessary during the whole software life-cycle.}
{Il est intéressant de constater que la plupart des mesures énumérées ne caractérisent pas la qualité du produit. Nombre d'entre elles ne sont utiles qu'en phase de maintenance. Afin d'éviter les efforts et les problèmes de maintenance, il est essentiel d'intégrer les mesures logicielles dès les premières phases du cycle de vie du logiciel. Nous sommes d'avis que la mesure logicielle est nécessaire tout au long de ce cycle de vie.}
{\ccite{zuse_framework_1998}[p.~491]}
Dans l'ouvrage de Bill Hetzel, on trouve plus de détails sur le sondage $(n=800)$ qui a permis de constituer le tableau~B.8~:
\quoteenfr{
In 1990 Software Quality Engineering conducted a large-scale survey aimed at measuring how industry was using software measurements and to benchmark what best companies and projects were doing. \elide The survey was distributed to eight hundred software organiza-
tions around the world. Respondents were asked to report their organization’s degree of usage of a representative list of selected software measures. Company practices were highly variable. A very
small percentage of companies reported regular use of most of the measures. Another one-third to one-half used none of the measures!
}{ En 1990, \textit{Software Quality Engineering} a mené une vaste enquête visant à mesurer l'utilisation des mesures logicielles dans l'industrie et à identifier les meilleures pratiques des entreprises et des projets les plus performants. \elide L'enquête a été diffusée auprès de huit cents organisations logicielles à travers le monde. Les répondants ont été invités à indiquer le degré d'utilisation, au sein de leur organisation, d'une liste représentative de mesures logicielles sélectionnées. Les pratiques des entreprises étaient très variables. Un très faible pourcentage d'entreprises ont déclaré utiliser régulièrement la plupart des mesures. Un tiers à la moitié n'en utilisaient aucune!
}{\ccite{hetzel_making_1993}[p.~6]}

En 2018, un sondage interroge 129~professionnels informatiques dans le domaine de la recherche sur l'usage des mesures, il y aurait au plus deux personnes connaissant des mesures sur la documentation~: 
\quoteenfr
{That is (1,0) indicates one person knew the metric, but no one actually used it. each metric was mentioned as known and as used, respectively.
\elide
Process Metrics: amount of documentation (1,0), app launch count (1,1), authors/committers (1,0), cycle time (3,3), development hours/story (1,1), development man [person] years (1,0), documentation (1,0) \elide.}
{Autrement dit, (1,0) indique qu'une personne connaissait la métrique, mais que personne ne l'a utilisée. Chaque métrique a été mentionnée comme connue et comme utilisée.
\elide Métriques de processus~: quantité de documentation (1,0), nombre de lancements d'applications (1,1), auteurs/contributeurs (1,0), temps de cycle (3,3), heures de développement par \textit{user story} (1,1), années-personnes de développement (1,0), documentation (1,0) \elide.}
{\cite{eisty_survey_2018}}
\vspace{-1\baselineskip}
}{}

\paragraphe{P019}{\fl{La documentation, comme plusieurs autres artefacts en informatique, ne dispose pas de mesures permettant un contrôle.} D'abord, les mesures populaires, comme le nombre de lignes de code, sont utiles une fois le projet terminé. Un article de 2020 analyse 81~articles provenant de conférences sur les langages de programmation et établit que 65~articles utilisent le nombre de lignes ou le nombre de lignes de code source \cite{alpernas_wonderful_2020}. Alain Abran donne quelques limitations de cette mesure \ccite{abran_metrology_2010}[p.~87]. Elle est disponible que lorsque le projet est terminé et elle est dépendante des technologies utilisées. Cela n'empêche pas Barry W. Boehm et P. N. Papaccio de fournir plusieurs résultats sur le coût relatif de la documentation \cite{boehm_understanding_1988}. Dans l'un de ses ouvrages, Barry W. Boehm indique qu'il utilise une liste de contrôle pour compter précisément le nombre de lignes de code source \ccite{boehm_cocomo2_2000}[p.~16]. Le même principe peut s'appliquer pour le nombre de lignes de texte d'un document.
\sep
Ensuite, les mesures ne permettent pas d'établir formellement le coût relatif de la documentation. Selon l'United States Air Force \ccite{horowitz_practical_1975}[p.~3-14] et la National Aeronautics and Space Administration \ccite{nasa_guide_1995}[p.~81], le coût relatif de la documentation se situerait autour de 10~\% et 11~\% respectivement. Selon Barry W. Boehm et P. N. Papaccio, cela avoisinerait souvent 60~\% \cite{boehm_understanding_1988}. D'ailleurs, ces résultats sont fondés sur les projets qui serviront par la suite à Barry W. Boehm dans la création du \textit{Constructive Cost Model} (COCOMO). 
\sep
Ainsi, ces mesures n'arrivent pas à effectuer un contrôle lors du développement ou à distinguer clairement l'effort exigé entre la documentation et le reste du développement. Cependant, il est possible de mesurer les documents contenant les modèles avec la méthode des points de fonction \footnote{Mesurer les points de fonction d'un document (spécification des exigences) exprimant un modèle permet d'obtenir des estimations fiables de l'effort nécessaire pour produire les autres artefacts.}, ce qui permet d'obtenir des estimations de l'effort requis pour produire d'autres artefacts. Cela confère à ces documents une importance accrue. \ccite{iso_19761_2011}[p.~\textit{V}]. 
}{ED13, ED14, ED15}{ED13}

\enonce{ED13}{Assertion}{Les métriques courantes sur le code source et la documentation sont utiles une fois le projet terminé.}{}{P019}
\explication{ED13}{l'assertion}{Les métriques courantes sur le code source et la documentation sont utiles une fois le projet terminé.}{
Un article de 1988 où Barry W. Boehm et P. N. Papaccio donne des résultats sur le coût de la documentation~:
\quoteenfr
{The volume of documentation with respect to lines of code tends to vary by application; Reference [70] reports a typical 28 pages of documentation per thousand instructions (pp/KDSI) for internal commercial programs and a typical 66 pp/KDSI for commercial software products of the same size (50 KDSI). The proportion of documentation-related to code-re: lated effort averaged about 60 :40 over the COCOMO data base of projects [23] and about 67:33 for large TRW projects [20].}
{Le volume de documentation par rapport au nombre de lignes de code varie généralement selon l'application ; la référence [70] indique une moyenne de 28 pages de documentation pour mille instructions (pp/KDSI) pour les programmes commerciaux internes et de 66 pp/KDSI pour les logiciels commerciaux de taille équivalente (50 KDSI). Le rapport entre l'effort consacré à la documentation et celui consacré au code est en moyenne d'environ 60:40 dans la base de données COCOMO [23] et d'environ 67:33 pour les grands projets TRW [20].}
{\cite{boehm_understanding_1988}}

En 2020, un article fait l'analyse de 81~articles provenant de différentes conférences en informatique. Parmi ces articles, 35~articles exposent la taille du code source, et 28 de ceux-là utilisent le nombre de lignes ou le nombre de lignes de code source pour le faire. 
\quoteenfr
{we surveyed recent papers from top-tier conferences in Programming Languages, Software Engineering,  and Systems: POPL, PLDI, OOPSLA, ICSE, ASE, ASPLOS,  OSDI, and SOSP.
\elide
Overwhelmingly, code is measured in lines of code. Of the  papers we surveyed, 80\% of the papers who employ code  metrics use loc or sloc.}
{Nous avons analysé des articles récents issus de conférences de premier plan dans les domaines des langages de programmation, du génie logiciel et des systèmes~: POPL, PLDI, OOPSLA, ICSE, ASE, ASPLOS, OSDI et SOSP.
\elide La taille du code est très largement mesurée en lignes. Parmi les articles analysés, 80 \% de ceux qui utilisent des métriques de code emploient le nombre de lignes (loc) ou le nombre de lignes de code source (sloc).}
{\cite{alpernas_wonderful_2020}}

Alain Abran écrit dans son ouvrage \textit{Software Metrics and Software Metrology} sur les limitations à compter les lignes de code source~:
\quoteenfr
{Developers of other types of software have had to use counts of Source Lines of Code (SLOC) as their size measure. But, as this size is only known accurately after the software is complete, it is difficult to use it for estimation. Moreover, being technology-dependent, SLOC counts are difficult to use for performance comparisons, especially across multiple technologies}
{Les développeurs d'autres types de logiciels ont dû utiliser le nombre de lignes de code source (SLOC) comme mesure de taille. Cependant, comme cette taille n'est connue avec précision qu'une fois le logiciel terminé, il est difficile de l'utiliser pour des estimations. De plus, étant donné sa dépendance à la technologie, le nombre de lignes de code source est difficile à utiliser pour comparer les performances, notamment entre différentes technologies.}
{\ccite{abran_metrology_2010}[p.~87]}

Aussi pour uniformiser l'usage de cette métrique Barry W. Boehm utilise une liste de contrôle créée par le Software Engineering Institute (SEI) \ccite{boehm_cocomo2_2000}[p.~16].
}{}

\enonce{ED14}{Assertion}{Il est difficile de connaître le coût de la documentation, même proportionnellement au code source.}{}{P019}
\explication{ED14}{sur l'assertion}{Il est difficile de connaître le coût de la documentation, même proportionnellement au code source.}{
En 1975, Barry W. Boehm produit un article ayant pour titre \titleen{The high cost of software}{Le coût élevé des logiciels} avec les données provenant du document \textit{(October 1974) Air Force-Industry Workshop on Software Costs held at USAF Electronics Systems Division.}, il indique que le coût de la documentation peut représenter 10~\% du coût total~: 
\quoteenfr
{Most of the software data come from the Air Force, primarily from the CCIP-85 study‘ and the recent Air Force Industry Software Cost Workshop, but they are probably fairly representative of other software activities elsewhere.
\elide
Amount of documentation. Experience indicates that documentation costs run about 10 percent of the total software development cost. For non automated documentation \elide.}
{La plupart des données logicielles proviennent de l'Armée de l'air, principalement de l'étude CCIP-85 et du récent atelier conjoint Armée de l'air-industrie sur les coûts logiciels, mais elles sont probablement assez représentatives d'autres activités logicielles.
\elide Quantité de documentation. L'expérience montre que les coûts de documentation représentent environ 10 \% du coût total de développement logiciel. Pour la documentation non automatisée \elide.}
{\ccite{horowitz_practical_1975}[p.~3-14]}

Un article de 1988 où Barry W. Boehm et P. N. Papaccio donnent des résultats sur le coût de la documentation qui avoisine les 60 \%~:
\quoteenfr
{\elide effort averaged about 60~:40 over the COCOMO data base of projects [23] and about 67:33 for large TRW projects [20].}
{\elide l'effort consacré à la documentation et celui consacré au code est en moyenne d'environ 60:40 dans la base de données COCOMO [23] et d'environ 67:33 pour les grands projets TRW [20].}
{\cite{boehm_understanding_1988}}

En 1995, la National Aeronautics and Space Administration (NASA) publie son \textit{Software Measurement Guidebook} avec un résultat sur le coût de la documentation de 11 \% (figure~B.6) \cite{nasa_guide_1995}~:
\quoteenfr
{These data are easy to collect in most software organizations. figure~6-7 shows the data collected from one large NASA organization.}
{Ces données sont faciles à collecter dans la plupart des entreprises de développement logiciel. La figure~6-7 présente les données recueillies auprès d'une grande organisation de la NASA.}
{\cite{nasa_guide_1995}}

\figmedium
{Répartition typique des ressources d'un projet logiciel affilié à la NASA}
{figure/chap1/documentation/nasa.png}
{(source~: \textit{Figure}~6-7 \ccite{nasa_guide_1995}[p.~81])}
{}{}{tb}

Dans l'appendice du guide, il y a le \titleen{Personnel Resources Form}{Le formulaire du temps travaillé} \ccite{nasa_guide_1995}[p.~115]. Ce formulaire confirme que certains artefacts ne sont pas considérés comme de la documentation, même s’ils se retrouvent inclus dans celle-ci. 
}{}

\enonce{ED15}{Assertion}{Les documents contenant des modèles peuvent servir à effectuer un contrôle.}{}{P019}
\explication{ED15}{sur l'assertion}{Les documents contenant des modèles peuvent servir à effectuer un contrôle.}{
La norme ISO 19761 intitulée \textit{Software engineering COSMIC: A functional size measurement method} indique~:
\quoteenfr{
The International Standard aims to meet the needs of 
\begin{itemize}
    \item a) software suppliers facing the task of translating customer requirements into the functional size of software  to be produced as a key activity in their project cost estimating,
    \item b) customers who want to know the functional size of delivered software as an important component of  measuring supplier performance.
\end{itemize}
}{La norme internationale vise à répondre aux besoins a) des fournisseurs de logiciels qui doivent traduire les exigences des clients en taille fonctionnelle du logiciel à produire, une activité clé dans l'estimation des coûts de leur projet, b) des clients qui souhaitent connaître la taille fonctionnelle du logiciel livré, un élément important pour mesurer la performance du fournisseur.}
{\ccite{iso_19761_2011}[p.~v]}
En 2016, Anandi Hira et Barry W. Boehm exposent l'efficacité des points de fonction dans les estimations de l'effort pour la maintenance logicielle et sa baisse d'efficacité lorsque cette métrique y est ajoutée~: 
\quoteenfr{
\elide In UCC’s development environment, FPs led to accurate effort estimates for projects adding new functions. Using a FPs to SLOC ratio definitely adds additional uncertainty, and therefore, cost estimators should not use this method for effort estimation if they expect or require accurate effort estimates.
}{\elide Dans l'environnement de développement d'UCC, les points de fonction ont permis d'obtenir des estimations d'effort précises pour les projets intégrant de nouvelles fonctionnalités. L'utilisation d'un ratio points de fonction/lignes de code source (SLOC) accroît indéniablement l'incertitude; par conséquent, les estimateurs de coûts ne devraient pas recourir à cette méthode s'ils attendent ou exigent des estimations d'effort précises.
}{\cite{hira_function_2016}}
\vspace{-1\baselineskip}
}{}

\paragraphe{P020}{\fl{Il faut continuellement maintenir la cohérence à l'intérieur d'un document et entre les documents tout au long du processus.}
La documentation accompagne plusieurs autres activités, comme le font la planification et les tests. Donc, ce n'est pas possible de circonscrire la manipulation des documents à une phase \ccite{schach_object_2011}[p.~17]. Plusieurs professionnels informatiques indiquent qu'ils doivent modifier au moins partiellement un document, peu importe la phase \cite{plosch_doc_2014}. Certains types de documents sont utilisés dans différentes activités. Cela pourrait entraîner des modifications supplémentaires. Le guide d'utilisateur et le guide de déploiement font partie des types de documents les plus successibles d'être largement utilisés \cite{aghajani_doc_2020}. S'ajoutent également les duplicata d'informations dispersés à travers les documents, phénomène souvent causé par l'éparpillement des informations à travers des documents, comme c'est le cas dans les documents d’architecture \cite{rost_software_2013}. Ainsi, la maintenance de la cohérence devient d'autant plus difficile à maintenir dans le temps.
}{ED16, ED17, ED18}{ED16}

\enonce{ED16}{Assertion}{Plusieurs activités de différentes phases modifient des documents.}{}{P020}
\explication{ED16}{de l'assertion}{Plusieurs activités de différentes phases modifient des documents.}{
Dans son ouvrage \textit{Object-Oriented and Classical Software Engineering}, Stephen R. Schach consacre une section à nommer les raisons pour lesquels la documentation n'est pas dans sa propre phase, il conclut~:
\quoteenfr
{Therefore, just as there is no separate planning phase or testing phase, there is no separate documentation phase. Instead, planning, testing, and documentation should be activities that accompany all other activities while a software product is being constructed.}
{Par conséquent, de même qu'il n'existe pas de phase de planification ou de phase de test distincte, il n'existe pas de phase de documentation distincte. Au contraire, la planification, les tests et la documentation doivent être des activités qui accompagnent toutes les autres activités lors du développement d'un produit logiciel.}
{\ccite{schach_object_2011}[p.~17]}
Un sondage de 2014 chez 88~professionnels informatiques au sujet de la documentation (figure~B.7) indique que les documents sont modifiés, peu importe la phase.
}{}

\figannexe{
\figmediumen
{\textit{Development activities in which respondents write or adapt software documentation}}
{figure/chap1/documentation/ploschfig.png}
{(source~: \textit{Figure}~8 \cite{plosch_doc_2014})}
{}{}{tb}{Phases dans lesquelles des professionnels modifient la documentation}
}

\enonce{ED17}{Assertion}{Un même document peut être utilisé lors de plusieurs activités.}{}{P020}
\explication{ED17}{l'assertion}{Un même document peut être utilisé lors de plusieurs activités.}{
Un sondage de 2020 chez 146~professionnels informatiques au sujet de la documentation (figure~B.8) montre que certaines catégories de documents sont utilisées pour différentes activités de développement.
}{}

\figannexe{
\figen
{\textit{Types of documentation perceived as useful by practitioners in the context of specific software engineering tasks (RQ2), according to the results of Survey-II.}}
{figure/chap1/documentation/aghajani.png}
{(source~: \textit{Figure}~3 \cite{aghajani_doc_2020})}
{}{}{H}{Utilisation des documents par phase selon des professionnels}
}

\enonce{ED18}{Assertion}{Il y a un éparpillement des informations à travers les documents d'architecture.}{}{P020}
\explication{ED18}{l'assertion}{Il y a un éparpillement des informations à travers les documents d'architecture.}{
Un sondage réalisé en 2013 auprès de 147~professionnels informatiques créant la documentation d'architecture confirme le phénomène de dispersion des informations, qui entraîne des duplicata ainsi que des problèmes de mise à jour \cite{rost_software_2013} (figure~B.09).
}{}

\paragraphe{P0210}{Il est usuellement convenu que la documentation n'est pas une priorité pour plusieurs professionnels informatiques. Le contrôle du processus de documentation et le contrôle de la qualité de la documentation sont difficiles à obtenir.}{}{}

\subsubsection{Après la réalisation}

\paragraphe{P021}{Une fois une première itération du développement terminée, les artefacts (document, logiciel) continuent à être utilisés et nécessitent des adaptations. Cette sous-section débute en exposant deux faits qui mettent en valeur certains écueils sur la documentation. Ensuite, la sous-section présente ces écueils et en décrit les répercussions.}{}{}

% Mettre un sous partie de la dééfinition de l'OQLF.
\paragraphe{P022}{\fl{Il y a une grande proportion de logiciels patrimoniaux, ce qui augmente par conséquent l'importance de la documentation.} Des logiciels patrimoniaux sont issus «\elide d'une version dépassée et qui continuent d'être utilisés, après des ajustements, en même temps que la technologie contemporaine d'une entreprise.~» \footnote{Plusieurs autres définitions font intervenir par exemple le niveau de criticité ou de connaissances de l'entité informatique \cite{rosenkranz_explaining_2024}. Néanmoins, ils peuvent appartenir à la définition proposée.}. Une des raisons de cette grande proportion est que la fin de vie d'une entité informatique est rarement planifiée. Ainsi, le référentiel des connaissances en gestion de projet (PMBOK, \textit{Project Management Body of Knowledge} en anglais) à la version 6 \ccite{pmi_guide_2017}[p.~19] et à la version 7 \ccite{pmi_guide_2021}[p.~19] ainsi que dans l'\titleen{Information Technology Infrastructure Library}{}  \ccite{itil_2019}[p.~195] survolent le sujet en quelques lignes, mais ne l'intègre pas aux procédés de gestion.
\sep
En plus, une estimation provenant d'une revue de littérature réalisée en 2024 $(n=60)$ sur le sujet \cite{rosenkranz_explaining_2024} indique qu'entre 39~\% et 50~\% des logiciels, dont une majorité des logiciels critiques, utilisés dans des pays industrialisés sont des logiciels patrimoniaux. Sans documentation appropriée, la maintenance ou l'adaptation de ces logiciels peut s'avérer plus difficile, voire impossible. 
}{ED19, ED20, ED21}{ED19}

\enoncewithoutexp{ED19}{Définition}{Patrimonial}{
«~\elide issus d'une génération ou d'une version dépassée et qui continuent d'être utilisés \elide.~» \cite{Gdt_patrimonial_0}
}{P022}

\figannexe{%
\figmediumen
{}
{figure/chap1/documentation/across.png}
{(modifié la \textit{Figure}~7 \cite{rost_software_2013})}
{}{}{tb}{Fréquence de la répartition documentaire selon des professionnels}
}

\enonce{ED20}{Assertion}{La fin de vie des entités informatiques est faiblement expliquée à l'intérieur des ouvrages de référence.}{}{P019}
\explication{ED20}{de l'assertion}{La fin de vie des entités informatiques est faiblement expliquée à l'intérieur des ouvrages de référence.}{
Plusieurs bases de connaissances importantes abordent très peu la fin de vie. Voici en exemple, deux versions du \titleen{Project Management Body of Knowledge}{}(PMBOK) et de l'\textit{Information Technology Infrastructure Library} (ITIL) (le terme employé dans ces ouvrages est celui de \textit{retirement} \footnote{L'OQLF traduit le terme \emphase{abandon} par \emphase{retirement}. \cite{Gdt_abandon_0}})~:
\begin{itemize}
    \item Apparaît deux fois dans la version 6 du PMBOK (2017)~: \enfr{Project life cycles are independent of product life cycles, which may be produced by a project. A product life cycle is the series of phases that represent the evolution of a product, from concept through delivery, growth, maturity, and to retirement.}{Le cycle de vie d'un projet est indépendant du cycle de vie d'un produit, même si celui-ci peut être induit par un projet. Le cycle de vie d'un produit correspond à la succession des phases qui représentent son évolution, de sa conception à sa mise hors service, en passant par sa commercialisation, sa croissance, sa maturité et son retrait du marché.}
    \ccite{pmi_guide_2017}[p.~19];
    \item Apparaît six fois dans la version 7 du PMBOK (2021)~: une figure explicite que le \textit{retirement} est le dernier projet \ccite{pmi_guide_2021}[p.~19];
    \item Apparaît cinq fois dans l'\textit{Information Technology Infrastructure Library} (ITIL)~: \enfr{Unless the retirement of a product/service is carefully planned, it could cause unexpected negative effects on customers or the organization that might otherwise have been avoided.}{À moins que le retrait d'un produit ou d'un service ne soit soigneusement planifié, il pourrait entraîner des effets négatifs inattendus sur les clients ou l'organisation, effets qui auraient pu être évités autrement.}
    \ccite{itil_2019}[p.~195].
\end{itemize}

}{}

\enonce{ED21}{Assertion}{Dans plusieurs pays industrialisés, il est estimé qu'entre 39~\% et 50~\% des systèmes d'information utilisés sont patrimoniaux.}{}{P022}
\explication{ED21}{de l'assertion}{Dans plusieurs pays industrialisés, il est estimé qu'entre 39~\% et 50~\% des systèmes d'information utilisés sont patrimoniaux.}{
En 2024, une revue de littérature s'intéresse aux entités informatiques dites patrimoniales. La revue de littérature porte sur 60~articles publiés sur une période s'étalant de 1973 à 2020. Il est décrit que~:
\quoteenfr
{A vast number of studies show that the information system inventory of organizations in some Western countries ranges between around 39~\% [2] to around 50~\% [3], [4]. Therefore, a significant part of the IS in Western industrialized countries should be considered outdated. However, these systems are not only a widespread economic problem, but also a problem for society as a whole, as legacy systems are also a significant problem for public authorities [5], [6], [7], [8].}
{De nombreuses études montrent que l'inventaire des systèmes d'information des organisations dans certains pays occidentaux se situe entre environ 39~\% [2] et environ 50~\% [3], [4]. Par conséquent, une part importante des systèmes d'information dans les pays industrialisés occidentaux doit être considérée comme obsolète. Cependant, ces systèmes constituent non seulement un problème économique généralisé, mais aussi un problème pour la société dans son ensemble, car les systèmes hérités représentent également un problème important pour les autorités publiques [5], [6], [7], [8].}
{\cite{rosenkranz_explaining_2024}}
Cette revue cite trois études ([2], [3], [4]), en voici des détails pour les deux premiers\footnote{Le troisième étant en allemand, il a été décidé de ne pas en faire de résumé.}~:
\begin{itemize}
\item {}[2] Sondage de 2016 visant les systèmes d'administration des politiques utilisés par les assureurs $(n=58)$ au Canada et aux États-Unis~: \enfr{policy administration system (PAS)
\elide LIMRA and Deloitte surveyed companies in the summer of 2016 to learn more about their current core PAS and plans for modernization. Fifty-eight companies in the United States and Canada responded; one company responded for two areas of their company.
\elide
Roughly two in five companies (39 percent) surveyed have one legacy PAS.}
{Système de gestion des polices d'assurance (SGA)
\elide
LIMRA et Deloitte ont mené un sondage auprès d'entreprises au cours de l'été 2016 afin d'en savoir plus sur leurs SGA actuels
et sur leurs projets de modernisation. Cinquante-huit entreprises aux États-Unis et au Canada ont répondu ; une compagnie a répondu pour deux domaines.
\elide
Environ deux entreprises sur cinq (39~\%) sondées possèdent encore un SGA ancien.}
\cite{deloitte_legacy_2013}
\item {}[3] Étude de 2014 portant sur les applications $(n=4130)$ utilisées par l'État du Texas : 
\enfr
{The state’s application portfolio comprises more
than 4,130 business applications supported by
the state’s systems. Of these, approximately 58\%
are listed as legacy because they contain legacy
components. \elide
Approximately 36\% of the entire state applica-
tion portfolio is deemed mission critical by agen-
cies, while 64\% is not mission critical.}
{Le parc d'applications de l'État comprend
plus de 4~130 applications d'affaires prises en charge par les systèmes de l’État. Parmi celles-ci, environ 58~\% sont considérées comme obsolètes parce qu'elles contiennent des composants obsolètes.
\elide
Environ 36~\% de l’ensemble du parc applicatif de l’État est jugé critique par les agences,
alors que 64~\% ne l’est pas.}
\cite{texas_legacy_2014}
\end{itemize}

}{}

\paragraphe{P023}{\fl{Durant une partie importante de la vie d'un logiciel, la fréquence des changements aux entités informatiques est rapide, il en va de même pour la documentation.} Cette rapidité est un défi en soi. Ce phénomène est énoncé comme un constat par plusieurs experts en informatique. Barry W. Boehm indique que cela représente un défi \cite{Selby2007-chap10}. Ian F. Sommerville en donne les raisons suivantes~:
\quoteenfrlist{
    \item New approaches to systems development in particular, construction by configuration. \elide
    \item The need for rapid software delivery. \elide
    \item The increasing rate of requirements change. \elide
    \item The need for improved ROI on software assets. \elide
}{
    \item Nouvelles approches du développement de systèmes, en particulier la construction par configuration. \elide
    \item Le besoin d'une livraison rapide de logiciels. \elide
    \item Le rythme croissant de l'évolution des exigences. \elide
    \item Le besoin d'un meilleur retour sur investissement des actifs logiciels. \elide
}{\cite{sommerville_integrated_2005}}
Pareillement au code source, le professionnel informatique sera amené à faire plusieurs de modifications à la documentation. Ainsi, il doit continuer à comprendre les répercussions de ces changements et à maintenir le tout cohérent. 
}{ED22}{ED22}

\enonce{ED22}{Axiome}{Les logiciels doivent s'adapter plus fréquemment.}{}{P023}
\explication{ED22}{de l'axiome}{Les logiciels doivent s'adapter plus fréquemment.}{
En 2005, Ian F. Sommerville dans l'article \textit{Integrated requirements engineering: a tutorial} explique quatre raisons pour lesquelles les changements aux entités informatiques vont devenir plus fréquent~:
\quoteenfrlist{
\item [][---] New approaches to systems development in particular, construction by configuration.
\item [][---] The need for rapid software delivery.
\item [][---] The increasing rate of requirements change.
\item [][---] The need for improved ROI on software assets.
}{
\item [][---] Nouvelles approches du développement de systèmes, en particulier la construction par configuration.
\item [][---] Le besoin d'une livraison rapide de logiciels.
\item [][---] Le rythme croissant de l'évolution des exigences.
\item [][---] Le besoin d'un meilleur retour sur investissement des actifs logiciels.
}{\cite{sommerville_integrated_2005}}

En 2007, Barry W. Boehm décrit dans un article les épreuves à surmonter~:
\quoteenfr
{The challenges we have discussed here of simultaneously dealing with rapid change, uncertainty and emergency, dependability, diversity, and interdependence will often be formidable, but there will also be opportunities to make significant contributions that will make a difference for the better.}
{Les défis que nous avons évoqués ici, à savoir la gestion simultanée des changements rapides, de l'incertitude et des situations d'urgence, de la fiabilité, de la diversité et de l'interdépendance, seront souvent redoutables, mais il y aura aussi des occasions d'apporter des contributions significatives qui feront une différence positive.}
{\cite{Selby2007-chap10}}

}{}

\paragraphe{P024}{\fl{Les changements aux contextes et aux entités informatiques peuvent finir par se répercuter négativement sur la documentation.} D'abord, les acquis du lectorat changent et cela devient plus probable avec les entités informatiques patrimoniales. Cette situation peut rendre l'interprétation d'un document plus incertaine. 
\begin{itemize}
    \item Entre les années 1980 et la fin des années 1990, les changements sur le lectorat des documents informatique étaient suffisamment importants pour devoir en tenir compte, par exemple en considérant le niveau d'expertise en plus du champ d'expertise \cite{mehlenbacher_documentation_2003}.
    \item Entre les années 1950 et les années 1990, l'utilisation répandue du diagramme de flux est devenue celle du diagramme de flux de données, puis celle du diagramme d'activités \cite{morris_flow_2006}. 
\end{itemize}
Ensuite, les documents peuvent ne pas être mis à jour à la suite de changements. Plusieurs sondages menés auprès de professionnels informatiques indiquent qu'il est pratique courante de ne pas mettre à jour la documentation \cite{forward_relevance_2002, lethbridge_how_2003, rost_software_2013}. Si les professionnels informatiques ne disposent pas de temps suffisant pour correctement faire la documentation lors de la réalisation, il est vraisemblable que cela soit de même après la réalisation.
}{ED23, ED24}{ED22}

\enonce{ED23}{Assertion}{La documentation peut devenir obsolète.}{}{P024}
\explication{ED23}{de l'assertion}{La documentation peut devenir obsolète.}{
En 2002, un sondage interroge 48~professionnels informatiques sur la documentation \cite{forward_relevance_2002}. Les professionnels informatiques affirment que la documentation est rarement mise à jour. 

En 2003, un article analyse trois études au sujet de la documentation \cite{lethbridge_how_2003}. La troisième étude s'intéresse aux mises à jour de la documentation. La figure~B.10 montre la fréquence subjective de mise à jour des documents selon les sondés $(n=48)$. Une hypothèse est émise à la fin de l'article : les professionnels informatiques mettent l'effort sur les documents qu'ils estiment en valoir la peine. 

En 2013, un sondage avec 147~experts en conception indique que l'obsolescence est le problème principal des documents d'architecture \cite{rost_software_2013}. 
\figmediumen
{\textit{The time between system changes and documentation updates for different documentation types.}}
{figure/chap1/documentation/lethbridge.png}
{(source : \textit{Figure} 1 \cite{lethbridge_how_2003})}
{}
{
\begin{tablenotes}
\footnotesize
\item[a] {Dans l'article, le diagramme est en couleur, néanmoins les fréquences demeurent en ordre ce qui permet un affichage en gris.}
\end{tablenotes}
}{tb}{Délai de la mise à jour de la documentation selon des professionnels}
}{}

\enonce{ED24}{Assertion}{Le temps peut changer le lectorat, les langages utilisés et les standards.}{}{P024}
\explication{ED24}{de l'assertion}{Le temps peut changer le lectorat, les langages utilisés et les standards.}{
Le lectorat peut changer. Brad Mehlenbacher, dans un article, décrit différentes problématiques avec le lectorat de la documentation informatique \cite{mehlenbacher_documentation_2003}. L'expertise en informatique s'éparpille et se démocratise. Dans les années 1980, le lectorat était divisé ainsi~: \enfr{expertise with computer, expertise with task domain, expertise with application}{expertise avec les ordinateurs, expertise dans les tâches du domaine, expertise avec les applications}. En 1998, un autre axe est ajouté ; \enfr{novice, transfer, expert}{novice, intermédiaire, expert}. 

Il arrive que les représentations graphiques changent. En 2006, un article décrit l'historique d'une représentation graphique, le diagramme de flux \cite{morris_flow_2006}. Le titre est \titleen{Flow diagrams~: rise and fall of the first software engineering Notation}{Diagrammes de flux~: essor et déclin de la première notation du génie logiciel}. Le diagramme d'activités est considéré comme le successeur au diagramme de flux et au diagramme de flux de données. 
\begin{itemize}
    \item Avant 1950, le diagramme de flux est déjà utilisé pour décrire le fonctionnement des machines ou des algorithmes. Le diagramme de flux est décrit comme cela ; \enfr{The flow-diagram itself "expresses the ‘algebra’ of the process"}{Le diagramme de flux en lui-même "exprime l’‘algèbre’ du processus"}.
    \item Au milieu des années 1950, il y a eu un essor de ce type de diagramme, entre autres, à l'intérieur des guides d'utilisateur. Un déclin s'ensuivit.
    \item Dans les années 1970 et 1980, le développement structuré popularise le diagramme de flux de données.
    \item Dans les 1990, c'est au tour de l'\textit{Unified Modeling Language} (UML) de populariser des types de diagrammes. Il y a, par exemple, le diagramme de séquence, le diagramme de classes et le diagramme d'activités.
\end{itemize}

Les standards peuvent changer. L’ANSI/ANS-10.3-1995 \textit{Documentation of Computer Software} produit par l'American National Standards Institute (ANSI) et l'American Nuclear Society (ANS) est un exemple de standard qui a été abandonné sans successeur \cite{ansi_103_1995}.
}{}

\paragraphe{P025}{\fl{La documentation continue d'être négligée, même en sachant les répercussions.} Cela fait intervenir le concept de dette technique. Une dette technique est le passif technique résultant de choix opportunistes dont le coût de mise à niveau augmente avec le temps. Déjà lors d'un sondage ($ n=69$), les professionnels informatiques ont indiqué que la documentation peut contribuer à la dette technique \cite{holvitie_technical_2018}. Philippe Kruchten inclut la documentation comme une des multiples activités de développements contribuant à la dette technique \cite{kruchten_technical_2012}. Par ailleurs, une revue de littérature ($n=89$) \cite{tamburri_success_2021} conclut même que cette dette conduit à l'échec de l'évolution et de la maintenance du logiciel. Pour plusieurs professionnels informatiques, la documentation est considérée comme une source non fiable d'information et ils préfèrent utiliser le code source comme source d'informations \cite{kruger_memorize_2024, souza_study_2005}. 
\sep
Pourtant, l'une des particularités de plusieurs documents informatiques est que le professionnel informatique fait partie de l’autorat et du lectorat \cite{rettig_nobody_1991}. C'est qu'un auteur n’est pas sensibilisé aux lectorats, même si un  professionnel informatique peut faire partie à la fois de l'autorat et du lectorat \cite{jansen_enriching_2009}.
}{ED25, ED26, ED27, ED28}{ED25}

\enonce{ED25}{Définition}{Dette technique}{
Passif technique résultant de choix opportunistes dont le coût de mise à niveau augmente avec le temps.
}{P025}
\explication{ED25}{de la définition}{Dette technique}{
La définition formulée est~: ensemble des artefacts résultants de choix opportunistes dont le coût de mise à niveau augmente avec le temps.

La définition de dette technique de Steve McConnell se retrouve dans plusieurs articles \cite{kruchten_technical_2012, avgeriou_reducing_2016,avgeriou_reducing_2016}. Elle précise la dette technique comme étant intentionnelle~:
\quoteenfr
{The definition of technical debt was later revisited on a number of occasions, usually to generalize what Cunningham had previously described for all applicable situations whilst categorizing its characteristics. Steve McConnell’s definition, which separates out intentional and unintentional accumulation of technical debt [7], has been widely adopted by academia (e.g., in [8,9] and recognized in [5]).}
{La définition de la dette technique a été réexaminée à plusieurs reprises par la suite, généralement pour généraliser la description antérieure de Cunningham à toutes les situations applicables, tout en catégorisant ses caractéristiques. La définition de Steve McConnell, qui distingue l’accumulation intentionnelle et non intentionnelle de dette technique [7], a été largement adoptée par le monde universitaire (par exemple, dans [8,9] et reconnue dans [5]).}
{\cite{holvitie_technical_2018}}

La littérature fournit ces informations qui ont servi à construire cette définition~:
\begin{itemize}
    \item Définition de Steve McConnell~: \quoteenfr
{\elide a design or construction approach that is expedient in the short term but that creates a technical context in which the same work will cost more to do later than it would cost to do now.}
{une approche de conception ou de construction qui est opportune à court terme, mais qui crée un contexte technique dans lequel le même travail coûtera plus cher à réaliser plus tard qu'il ne coûterait à réaliser maintenant.}
{\cite{avgeriou_reducing_2016}}
    \item Définition de passif est~: «~Ensemble des dettes et des charges, évaluables en argent, qui grèvent un patrimoine.~» \cite{Larousse_passif_0}.
\end{itemize}

Le terme \emphase{dette technique} s'est imposé dans l'usage, mais le terme \emphase{passif technique} est vraisemblablement plus approprié.
}{}

\enonce{ED26}{Assertion}{La documentation peut contribuer à la dette technique menant à l'échec de l'évolution ou de la maintenance d'un logiciel.}{}{P025}
\explication{ED26}{de l'assertion}{La documentation peut contribuer à la dette technique menant à l'échec de l'évolution ou de la maintenance d'un logiciel.}{
En 2018, dans un sondage mené auprès de 184~professionnels informatiques (figure~B.11) révèle notamment que la documentation inadéquate est corrélée à l'accroissement de la dette technique \cite{holvitie_technical_2018}. 
\figmediumen
{\textit{Causes of concrete technical debt (based on a subset from total answers}, $N_i = 69$)}
{figure/chap1/documentation/holvitie.png}
{(source~: \textit{Figure}~4 \cite{holvitie_technical_2018})}
{}{}{tb}{Causes de la dette technique selon des professionnels}

Philippe Kruchten (ancien directeur du RUP à Rational Rose) \textit{et al.} incluent la documentation comme l'une des multiples activités de développement contribuant à la dette technique~:
\quoteenfr
{Architecture plays a significant role in the development of large systems, together with other development activities, such as documentation and testing (which are often lacking). These activities can add significantly to the debt and thus are part of the technical debt landscape.}
{L'architecture joue un rôle important dans le développement des grands systèmes, de même que d'autres activités de développement, telles que la documentation et les tests (souvent négligés). Ces activités peuvent considérablement accroître la dette technique et font donc partie intégrante de ce phénomène.}
{\cite{kruchten_technical_2012}}

Publiée en 2021, une revue de littérature analyse les critères de succès ou d'échec en génie logiciel. La revue de littérature porte sur l'analyse de 89~articles publiés sur une période s'échelonnant de 1970 à 2019. Un des résultats observés est que~:
\quoteenfr
{Documentation quality: From the perspective of software maintenance and evolution, documentation is a discriminant in successful or failing software projects [32]–[34].}
{Qualité de la documentation~: Du point de vue de la maintenance et de l'évolution des logiciels, la documentation est un facteur discriminant dans les projets logiciels réussis ou en échec [32]–[34].}
{\cite{tamburri_success_2021}}

}{}

\enonce{ED27}{Assertion}{Plusieurs professionnels informatiques considèrent la documentation comme une source non fiable d'informations.}{}{P025}
\explication{ED27}{l'assertion}{Plusieurs professionnels informatiques considèrent la documentation comme une source non fiable d'informations.}{
En 2005, 76~professionnels informatiques s'occupant de la maintenance sont sondés \cite{souza_study_2005}. Selon eux, le code source a beaucoup plus d'importance que de nombreux autres artefacts qui sont généralement compris à l'intérieur de la documentation, les résultats sont dans la figure~B.12.
En 2024, un sondage est mené sur 37~professionnels informatiques à propos de la disponibilité de la connaissance \cite{kruger_memorize_2024}. Les sources de connaissances les plus précieuses, classées en ordre d'importance~: \sfren{connaissances propres}{own knowledge}, \sfren{connaissance des collaborateur}{knowledge of collaborators}, \sfren{code source}{source code}, \sfren{documentation}{documentation}, \sfren{système de contrôle de version}{version control system}, \sfren{outils d'analyse}{analysis tools}. 
}{}

\figannexe{
\figgm
{Importance des artefacts selon des professionnels}
{figure/chap1/documentation/souza.png}
{(source~: \textit{Table} 1 \cite{souza_study_2005})}
{}{}{H}
%The second to last column give the percentage of “very important” notes among maintainers that knew the documentation artifact (i.e. total minus last column).
}

\enonce{ED28}{Assertion}{Plusieurs professionnels informatiques ne se préoccupent pas du lectorat, bien qu'ils fassent partie eux-mêmes du lectorat.}{}{P025}
\explication{ED28}{de l'assertion}{Plusieurs professionnels informatiques ne se préoccupent pas du lectorat, bien qu'ils fassent partie eux-mêmes du lectorat.}{
En 2009, une expérimentation est menée sur l'utilisation d'un instrument dont l'objectif est d'améliorer la documentation d'architecture \cite{jansen_enriching_2009}. Dans les limitations énoncées, il y a le découragement des auteurs. Le bénéfice à obtenir une bonne documentation profite surtout aux lectorats, et à peu près pas aux auteurs. 

En 1991, un article d'opinion y va d'un titre provocateur \titleen{Nobody Reads Documentation}{Personne ne lit la documentation} \cite{rettig_nobody_1991}. L'auteur Marc Rettig affirme que les professionnels informatiques se plaignent de la documentation et qu'ils ne peuvent en vouloir qu'à eux-mêmes. Il écrit \enfr{Physician, heal thyself}{Médecin, guéris-toi toi-même}. 

En 2012, Qian Hu décrit les améliorations qu'il a apportées à son ours de rédaction \cite{qian_software_2012}. Son objectif principal était de faire comprendre aux étudiants en informatique qu'ils étaient incapables de maintenir un logiciel sans la documentation. Ses étudiants ne disposaient que du code source et des tâches de maintenance et d'évolutivité que Qian Hu leur avait attribuées. Par la suite, les étudiants devaient s'entendre sur la documentation manquante et l'écrire. En expérimentant le rôle du lectorat, ils furent mieux à même de comprendre l'importance de la documentation.
}{}

\paragraphe{P0250}{
Le contexte actuel comportant une grande proportion de logiciels patrimoniaux et une fréquence élevé de changements rend la documentation d'autant plus importante. Car, la documentation s’avère être une source d’information pour le professionnel informatique qui devra maintenir ou adapter les logiciels. Néanmoins, cette documentation doit demeurer de qualité.
}{}{}

\subsection{Des approches}

\paragraphe{P026}{Parmi les approches, il y a des instruments, des méthodes et des techniques. Il existe quelques propositions d'instruments qui combinent méthodes et techniques, elles sont présentées au début. Cependant, la majorité des instruments utilisés présentent certaines difficultés sur l'accessibilité et les fonctions offertes. Ainsi, les méthodes et les techniques sont là pour expliquer le fonctionnement que pourraient avoir des instruments pour la documentation informatique. Les méthodes proviennent de l'informatique. Les techniques proviennent de l'informatique, du traitement automatique des langues et de l'apprentissage profond.}{}{}

\paragraphe{P027}{Voici les définitions d'instrument, de méthode et de technique~:
\begin{itemize}
    \item instrument~: «~Objet fabriqué servant à exécuter [quelque chose], à faire une opération.~» \cite{Robert_instrument_0};
    \item méthode~: «~Ensemble de démarches raisonnées, suivies pour parvenir à un but.~»  \cite{Robert_methode_0};
    \item technique~: «~Ensemble des applications de la science dans le domaine de la production.~» \cite{Larousse_technique_0}.
\end{itemize}
\vspace{-1\baselineskip}
}{AD01, AD02, AD03}{AD01}

\enoncewithoutexp{AD01}{Définition}{Instrument}{
Objet fabriqué servant à exécuter [quelque chose], à faire une opération.\cite{Robert_instrument_0} 
}{P027}

\enoncewithoutexp{AD02}{Définition}{Méthode}{
Ensemble de démarches raisonnées, suivies pour parvenir à un but. \cite{Robert_methode_0}
}{P027}

\enoncewithoutexp{AD03}{Définition}{Technique}{
Ensemble des applications de la science dans le domaine de la production.\cite{Larousse_technique_0}
}{P027}

\subsubsection{Des instruments}

\paragraphe{P028}{En premier sont présentées des propositions de solutions pour maintenir cohérent la documentation. S'ensuit une présentation des principaux instruments utilisés en documentation qui pointent vers certaines difficultés d'accessibilité.}{}{}

\paragraphe{P029}{\fl{Il existe des propositions de logiciels qui combinent des techniques informatiques visant à améliorer la qualité de la documentation.} Un article, publié en 2002, intitulé \titleen{Lightweight validation of natural language requirements}{} propose un prototype qui regroupe langage formel et modélisation informatique \cite{gervasi_lightweight_2002}. D'une façon sommaire, voici les huit étapes comprenant la préparation et l'exécution de la solution~:
\begin{enumerate}
    \item Mise en place d'un ensemble de règles pour le langage appelé aussi grammaire.
    \item Mise en place d'un ensemble de règles pour les propriétés du modèle.
    \item Mise en place d'un ensemble de règles pour extraire le modèle.
    \item Prétraitement sur le texte.
    \item Analyse lexicale et syntaxique du texte grâce à la grammaire.
    \item Transformation de l'arbre syntaxique en modèle à partir des règles d'extraction,
    \item Application des règles pour les propriétés du modèle.
    \item Validation et correction du document par l'utilisateur, puis retour à l'étape 4.
\end{enumerate}
Dans ce procédé, c'est l'utilisateur qui détermine si le document est valide ou non et qui décide s'il a besoin d'itérer. Ainsi, il demeure en contrôle. À propos de leur approche, les auteurs affirment~:
\quoteenfr
{The low cost of our approach, both in terms of human training and of computational resources needed, makes it particularly well suited for the initial introduction of formal methods in an organization.}
{Le faible coût de notre approche, tant en termes de formation humaine que de ressources informatiques nécessaires, la rend particulièrement adaptée à l'introduction initiale de méthodes formelles dans une organisation.}
{\cite{gervasi_lightweight_2002}}
\vspace{-1\baselineskip}
}{AD04}{AD04}

\enonce{AD04}{Assertion}{Il existe des propositions de logiciels qui combinent des techniques informatiques visant à améliorer la qualité de la documentation.}{}{P029}
\explication{AD04}{l'assertion}{Il existe des propositions de logiciels qui combinent des techniques informatiques visant à améliorer la qualité de la documentation.}{
Dans un article publié en 2002, Vincenzo Gervasi1 et Bashar Nuseibeh proposent un prototype visant à valider les documents d'exigences. Le diagramme dans la figure~B.13 décrit le procédé \cite{gervasi_lightweight_2002}. 
\figmsen
{\textit{Set-up phase and production phase in our approach.}}
{figure/chap1/documentation/processlw.png}
{(source~: \textit{Figure}~1 \cite{gervasi_lightweight_2002})}
{}{}{H}{Procédé de formalisation d'un document}
Ensuite, les auteurs décrivent chacune des huit étapes présentées dans la figure~B.13~:
\quoteenfrlist {
    \item[][---] Defining a style a structure and a language for the requirements document. This step can be undertaken either normatively, i.e. as the production of a prescriptive style manual for the requirements document (and in this case a syntax-guided editor can be used to support requirements writing, as in [16]), or descriptively, i.e. as an adaptation of the capabilities of a parsing tool to an already existing document written in a defined style (as in the case of the experience we report in this paper).
    \item[][---] Selecting desirable properties to check. Which properties of a certain document or system described in a document are ‘interesting’ depends on the particular context of the analysis. As is common with lightweight formal methods, partial validation is usually acceptable at this stage.
    \item[][---] Defining one or more models against which the properties selected in the previous step can be checked. Properties are always relative to models, i.e. abstractions of the document or of the system described in it, which collect in an analysable structure the information needed to check the property. For example, a connection property among system components can be checked against a model describing all the communications among system components.
    \item[][---] Pre-processing the requirements document, to handle format, structure and typographical details, and to translate the requirements document to a canonical form amenable to later processing.
    \item[][---] Parsing the NL text of the requirements, leading to an analysable representation of the semantic content of the text. Again, parsing can be (and usually is) partial, to help reduce the cost of validation, as long as this does not interfere with the collection of the information needed to perform the validation.
    \item[][---] Building the models defined in step 3 above, using the information collected during the parsing process. It is possible to build models of the requirements document (for example, distribution of topics among sections of the document) and of the system described by the requirements (for example, a model of the communication paths in a distributed system).
    \item[][---] Checking that the models satisfy chosen properties. As in the previous step, it is possible to check properties of the document (in our previous example, consistency of topics inside a single section) and of the system (for example, the existence of disjoint components in the communication paths model).
    \item[][---] Evaluating findings and revising the requirements specification accordingly. It is particularly important that the validation checks provide as much detail as possible about the point of and the reason for a failure (i.e. about the circumstances in which a validation property was violated). Like the counterexamples provided by other formal methods, this information helps the requirements engineer to identify and fix errors that cause violations.
} {
    \item[][---] Définition d'un style, d'une structure et d'un langage pour le document d'exigences. Cette étape peut être menée de manière normative, c'est-à-dire par la production d'un guide de style prescriptif pour le document d'exigences (et dans ce cas, un éditeur syntaxique peut être utilisé pour faciliter la rédaction des exigences, comme dans [16]), ou de manière descriptive, c'est-à-dire par l'adaptation des capacités d'un outil d'analyse syntaxique à un document existant rédigé dans un style défini (comme dans le cas de l'expérience que nous rapportons dans cet article).
    \item[][---] Sélection des propriétés pertinentes à vérifier. Les propriétés «~intéressantes~» d'un document ou d'un système décrit dans un document dépendent du contexte spécifique de l'analyse. Comme c'est souvent le cas avec les méthodes formelles légères, une validation partielle est généralement acceptable à ce stade.
    \item[][---] Définition d'un ou plusieurs modèles permettant de vérifier les propriétés sélectionnées à l'étape précédente. Les propriétés sont toujours relatives à des modèles, c'est-à-dire des abstractions du document ou du système qu'il décrit, qui regroupent dans une structure analysable les informations nécessaires à la vérification de la propriété. Par exemple, une propriété de connexion entre les composants du système peut être vérifiée par rapport à un modèle décrivant toutes les communications entre ces composants.
    \item[][---] Prétraitement du document d'exigences~: gestion du format, de la structure et des détails typographiques, et traduction du document d'exigences dans un format canonique adapté au traitement ultérieur.
    \item[][---] Analyse syntaxique du texte en langage naturel des exigences, aboutissant à une représentation analysable du contenu sémantique du texte. Là encore, l'analyse syntaxique peut être (et est généralement) partielle, afin de réduire le coût de la validation, à condition que cela n'entrave pas la collecte des informations nécessaires à cette validation.
    \item[][---] Construction des modèles définis à l'étape 3 ci-dessus, à partir des informations collectées lors de l'analyse syntaxique. Il est possible de construire des modèles du document d'exigences (par exemple, la répartition des sujets entre les sections du document) et du système décrit par les exigences (par exemple, un modèle des chemins de communication dans un système distribué).
    \item[][---] Vérification que les modèles satisfont aux propriétés choisies. Comme à l'étape précédente, il est possible de vérifier les propriétés du document (dans notre exemple précédent, la cohérence des sujets au sein d'une même section) et du système (par exemple, l'existence de composants disjoints dans le modèle des chemins de communication).
    \item[][---] Évaluation des résultats et révision du cahier des charges en conséquence. Il est particulièrement important que les contrôles de validation fournissent autant de détails que possible sur le point et la raison d'une défaillance (c'est-à-dire sur les circonstances dans lesquelles une propriété de validation a été violée). À l'instar des contre-exemples fournis par d'autres méthodes formelles, ces informations aident l'ingénieur en exigences à identifier et à corriger les erreurs qui entraînent des violations.
}{\cite{gervasi_lightweight_2002}}
\vspace{-1\baselineskip}
}{}

\paragraphe{P030}{\fl{Dans les fonctionnalités demandées, plusieurs professionnels informatiques veulent que les instruments gèrent les relations entre les documents et à l'intérieur de ceux-ci.} L'objectif étant que, même lors des mises à jour, la cohérence soit maintenue. Il est même suggéré, dans une revue de littérature sur la documentation, que les documents deviennent exécutables \cite{theunissen_mapping_2022}. De cette façon, il n'y aurait plus de désynchronisation et elle serait vérifiable avec des tests. Un autre article propose différentes solutions pour y arriver (tableau 1.3) et comprend méthodes et techniques \cite{correia_patterns_2009}. 
\tab
{Sommaire des problèmes et des solutions}
{figure/chap1/documentation/instrument sol-fr.png}
{[Traduction libre] (inspiré de \cite{correia_patterns_2009})}
{}
{
\begin{tablenotes}
\footnotesize
\item[a] {Il peut exister plusieurs solutions pour un même problème, elles sont appelées variantes et peuvent s'exécuter sous des conditions similaires ou différentes.}
\end{tablenotes}
}{tb}
}{AD05}{AD05}

\figannexe{
\tab
{Sommaire des problèmes et solutions en anglais}
{figure/chap1/documentation/instrument sol.png}
{(inspiré de \cite{correia_patterns_2009})}
{}
{}{tb}
}
\enonce{AD05}{Assertion}{Un instrument de documentation devrait maintenir les relations dans les documents.}{}{P030}
\explication{AD05}{l'assertion}{Un instrument de documentation devrait maintenir les relations dans les documents.}{
En 2009, l'article \titleen{Patterns for consistent software documentation} fait l'exercice de décrire ce qu'un instrument devrait utiliser pour éviter les contradictions dans les mises à jour de la documentation \cite{correia_patterns_2009}. Le tableau~B.9 regroupe les solutions proposées par problème.

En 2022, une revue de littérature suggère comme caractéristique que la documentation devrait être exécutable \cite{theunissen_mapping_2022}. La revue porte sur 63~articles s'étalant de 2001 à 2018 dans le contexte des méthodes agiles. Trois avantages de cette caractéristique sont mis en évidence~: 
\quoteenfrlist {
    \item[][---] Executable documentation is never out-of-sync, it’s evolution is naturally connected to the evolution of the other parts of the software. Executable documentation does not just describe the software, but it is part of the software.
    \item[][---] Executable documentation can be tested. If it does not lead to the desired results, then something must be wrong. In that respect, executable documentation is just like source-code.
    \item[][---] The process of creating executable documentation is not intrusive, i.e. developers do not stop their work to take care of documentation; coding and documenting are part of the same task.
} {
\item[][---] La documentation exécutable est toujours synchronisée ; son évolution est naturellement liée à celle des autres composants du logiciel. Elle ne se contente pas de décrire le logiciel, elle en fait partie intégrante.
\item[][---] La documentation exécutable peut être testée. Si elle ne produit pas les résultats escomptés, c’est qu’il y a un problème. À cet égard, elle est comparable au code source.
\item[][---] Le processus de création de la documentation exécutable n’est pas intrusif~: les développeurs n’interrompent pas leur travail pour s’en occuper. Coder et documenter font partie intégrante de la même tâche.
}{\cite{theunissen_mapping_2022}}
\vspace{-1\baselineskip}
}{}

\paragraphe{P031}{\fl{Les détails sur le fonctionnement des instruments peuvent être inaccessibles.} D'abord, un article de 2002 relève lors d'un sondage $(n=60)$ les instruments de documentation utilisés par des professionnels informatiques \cite{forward_relevance_2002}. Plus récemment, une revue de littérature $(n=63)$ de 2022 relève aussi les instruments abordés dans les articles sur lesquels portent cette même revue de littérature. Le tableau 1.4 compare les deux relevés d'instruments. À l'intérieur figurent les instruments Word et Confluence qui sont très utilisés et demeurent des solutions propriétaires. 
\sep
Ensuite, dans les éléments qui peuvent aider à comprendre le fonctionnement d'un instrument, il y a la documentation (à la condition qu'elle soit de qualité [à jour, cohérente, etc.]) et le code source. Ces deux artefacts peuvent venir avec des conditions d'accessibilité. Pour un utilisateur d'une solution propriétaire, il arrive que la documentation nécessite un coût supplémentaire, tandis que le code source lui sera rarement accessible.
\tab
{Relevés des catégories d'instruments et leurs fréquences.}
{figure/chap1/documentation/instrument.png}
{(inspiré de \cite{forward_relevance_2002, theunissen_mapping_2022})}
{}
{
\begin{tablenotes}
\footnotesize
\item[a] {Des exemples sont fournis entre parenthèses.}
\item[b] {La fréquence d'apparition de la catégorie d'instrument dans les articles est présentée entre chevrons.}
\end{tablenotes}
}{tb} 
}{AD06}{AD06}

\enonce{AD06}{Assertion}{Entre les années 2002 et 2018, les instruments restent relativement similaires, à l'exception de l'utilisation des sites wiki.}{}{P031}
\explication{AD06}{l'assertion}{Entre les années 2002 et 2018, les instruments restent relativement similaires, à l'exception del'utilisation des sites wiki.}{
En 2022, une revue de littérature porte sur 63~articles dans le contexte des méthodes agiles \cite{theunissen_mapping_2022}. Les instruments utilisés sont relevés à l'intérieur d'articles publiés entre 2004 et 2018. Leurs descriptions et la fréquence d'articles correspondants entre parenthèses sont données ici~:
\begin{itemize}
    \item \textit{Word, Excel, Powerpoint etc.~: Used for reports and intended for long term usage} (40);
    \item \textit{Wiki~: Wiki-likes, Confluence} (39);
    \item \textit{Source control~: Tools for Git, Mercurial, SVN} (5);
    \item \textit{Scripts~: Executable lines of code to run tasks} (9);
    \item \textit{Markdown editors~: Light weight text editor, often used without editor but edited directly} (2);
    \item \textit{Verbal communication~: The proverbial water-cooler conversation} (7).
\end{itemize}

Un sondage de 2002 est réalisé auprès de 60~professionnels informatiques au sujet de la documentation \cite{forward_relevance_2002}. Les instruments utilisés par les participants sont relevés, les voici avec la fréquence~:
\begin{itemize}
    \item \textit{MS Word (and other word processors)} (22);
    \item \textit{Javadoc and similar tools (Doxygen, Doc++)} (21);
    \item \textit{Text Editors} (9);
    \item \textit{Rational Rose} (5);
    \item \textit{Together (Control Centre, IDE)} (3).
\end{itemize}
\vspace{-1\baselineskip}
}{}

\subsubsection{Des méthodes informatiques}

\paragraphe{P032}{Les méthodes en informatique peuvent comprendre des méthodes informatiques de développement et des méthodes informatiques de documentation. Toutefois, certaines méthodes peuvent emprunter des techniques à d'autres méthodes. C'est le cas des deux méthodes suivantes~: \emphase{Software design documentation} et \emphase{Literate programming}. 
}{}{}

\paragraphe{P033}{Ces méthodes nécessitent de définir les concepts suivants~:
\begin{itemize}
    \item organiser~: «~Doter d'une structure, d'une constitution déterminée, d'un mode de fonctionnement.~» \cite{Robert_organiser_0};
    \item module~: est un composant comprenant une interface; 
    \item couplage~: mesure l'interdépendance entre des composants logiciels;
    \item adaptabilité~: la capacité d'une entité à être modifié afin de répondre à des changements d'environnements ou de besoins.
\end{itemize}
Il est utilisé le terme \emphase{cohérence} à \emphase{cohésion} dans le mémoire. Il s'est avéré que le concept de cohésion comme celui de couplage possède une grande variabilité en informatique appliquée \cite{chami_cc_2008}. Lorsque le terme cohésion est décrit, il fait souvent référence à l'utilisation d'une logique (séquentielle, temporelle, fonctionnelle, etc.). Ainsi, le terme cohérence est plus approprié que cohésion. D'ailleurs, Bertrand Meyer indique que la cohérence est la caractéristique désirée et résultant d'une encapsulation appropriée \ccite{meyer_oo_1997}[p.~191].
}{AD07, AD08, AD09, AD10, AD11}{AD07}

\enoncewithoutexp{AD07}{Définition}{Organiser}{
Doter d'une structure, d'une constitution déterminée, d'un mode de fonctionnement. \cite{Robert_organiser_0}
}{P033}

\enonce{AD08}{Définition}{Module}{
Est un composant comprenant une interface. 
}{P033}
\explication{AD08}{de la définition}{Module}{
La définition formulée est~: est un composant comprenant une interface. 

Voici des descriptions qui proviennent de la norme ISO 24765\hc2017 et de l'ouvrage \textit{Documenting software architectures: views and beyond} de Paul Clements~:
\begin{itemize}
\item description 1~: 
\enfr{Note 1 to entry: The terms ‘module’, ‘component,’ and ‘unit’ are often used interchangeably or defined to be subelements of one another in different ways depending upon the context. The relationship of these terms is not yet standardized.}{Note 1~: Les termes «~module~», «~composant~» et «~unité~» sont souvent employés indifféremment ou définis comme des sous-éléments les uns des autres de différentes manières selon le contexte. La relation entre ces termes n’est pas encore normalisée.~»}
\ccite{iso_24765_2017}[p.~281];

\item description 2~: \enfr{A module refers first and foremost to a unit of implementation. Parnas’s foundational work in module design (Parnas 1972) used information hiding as the criterion for allocating responsibility to a module. Information that was likely to change over the lifetime of a system, such as the choice of data structures or algorithms, was assigned to a module, which had an interface through which its facilities were accessed.}{Un module désigne avant tout une unité d'implémentation. Les travaux fondateurs de Parnas sur la conception modulaire (Parnas, 1972) ont utilisé le principe d'encapsulation comme critère d'attribution des responsabilités à un module. Les informations susceptibles d'évoluer au cours du cycle de vie d'un système, telles que le choix des structures de données ou des algorithmes, étaient attribuées à un module, lequel disposait d'une interface permettant d'accéder à ses fonctionnalités.} 
\ccite{clements_documenting_2014}[p.~29-31];

\item description 3~: \enfr{A component refers to a runtime entity. Szyperski says that a component «~can be deployed independently and is subject to composition by third parties~» (Szyperski 1998, p.~30). The emphasis is clearly on the finished product and not on the implementation considerations that went into it.}{Un composant désigne une entité d'exécution. Szyperski précise qu'un composant «~peut être déployé indépendamment et est susceptible d'être composé par des tiers~» (Szyperski, 1998, p.~30). L'accent est clairement mis sur le produit fini et non sur les considérations d'implémentation qui ont contribué à sa réalisation.}
\ccite{clements_documenting_2014}[p.~29-31].
\end{itemize}
\vspace{-1\baselineskip}
}{}

\enonce{AD09}{Assertion}{Le terme cohérence est préférable à celui de cohésion pour discuter de modularité.}{}{P033}
\explication{AD09}{de l'assertion}{Le terme cohérence est préférable à celui de cohésion pour discuter de modularité.}{
Avant tout, les définitions de cohésion et de couplage apportent de la confusion en informatique, une analyse de l'évolution de ces concepts datant de 2008 soulève l'ambiguïté existant~:
\quotefr 
{Le grand nombre de recherches effectuées sur le couple C\&C [(cohésion et couplage)], et ayant comme objectif général l'éclaircissement de ces deux concepts, est une preuve de l'utilité et du grand emploi du couple C\&C en GL. Reste à savoir si cette utilisation est efficace ou non. Nous pensons qu'elle ne l'est pas vraiment en raison de l'ambiguïté qui existe encore autour de ces deux concepts.}
{\ccite{chami_cc_2008}[p.~132]}

Ensuite, les différentes définitions de cohésion laissent transparaître qu'il s'agit d'une mesure appliquée selon une logique déterminée~:
\begin{itemize}
    \item définition 1 de cohésion~: «~Propriété d'un ensemble dont toutes les parties sont solidaires \elide ~»
    \cite{Larousse_cohesion_0};
    \item définition 2 de cohésion~: «~Caractère d'une pensée, d'un exposé, etc., dont toutes les parties sont liées logiquement les unes aux autres \elide ~»
    \cite{Larousse_cohesion_0};
    \item définition 3 de cohésion~: \enfr{\elide in software design, a measure of the strength of association of the elements within a module \elide.}{\elide en conception logicielle, une mesure de la force d'association des éléments au sein d'un module \elide.}
    \cite{iso_26514_2021};
    \item définition 4 de cohésion~:
    \enfr{Cohesion of a component means that its functions and data are semantically related (they relate to a common purpose), instead of being arbitrary collections of elements.}{La cohésion d'un composant signifie que ses fonctions et ses données sont liées sémantiquement (elles se rapportent à un objectif commun), au lieu d'être des collections arbitraires d'éléments.}
    \ccite{lano_arch_2023}[p.~24];
    \item description de cohésion~: \enfr{Cohesion
    measures how strongly the responsibilities of a module are related.}{La cohésion mesure le degré de cohésion entre les responsabilités d'un module.}
    \ccite{bass_software_2021}[p.136].
\end{itemize}
Pour finir, Bertrand Meyer, dans son livre \textit{Object-Oriented  Software Construction}, affirme qu'un des critères souhaités pour la modularité est la cohérence~:
\quoteenfr
{\elide we know from earlier discussion how important information hiding is to the design of coherent and flexible architectures.
}{\elide nous savons, d'après une discussion précédente, à quel point l'encapsulation est importante pour la conception d'architectures cohérentes et flexibles.}
{\ccite{meyer_oo_1997}[p.~191]}
}{}

\enonce{AD10}{Définition}{Couplage}{Mesure l'interdépendance entre des composants logiciels.}{P033}
\explication{AD10}{de la définition}{Couplage}{
La littérature fournit ces informations qui ont servi à construire cette définition~:
\begin{itemize}
    \item définition de coupler~: «~Assembler des choses, les réunir, les relier afin de constituer un seul élément.~»
    \cite{Larousse_coupler_0};
    \item description 1 de coupler~: 
    \enfr{\elide in software design, a measure of the interdependence among modules in a computer program \elide.}{\elide en conception logicielle, une mesure de l'interdépendance entre les modules d'un programme informatique \elide.}
    \cite{iso_26514_2021};
    \item description 2 de coupler~: 
    \enfr{We can quantify this overlap by measuring the probability that a modification to one module will propagate to the other. This relationship is called coupling \elide.}{Nous pouvons quantifier ce chevauchement en mesurant la probabilité qu'une modification apportée à un module se propage à l'autre. Cette relation est appelée couplage \elide.}
    \ccite{bass_software_2021}[p.136].
\end{itemize}

}{}

\enonce{AD11}{Définition}{Adaptabilité}{La capacité d'une entité à être modifié afin de répondre à des changements d'environnements ou de besoins.}{P033}
\explication{AD11}{de la définition}{Adaptabilité}{
La définition formulée est~: la capacité d'une entité à être modifié afin de répondre à des changements d'environnements ou de besoins.

La littérature fournit ces informations qui ont servi à construire cette définition~:
\begin{itemize}
    \item définition d’évolution~: «~Ensemble de ces modifications, stade atteint dans ce processus, considéré comme un progrès \elide.~» \cite{Larousse_evolution_0};
    \item \sfren{la modifiabilité}{modifiability}~: \enfr{Modifiability is about change, and our interest in it is to lower the cost and risk of making changes.}{La modifiabilité est liée au changement, et notre intérêt pour elle est de réduire le coût et le risque des changements.}
    \ccite{bass_software_2021}[p.~117];
    \item \sfren{la maintenabilité}{maintainability}~: \enfr{The capability of the software product to be modified. Modifications may include corrections, improvements or adaptation of the software to changes in environment, and in requirements and functional specifications.}{La capacité du produit logiciel à être modifié. Les modifications peuvent inclure des corrections, des améliorations ou une adaptation du logiciel aux changements d'environnement, d'exigences et de spécifications fonctionnelles.}
    \ccite{iso_9126part1_2001}[p.~10];
    \item définition d'adaptatif~: «~Qui facilite ou constitue une adaptation.~» \cite{Antidote_0};
    \item définition de \emphase{s’adapter à}~: «~s’habituer à (un environnement) en modifiant son comportement en fonction de (cet environnement).~» \cite{Antidote_0};
    \item définition d'adaptabilité~: «~Caractère de ce qui est adaptable.~» \cite{Antidote_0};
    \item définition d'adaptable~: «~Qui peut s’adapter, que l’on peut adapter.~» \cite{Antidote_0}
\end{itemize}
L'évolution est souvent associée à l’idée de progrès. Cependant, un changement à un logiciel n'est pas forcément un progrès. 
}{}

\paragraphe{P034}{\fl{Une méthode est d'utiliser les techniques qui améliorent l'adaptabilité des documents comme il est possible de le faire pour le logiciel.} D'abord, la programmation structurée et l'encapsulation avec la spécification des antécédents et des conséquents permettent de diminuer le couplage et de maintenir la cohérence et, ainsi d'obtenir une modularité qui améliore l'adaptabilité d'un logiciel \ccite{bass_software_2021}[p.~120]. 
\sep
Un article de S. D. Hester et David L. Parnas présente une méthode, la \emphase{Software design documentation}, qui applique ces techniques à la documentation \cite{hester_using_1981}. Constatant la faible adhésion à cette méthode, David L. Parnas, Paul Clements et David Weiss proposent un exemple plus réaliste en s'appuyant sur le Operational Flight Program de l'avion A-7E \cite{parnas_modular_1985}. 
\sep
Néanmoins, même lors de l'utilisation de langage formel, ces techniques sont peu efficacement  utilisées. Une analyse réalisée sur cent logiciels en code source libre sur GitHub détermine qu'il faudrait un effort important pour diminuer le couplage et pour augmenter la cohérence \cite{candela_using_2016}.
}{AD12, AD13, AD14}{AD12}

\enonce{AD12}{Assertion}{La programmation structurée et l'encapsulation permettent d'améliorer l'adaptabilité d'un logiciel.}{}{P034}
\explication{AD12}{de l'assertion}{La programmation structurée et l'encapsulation permettent d'améliorer l'adaptabilité d'un logiciel.}{
Dans leur ouvrage \textit{Software Architecture in Practice}, Len Bass, Paul Clements et Rick Kazman explicitent, à l'aide d'un diagramme (figure~B.14), comment accroître la cohérence, diminuer la couplage ou différer la liaison afin d'améliorer l'adaptabilité \ccite{bass_software_2021}[p.~120].
\fig
{Approches pour augmenter l'adaptabilité}
{figure/chap1/documentation/tactic.png}
{(source~: \textit{Figure}~8.3 \cite{bass_software_2021})}
{}{}{H}
Ces techniques proviennent de la programmation structurée et de l'encapsulation~:
\quoteenfr
{Edsger Dijkstra’s 1968 paper on the T.H.E. operating system introduced the concept of layers [Dijkstra 68]. The early work of David Parnas laid many conceptual foundations, including information hiding [Parnas 72], program families [Parnas 76], the structures inherent in software systems [Parnas 74], and the uses structure to build subsets and supersets of systems [Parnas 79]. \elide
In 1972, Dijkstra and Hoare, along with Ole-Johan Dahl, argued that programs should be decomposed into independent components with small and simple interfaces. They called their approach structured programming, \elide.}
{L’article d’Edsger Dijkstra de 1968 sur le système d’exploitation T.H.E. a introduit le concept de couches [Dijkstra 68]. Les premiers travaux de David Parnas ont posé de nombreux fondements conceptuels, notamment l’encapsulation de l’information [Parnas 72], les familles de programmes [Parnas 76], les structures inhérentes aux systèmes logiciels [Parnas 74] et l’utilisation de ces structures pour construire des sous-ensembles et des sur-ensembles de systèmes [Parnas 79]. \elide En 1972, Dijkstra et Hoare, avec Ole-Johan Dahl, ont soutenu que les programmes devaient être décomposés en composants indépendants dotés d’interfaces simples et réduites. Ils ont nommé leur approche « programmation structurée » \elide.}
{\ccite{bass_software_2021}[p.~21]}
\vspace{-1\baselineskip}
}{}

\enonce{AD13}{Assertion}{Il est possible d'organiser les documents comme les modules logiciels.}{}{P034}
\explication{AD13}{de l'assertion}{Il est possible d'organiser la documentation comme les modules logiciels}{
En 1981, S. D. Hester et David L. Parnas ont publié un article intitulé  \textit{Using Documentation as a Software Design Medium}  \cite{hester_using_1981} qui résume cette méthode. Selon l'article, voici les avantages qu'on retire de son utilisation~: 
\quoteenfrlist
{
    \item [i)] Ease of change. System functions that are likely to change are identified and information hiding is applied to minimize the amount of software affected by a change in these functions.
    \item [ii)] Control of the information about the functions of the system. A carefully structured requirements document is to be maintained throughout the life of the project.
    \item [iii)] Ordering of the development steps to meet the project objectives. Documentation of the useful subsets of the system and the dependencies between the software modules serve to guide the scheduling of the development effort.
    \item [iv)] Making the agreements between developers explicit. Misunderstandings are avoided and a smoother system integration is achieved by documenting the interfaces between the software modules of individual developers.
}
{
    \item [i)] Facilité de modification. Les fonctions système susceptibles d'évoluer sont identifiées et l'encapsulation des informations est mise en œuvre afin de minimiser l'impact d'une modification de ces fonctions sur le logiciel.
    \item [ii)]Maîtrise des informations relatives aux fonctions du système. Un document de spécifications structuré avec soin est maintenu à jour tout au long du projet.
    \item [iii)]Planification des étapes de développement en fonction des objectifs du projet. La documentation des sous-ensembles utiles du système et des dépendances entre les modules logiciels permet d'orienter la planification de développement.
    \item [iv)]Clarification des accords entre développeurs. La documentation des interfaces entre les modules logiciels des différents développeurs permet d'éviter les malentendus et d'assurer une meilleure intégration du système.
}
{\cite{hester_using_1981}}

En 1985, David L. Parnas, Paul Clements et David M. Weiss ont fourni des raisons expliquant l'inutilisation de la méthode et donnent finalement un exemple réaliste de son application \cite{parnas_modular_1985} en s'appuyant sur l'Operational Flight Program de l'avion A-7E. 
}{}

%How far is the modularization of open-source projects from an optimal modularization in terms of cohesion and coupling? \newline  [Answer :]
%Do open-source developers seek a fair compromise between cohesion and coupling when defining their modularization? \newline  [Answer:]
\enonce{AD14}{Assertion}{L'adaptabilité de plusieurs composants logiciels peut encore être amélioré.}{}{P034}
\explication{AD14}{l'assertion}{L'adaptabilité de plusieurs composants logiciels peut encore être amélioré.}{
Un article présente une étude sur le couplage et la cohérence du code source \cite{candela_using_2016}. Les auteurs analysent cent projets code source libre et leur appliquent des mesures de couplage et de cohérence, à l'aide de neuf  formules distinctes. Voici leurs conclusions qui confirment que la cohérence et le couplage pourraient être améliorés et donc l'adaptabilité~: 
\quoteenfrlist{
\item The achieved results highlight that the modularization of open-source systems is generally far from an optimal modularization in terms of both structural and conceptual cohesion/coupling. With very few exceptions, the refactoring effort needed to reach an optimal modularization from the original design as assessed by the MoJoFM is high.
\item Overall, open-source systems tend to give higher priority to cohesion over coupling. Our analysis shows that this is mainly due to the clear separation between the application layers pursued by developers when modularizing their system. Also, the results show that package cohesion and coupling might not be the main factors considered by developers when partitioning the classes of a system.
}{
\item Les résultats obtenus montrent que la modularisation des systèmes code source libre est généralement loin d'être optimale, tant du point de vue de la cohérence/du couplage structurel que conceptuel. À quelques rares exceptions près, l'effort de refactorisation nécessaire pour atteindre une modularisation optimale à partir de la conception originale, telle qu'évaluée par MoJoFM, est important.
\item Globalement, les systèmes en code source libre tendent à privilégier la cohérence par rapport au couplage. Notre analyse montre que cela est principalement dû à la séparation nette entre les couches applicatives recherchée par les développeurs lors de la modularisation de leur système. De plus, les résultats indiquent que la cohérence et le couplage des packages ne sont peut-être pas les principaux facteurs pris en compte par les développeurs lors du partitionnement des classes d'un système.
}{\cite{candela_using_2016}}

}{}

\paragraphe{P035}{\fl{La programmation littéraire\footnote{Cette traduction de \it{literate} fait référence à la langue utilisée dans les essais (scientifiques), par opposition à la langue littéraire (romans, poésie) et aux langages formels (utilisés en mathématiques et en programmation).} est une méthode qui veut aider à la compréhension du lectorat.} Donald E. Knuth, l'auteur de la programmation littéraire, veut prioriser le lectorat à la machine \ccite{knuth_literate_1992}[p.~99]. Dans un même document, les explications (documentation) sont jumelées à des instructions (code source). Puis un instrument s'occupe de traiter ces informations. D'ailleurs, les instruments qui permettent d'exécuter de la programmation littéraire ont continué de s'adapter pour ajouter de nouvelles fonctionnalités comme \cite{pieterse_case_2004}~:
\begin{itemize}
    \item utiliser d'autres langages dans les documents en entrées;
    \item fournir d'autres formats de sorties ;
    \item utiliser des hyperliens pour faire des références croisées ;
    \item fournir des intégrations avec d'autres instruments.
\end{itemize}
Toutefois, maintenir la cohérence lors de modifications demeure problématique, conclut l'auteur d'Elucidator un instrument de programmation littéraire \cite{normark_elucidative_2000}. Le \mbox{Theme-based} literate programming propose l'idée de faire correspondre les informations à un modèle et d'utiliser différentes techniques pour maintenir le modèle cohérent \cite{kacofegitis_theme-based_2002}. Cette approche propose l'utilisation de ces trois langages~:
\begin{itemize}
    \item langage de balisage extensible (XML, \textit{Extensible Markup Language} en anglais) pour représenter le document ;
    \item langage de transformation XML (XSLT, \textit{eXtensible Stylesheet Language Transformations} en anglais) pour spécifier les transformations ;
    \item langage pour la définition de type de document (DTD, \textit{Document Type Definitions} en anglais) pour spécifier le modèle.
\end{itemize}
Depuis plusieurs années, un instrument du nom de Jupyter notebook gagne en popularité. Cet instrument permet de créer une documentation à partir de texte et de code source exécutable. Plusieurs l'associent à de la programmation littéraire \cite{pimentel_understanding_2021}. Pourtant, l'objectif n'est pas de créer une entité informatique, mais de partager des idées \cite{granger_jupyter_2021}. Par conséquent, les Jypiter notebooks ont une faible reproductibilité, au plus 15~\%, pour un échantillon $(n=1~159~568)$ provenant de GitHub \cite{pimentel_understanding_2021} et encourageraient de mauvaises pratiques de programmation \cite{perkel_jupyter_2018}.
}{AD15, AD16, AD17, AD18}{AD15}

\enonce{AD15}{Citation}{La programmation littéraire cherche à expliquer en premier aux humains et en second à la machine.}{}{P035}
\explication{AD15}{de la citation}{La programmation littéraire cherche à expliquer en premier aux humains et en second à la machine.}{
En 1984, Donald E. Knuth publie un article intitulé \titleen{Literate Programming}{} \footnote{Il a développé TeX sur lequel repose LaTeX. LaTeX est le langage utilisé pour écrire ce mémoire.}, il y indique que~: 
\quoteenfr
{Instead of imagining that our main task is to instruct a computer what to do, let us concentrate rather on explaining to human beings what we want a computer do.}
{Au lieu d’imaginer que notre tâche principale consiste à donner des instructions à un ordinateur sur ce qu’il doit faire, concentrons-nous plutôt sur le fait d’expliquer aux êtres humains ce que nous voulons qu’un ordinateur fasse.}
{\ccite{knuth_literate_1992}[p.~99]}
\vspace{-1\baselineskip}
}{}

\enonce{AD16}{Assertion}{Les instruments qui exécutent la programmation littéraire ont continué de s'adapter.}{}{P035}
\explication{AD16}{l'assertion}{Les instruments qui exécutent la programmation littéraire ont continué de s'adapter.}{
En 2004, un article résume six différents environnements pour la programmation littéraire \cite{pieterse_case_2004}~:
\begin{itemize}
    \item \textit{the structure and dataflow of WEB (the first LPE)};
    \item \textit{the structure and dataflow of LIPED (a language independent LPE)};
\end{itemize}
\figmediumens
{}
{figure/chap1/documentation/lp.png}
{(modifiée de \cite{pieterse_case_2004})}
{}{}{H}{Représentations d'environnements de programmation littéraire}
\begin{itemize}
    \item \textit{the structure and dataflow of AOPS (an OOP LPE that includes a hypertext browser)};
    \item \textit{the structure and dataflow of LPW (a LPE that integrates a third party editor)};
    \item \textit{the structure and dataflow of warp (a LPE that integrates third party programming environment)};
    \item \textit{the structure and dataflow of Elucidator (a tool to enhance a programming environment for elucidative programming)}.
\end{itemize}
La figure~B.15 présente les diagrammes de quatre de ces environnements, dont voici la description textuelle~: 
\begin{itemize}
    \item Le diagramme a) emploie deux processus distincts pour extraire le code exécutable ou le document typographique.
    \item Le diagramme b) montre qu'il est possible de spécifier le langage de programmation et le langage de formatage du document. Cela correspond aux \textit{Language Specifications} et \textit{Typographic Specifications}.
    \item Le diagramme c) représente un environnement où il est possible de naviguer entre les relations en utilisant un fureteur. 
    \item Le diagramme d) montre que l'instrument Elucidator s'occupe de maintenir les relations entre les concepts provenant de deux endroits.
\end{itemize}

}{}

\enonce{AD17}{Assertion}{Maintenir la cohérence lors de modifications avec la méthode de programmation littéraire est problématique.}{}{P035}
\explication{AD17}{de l'assertion}{Maintenir la cohérence lors de modifications avec la méthode de programmation littéraire est problématique.}{
L'auteur d'Elucidator indique à quel point cela vient difficile de maintenir la cohérence lors de modifications~:
\quoteenfr
{On the negative side, it is less trivial to ensure a proper updating of the documentation upon program modifications, and vice versa. The problem is amplified by the fact that a single program fragment may be addressed at multiple places in the documentation.}
{En revanche, il est plus complexe de garantir la mise à jour de la documentation suite à des modifications du programme, et inversement. Ce problème est accentué par le fait qu'un même fragment de programme peut être mentionné à plusieurs endroits dans la documentation.}
{\cite{normark_elucidative_2000}}

Le \titleen{Theme-based literate programming}{} tente de maintenir la cohérence en utilisant modèles et techniques \cite{kacofegitis_theme-based_2002}. Ces techniques comprennent du XML, du XSLT et du DTD et, comme décrits, ils servent à représenter et traiter formellement les informations contenues dans les documents~:
\quoteenfrlist {
    \item[][---] We use XML [12] as the fundamental representation of LP documents \elide.
    \item[][---] XML parsing tools are used in our authoring tool and XSLT [6] transformations are used to perform much of the processing \elide.
    \item[][---] Document Type Definitions (DTD) allow a grammar for individual chunks and their relationships to be specified.
} {
\item[][---] Nous utilisons XML [12] comme représentation fondamentale des documents LP \elide.
\item[][---] Notre outil de création utilise des outils d'analyse XML et des transformations XSLT [6] pour effectuer une grande partie du traitement \elide.
\item[][---] Les définitions de type de document (DTD) permettent de spécifier une grammaire pour chaque segment de données et leurs relations.
}{\cite{kacofegitis_theme-based_2002}}
\vspace{-1\baselineskip}
}{}

\enonce{AD18}{Assertion}{Plusieurs individus croient à tort que l'instrument Jupyter est de la programmation littéraire.}{}{P035}
\explication{AD18}{de l'assertion}{Plusieurs individus croient à tort que l'instrument Jupyter est de la programmation littéraire.}{
Voici la description provenant de la documentation officielle de Jupyter qui combine deux composantes, elle provient de la documentation officielle~:
\quoteenfrlist{
    \item A web application: A browser-based editing program for interactive authoring of computational notebooks which provides a fast interactive environment for prototyping and explaining code, exploring and visualizing data, and sharing ideas with others
    \item Computational Notebook documents: A shareable document that combines computer code, plain language descriptions, data, rich visualizations like 3D models, charts, mathematics, graphs and figures, and interactive controls.
}{
\item Une application web~: un programme d’édition basé sur un navigateur permettant la création interactive de carnets de calcul. Il offre un environnement interactif rapide pour le prototypage et l’explication de code, l’exploration et la visualisation de données, ainsi que le partage d’idées.
\item Documents de carnets de notes pour le calcul~: un document partageable qui combine du code informatique, des descriptions en langage clair, des données, des visualisations riches (modèles 3D, diagrammes, formules mathématiques, graphiques et figures) et des commandes interactives.
}{\cite{jupyter_doc_2026}}

Un article décrit l'analyse de carnets Jupyter (1~159~568) provenant de répertoires GitHub (264~020). L'article utilise 24 fois le terme \emphase{literate programming} et décrit la popularité de Jupyter et les problèmes de reproductibilité~:
\quoteenfrlist {
    \item[][---] Jupyter Notebook is the most widely-used system for interactive literate programming (Shen 2014). 
    \item[][---] \elide The system was released in 2013, and today there are over 9 million notebooks in GitHub (Parente 2020).
    \item[][---] \elide most notebooks do not test their code and that a large number of notebooks has characteristics that hinder the reasoning and the reproducibility, such as out-of-order cells, non-executed code cells, and the possibility of hidden states.
    \item[][---] \elide We achieved a reproducibility rate that ranged from 4.90\% to 15.04\%.
} {
\item[][---] Jupyter Notebook est le système le plus utilisé pour la programmation littéraire interactive (Shen 2014).
\item[][---] \elide Le système a été lancé en 2013 et compte aujourd'hui plus de 9 millions de cahiers sur GitHub (Parente 2020).
\item[][---] \elide La plupart des carnets de notes ne testent pas leur code et un grand nombre d'entre eux présentent des caractéristiques qui entravent le raisonnement et la reproductibilité, tels que des cellules mal ordonnées, des cellules de code non exécutées et la possibilité d'états cachés.
\item[][---] \elide Nous avons obtenu un taux de reproductibilité compris entre 4,90~\% et 15,04~\%.
}{\cite{pimentel_understanding_2021}}

Un article cite des problèmes dans l'utilisation de code avec Jupyter, la modularité y semble particulièrement problématique~:
\quoteenfr{
Joel Grus \elide Jupyter notebooks also encourage poor coding practice, he says, by making it difficult to organize code logically, break it into reusable modules and develop tests to ensure the code is working properly. \elide Those aren’t insurmountable issues, Grus concedes, but notebooks do require discipline when it comes to executing code \elide.
}{Selon Joel Grus, les carnets Jupyter encouragent également les mauvaises pratiques de programmation, car ils rendent difficiles l'organisation logique du code, son découpage en modules réutilisables et le développement de tests pour garantir son bon fonctionnement. Grus concède que ces problèmes ne sont pas insurmontables, mais que l'utilisation des carnets exige de la rigueur lors de l'exécution du code.}
{\cite{perkel_jupyter_2018}}

Un autre article dont les auteurs sont Brian E. Grange et Fernando Pérez \footnote{Fernando Pérez est cofondateur du projet Jupyter.} explique qu'il s'agit de \emphase{computational narrative} et non de literate programming. Il s'en distingue par le fait que l'objectif est de faire comprendre au lectorat un modèle et non de permettre à une machine de l'exécuter~:
\quoteenfr
{These human uses of computational narratives illustrate how and why the ways that Jupyter notebooks are used are so dramatically different from traditional software engineering tools where the goal is for one group of people to write software that is subsequently deployed to and used by an entirely different group of users. While Knuth’s literate programming paradigm weaves human-oriented documentation into the software engineering process, a computational narrative is distinct in its incorporation of interactive computing as the central element. The outcome here is not a software product but ideas and understanding that are “deployed” to other humans.}
{Ces exemples d'utilisation humaine des récits informatiques illustrent en quoi les utilisations des carnets Jupyter diffèrent radicalement des outils de génie logiciel traditionnels, dont l'objectif est qu'un groupe de personnes développe un logiciel destiné ensuite à être déployé et utilisé par un tout autre groupe d'utilisateurs. Alors que le paradigme de programmation littéraire de Knuth intègre une documentation orientée utilisateur au processus de génie logiciel, le récit informatique se distingue par son approche centrée sur l'informatique interactive. Le résultat n'est pas ici un produit logiciel, mais des idées et une compréhension « déployées » auprès d'autres personnes.}
{\cite{granger_jupyter_2021}}
\vspace{-1\baselineskip}
}{}

\subsubsection{Des techniques informatiques (donc, formelles)}

\paragraphe{P036}{Dans les techniques informatiques, il y a les langages formels. L'utilisation de langages formels permet de traiter formellement les documents avec des machines, comme ce fût le cas avec la programmation littéraire. À cela, s'ajoutent quelques exemples d'utilisation de modèles relationnels pour maintenir la cohérence entre les documents et à l'intérieur de ceux-ci.}{}{}

\paragraphe{P037}{Avant d'aborder les techniques utilisées pour documenter en informatique appliquée, il est nécessaire de définir ces concepts~:  
\begin{itemize}
    \item langage~: système pour communiquer comprenant des symboles et des règles;
    \item langage formel~: est composé d'ensembles explicites de symboles et d'un ensemble explicite de règles appelé grammaire formelle;
    \item langage naturel~: est un langage non formel;
    \item architecture informatique~: organisation d'entités informatiques résultant de la phase de conception;
    \item langage de description d'architecture (ADL, \textit{architecture description language} en anglais)~: langage avec une syntaxe et une sémantique spécifiées permettant de décrire une architecture;
    \item ontologie informatique~: «~corpus structuré de concepts, qui est modélisé dans un langage permettant l'exploitation par un ordinateur des relations sémantiques ou taxonomiques établies entre ces concepts.~» \footnote{Une \emphase{ontologie informatisée} a pour source (origine) une \emphase{ontologie générale} à laquelle il est présumé qu’elle tente d’être le plus fidèle possible.} \cite{Gdt_OntInfo_1}.
\end{itemize}
\vspace{-1\baselineskip}
}{AD19, AD20, AD21, AD22, AD23, AD24}{AD19}

\enonce{AD19}{Définition}{Langage}{
Système pour communiquer comprenant des symboles et des règles.
}{P037}
\explication{AD19}{de la définition}{Langage}{
La définition formulée est~: système pour communiquer comprenant des symboles et des règles.

La littérature fournit ces informations qui ont servi à construire cette définition~:
\begin{itemize}
  \item définition 1 de langage~: «~Faculté propre à l'être humain d'exprimer sa pensée ou de communiquer au moyen d'une langue et de comprendre les messages d'autrui.~» \cite{Gdt_langage_0} ;
  \item note sur la définition 1 de langage~: «~Le langage comporte différentes composantes (phonologique, lexicale, morphosyntaxique, pragmatique, etc.) et se développe dans les volets réceptif (compréhension) et expressif (production).» \cite{Gdt_langage_0} ;
  \item définition 2 de langage~: «~Tout système de signes permettant la communication.~» \cite{Robert_langage_0} ;
  \item définition de langue~: «~Système d'expression et de communication par des moyens phonétiques et éventuellement graphiques, commun à un groupe social.~» \cite{Robert_langue_0}.
\end{itemize}

}{}

\enonce{AD20}{Définition}{Langage formel}{
Est composé d'ensembles explicites de symboles et d'un ensemble explicite de règles appelé grammaire formelle.
}{P037}
\explication{AD20}{de la définition}{Langage formel}{
La définition formulée est~: est composé d'ensembles explicites de symboles et d'un ensemble explicite de règles appelé grammaire formelle.

La littérature fournit ces informations qui ont servi à construire cette définition~:
\begin{itemize}
  \item définition 1 de langage formel~: «~Langage qui utilise un ensemble de termes et de règles syntaxiques pour permettre de communiquer sans aucune ambiguïté.~»
  \cite {Gdt_lf_0};
  \item définition 2 de langage formel~: 
  \enfr{A formal language is a well-defined set of words from a given alphabet. Usually it is defined by a grammar, which is a set of rules determining the membership of the language.}{Un langage formel est un ensemble bien défini de mots appartenant à un alphabet donné. Il est généralement défini par une grammaire, c'est-à-dire un ensemble de règles déterminant l'appartenance des mots au langage.}
  \ccite{roggenbach_formal_2021}[p.~6] ;
  \item définition 3 de langage formel~: 
  \enfr{1) a set of symbols (the alphabet of his language) and (2) a set of formation rules determining which sequences of symbols from his alphabet are wffs of his language.}{1) un ensemble de symboles (l'alphabet de sa langue) et (2) un ensemble de règles de formation déterminant quelles séquences de symboles de son alphabet sont des formules bien formées de sa langue.}
  \ccite{hunter_metalogic_1973}[p.~4];
  \item description de langage formel~:
  \enfr{A formal grammar generates a formal language, which is set of finite length sequences of symbols created by applying the production rules of the grammar.}
  {Une grammaire formelle génère un langage formel, qui est un ensemble de séquences de symboles de longueur finie créées en appliquant les règles de production de la grammaire.}
  \ccite{oregan_mathematics_2020}[p.~188]
  \item définition 4 de langage formel~:
  \enfr{\elide language whose rules are explicitly established prior to its use \elide.}
  {\elide un langage dont les règles sont explicitement établies avant son utilisation \elide.}
 \cite{iso_24765_2017}.
\end{itemize}

}{}

\enonce{AD21}{Définition}{Langage naturel}{
Est un langage non formel.
}{P037}
\explication{AD21}{de la définition}{Langage naturel}{
La définition formulée est~: est un langage non formel.

La littérature fournit ces informations qui ont servi à construire cette définition~:
\begin{itemize}
    \item définition 1 de langage naturel~: \enfr{\elide language whose rules are based on current usage without being specifically prescribed \elide.}{\elide langue dont les règles sont basées sur l'usage courant sans être spécifiquement prescrites \elide.} \cite{iso_24765_2017};
    \item définition 2 de langage naturel~: «~Langage humain par opposition aux langages de programmation.~» \cite{Gdt_ln_0}.
\end{itemize}
Dans le mémoire, il est souvent préféré le terme \emphase{langage non formel} à \emphase{langage naturel}, mais, puisque l'usage de langage naturel est répandu, il devait être défini.
}{}

\enonce{AD22}{Définition}{Architecture informatique}{
Organisation d'entités informatiques résultant de la phase de conception.
}{P037}
\explication{AD22}{de la définition}{Architecture informatique}{
La définition formulée est~: organisation d'entités informatiques résultant de la phase de conception.

La littérature fournit ces informations qui ont servi à construire cette définition~:
\begin{itemize}
    \item définition d'architecture~: \enfr{\elide fundamental concepts or properties of an entity in its environment (3.13) and governing principles for the realization and evolution of this entity and its related life cycle processes \elide.}{\elide concepts ou propriétés fondamentaux d'une entité dans son environnement (3.13) et principes directeurs de la réalisation et de l'évolution de cette entité et de ses processus de cycle de vie associés \elide.}
    \ccite{iso_42010_2022}[p.~2];
    \item définition d'architecture d'un système~: \enfr{The software architecture of a system is the set of structures needed to reason about the system. These structures comprise software elements, relations among them, and properties of both.}{L'architecture logicielle d'un système est l'ensemble des structures nécessaires pour raisonner sur ce système. Ces structures comprennent des éléments logiciels, les relations entre eux et leurs propriétés.}
    \ccite{bass_software_2021}[p.~1];
    \item description d'\textit{Architecture versus Design}~: \enfr{Architecture is design, but not all design is architecture. That is, many design decisions are left unbound by the architecture—it is, after all, an abstraction—and depend on the discretion and good judgment of downstream designers and even implementers.}{L'architecture est une conception, mais toute conception n'est pas de l'architecture. Autrement dit, de nombreuses décisions de conception ne sont pas liées à l'architecture — qui est, après tout, une abstraction — et dépendent du discernement et du bon jugement des concepteurs et même des implémenteurs en aval.}
    \ccite{bass_software_2021}[p.~3];
    \item description de \textit{software system architecture}~: \enfr{[Cite Barry W. Boehm] A collection of software and system components, connections, and constraints. A collection of system stakeholders' need statements. A rationale which demonstrates that the components, connections, and constraints define a system that, if implemented, would satisfy the collection of system stakeholders need statements.}
    {[Cite Barry W. Boehm] Un ensemble de composants logiciels et système, de connexions et de contraintes. Un ensemble d'énoncés de besoins des parties prenantes du système. Une justification démontrant que les composants, les connexions et les contraintes définissent un système qui, s'il était implémenté, satisferait l'ensemble des énoncés de besoins des parties prenantes du système.}
    \ccite{printz_architecture_2012}[p.~61];
    \item description d'architecture~:\enfr
    {There is little meaningful difference between «~design~» and «~architecture~». For clarity, this discussion will use «~design~» for the process and «~architecture~» for its result (then we do not need «~architect~» as a verb). The same convention can be applied to «~implementation~» and «~code~».}
    {Il y a peu de différence significative entre «~conception~» et «~architecture~». Par souci de clarté, nous utiliserons ici «~conception~» pour désigner le processus et «~architecture~» pour son résultat [le verbe «~architecturer~» devient alors superflu]. La même convention s’applique à «~implémentation~» et «~code~».}
    \ccite{meyer_agile_2014}[p.~37];
\end{itemize}

}{}

\enonce{AD23}{Définition}{Langage de description d'architecture (ADL, \textit{Architecture Description Language} en anglais)}{
Langage avec une syntaxe et une sémantique spécifiées qui décrit l'architecture d'entité informatique.
}{P037}
\explication{AD23}{de la définition}{Langage de description d'architecture (ADL, \textit{Architecture Description Language} en anglais))}{
La norme ISO 42010 intitulée \titlefr{Logiciel, systèmes et entreprise — description de l'architecture} définit les langages de description d'architecture (ADL, \textit{Architecture Description Languages} en anglais)~:
\quoteenfr
{An ADL is a specified syntax and semantics intended for use in describing the architecture of an entity of interest. An ADL is a language for stakeholders, including those involved in the architecting effort that allows the expression of architecture considerations by means of AD elements pertaining to the entity of interest, and the architecting context. An AD can use more than one ADL, even a different ADL for each viewpoint, or even distinct ADLs for each model kind specified by a single architecture viewpoint.}
{Un langage de description architecturale (ADL) est une syntaxe et une sémantique spécifiées, destinées à décrire l'architecture d'une entité donnée. L'ADL est un langage utilisé par les parties prenantes, notamment celles impliquées dans la conception architecturale, pour exprimer les considérations architecturales au moyen d'éléments de description architecturale relatifs à l'entité et au contexte de conception. Une description architecturale peut utiliser plusieurs ADL, voire une ADL différente pour chaque point de vue, ou encore des ADL distinctes pour chaque type de modèle spécifié par un même point de vue architectural.}
{\ccite{iso_42010_2022}[p.~2]}

La norme comprend un modèle conceptuel qui peut tenir compte des différentes déclinaisons d'ADL, des exemples se retrouvent entre parenthèses~:
\quoteenfrlist{
\item [] [---] a single model kind (e.g. Darwin[37]); 
\item [] [---] multiple model kinds (e.g. UML, Rapide, AADL); 
\item [] [---] single viewpoint (e.g. ACME[33]); 
\item [] [---] multiple viewpoints (e.g. ArchiMate, SysML).
}{
\item [] [---] un seul type de modèle (par exemple, Darwin[37]);
\item [] [---] plusieurs types de modèles (par exemple, UML, Rapide, AADL);
\item [] [---] un seul point de vue (par exemple, ACME[33]);
\item [] [---] plusieurs points de vue (par exemple, ArchiMate, SysML).
}{\ccite{iso_42010_2022}[p.~39]}
La norme ne prescrit pas l'usage de la modélisation formelle~:
\quoteenfr
{This document does not prescribe the extent or expectation regarding the use of formal modelling methods in an AD. While informal methods can be used effectively, formal modelling methods are often less ambiguous.}
{Ce document ne prescrit pas l'étendue ni les attentes concernant l'utilisation des méthodes de modélisation formelles dans une description d'architecture. Si les méthodes informelles peuvent être utilisées efficacement, les méthodes de modélisation formelles sont souvent moins ambiguës.}
{\ccite{iso_42010_2022}[p.~20]}

}{}

\enoncewithoutexp{AD24}{Définition}{Ontologie informatique}
{«~Corpus structuré de concepts, qui est modélisé dans un langage permettant l'exploitation par un ordinateur des relations sémantiques ou taxonomiques établies entre ces concepts.~» \cite{Gdt_OntInfo_1}}{P037}

\paragraphe{P038}{\fl{Les langages formels permettent à des instruments d'assurer entre autres l'adaptabilité des documents.} D'abord, le XML est un exemple de langage formel qui permet d'exprimer des informations selon une structure. Il est une déclinaison du Standard Generalized Markup Language dont l'objectif est de permettre le partage de document \cite{kilpelainen_sgml_2001}. 
\sep
Ensuite, il y a les langages de schémas qui servent à vérifier les XML. Une norme sur les langages de schémas en comprend six, dont la définition des types de document (DTD, \textit{Document Type Definition} en anglais), la définition du schéma XML (XSD, \textit{XML Schema Definition} en anglais) et le Schematron \cite{ISO_19757p11_2011}. 
\sep
Néanmoins, le XML et les langages de schéma ne permettent pas d'exprimer pleinement la sémantique des données, et ce, même pour un langage de schéma aussi poussé que le Schematron \cite{schematron_semantic}. Cela entraîne des difficultés pour réaliser des transformations adéquates vers le format Web Ontology Language (OWL) qui sert aux ontologies informatiques \cite{aussenac_ecise_2022}. 
\sep
Pour finir, le XML, le DTD ou le XSD sont utilisés par deux instruments de documentation et une architecture de documentation. Ce sont DocBook, S1000D et l'architecture des types d’informations Darwin (DITA, \textit{Darwin Information Type Architecture} en anglais)  \ccite{farazi_model-based_2017}[][self_dita_2011][p.~226]. Dans les fonctionnalités offertes par le DITA, il y a au moins la modularité et la réutilisation de contenu qui aident à l'adaptabilité \ccite{self_dita_2011}[p.~222]. 
}{AD25, AD26, AD27, AD28, AD29}{AD25}

\enonce{AD25}{Assertion}{Le XML est un langage formel qui permet d'exprimer des informations selon une structure.}{}{P038}
\explication{AD25}{de l'assertion}{Le XML est un langage formel qui permet d'exprimer des informations selon une structure.}{
%L'origine du XML vise à transmettre des informations comme cela est décrit dans cet article~: 
L'origine du XML vise à transmettre des informations comme cela est décrit~:
\quoteenfr
{The Standard Generalized Markup Language (SGML) [12, 13] promotes the interchangeability and application-independent management of electronic documents by providing a syntactic metalanguage for the definition of textual markup systems. The Extensible Markup Language (XML) [2] is, essentially, a simplified and more restrictive version of SGML. The goal of XML is to allow SGML documents to be served, received, and processed on the Web. It is the proposed syntactic metalanguage for the specification of document grammars for W3 documents.}
{Le langage SGML (Standard Generalized Markup Language) [12, 13] favorise l'interchangeabilité et la gestion indépendante des applications des documents électroniques en fournissant un métalangage syntaxique pour la définition des systèmes de balisage textuel. Le langage XML (Extensible Markup Language) [2] est, en substance, une version simplifiée et plus restrictive du SGML. L'objectif du XML est de permettre la diffusion, la réception et le traitement des documents SGML sur le Web. Il s'agit du métalangage syntaxique proposé pour la spécification des grammaires de documents W3C.}
{\cite{kilpelainen_sgml_2001}}

Le XML a permis de structurer les informations avec des métalangages, un article explique son utilisation dans le Darwin Information Type Architecture (DITA)~:
\quoteenfr
{SGML and XML are recognized as meta languages that allow communities of data owners to describe their information assets in ways that reflect how they develop, store, and process that information.}
{SGML et XML sont reconnus comme des métalangages qui permettent aux communautés de propriétaires de données de décrire leurs actifs informationnels de manière à refléter la façon dont elles développent, stockent et traitent ces informations.}
{\cite{day_introduction_2005}}

L'utilisation du XML s'est répandue, comme l'explique le guide du DITA~:
\quoteenfr
{The development of XML in the late 1990s transformed the way in which knowledge was stored. XML permitted structured information standards to be created for the storage of knowledge and data for all  types of industries.}
{Le développement du XML à la fin des années 1990 a transformé la manière dont les connaissances étaient stockées. Le XML a permis la création de normes d'information structurées pour le stockage des connaissances et des données dans tous les secteurs d'activité.}
{\ccite{self_dita_2011}[p.~225]}
\vspace{-1\baselineskip}
}{}

%\figannexe{
%\tabmen
%{\textit{Use of a combination of schematypens}}
%{figure/chap1/documentation/xsd.png}
%{(source~: \ccite{ISO_19757p11_2011}[p.~5])}
%{}{}{tb}{Langages de schéma référés par le standard ISO~19757}
%}

\enonce{AD26}{Assertion}{Les langages de schéma XML (DTD, XSD, Schematron) permettent d'exprimer des vérifications sur le XML.}{}{P038}
\explication{AD26}{de l'assertion}{Les langages de schéma XML (DTD, XSD, Schematron) permettent d'exprimer des vérifications sur le XML.}{
La norme \textit{Information technology — Document Schema Definition Languages (DSDL)} soutient six langages de schéma XML~: DTD, W3C XML Schema, RELAX NG, RELAX NG - compact syntax, Schematron et NVDL.

Un article explique les grammaires et les exceptions qui se trouvent dans le SGML et le XML~:
\quoteenfr
{Both SGML and XML allow users to define document-type definitions (DTDs), which are essentially extended context-free grammars expressed in a notation that is similar to extended Backus–Naur form. The right-hand side of a production, called a content model, is both an extended and a restricted regular expression.}
{SGML et XML permettent tous deux de définir des définitions de type de document (DTD), qui sont essentiellement des grammaires hors contexte étendues, exprimées dans une notation similaire à la forme Backus-Naur étendue. La partie droite d'une production, appelée modèle de contenu, est à la fois une expression régulière étendue et une expression régulière restreinte.}
{\cite{kilpelainen_sgml_2001}}
Ces différents langages répondent à des besoins différents de vérifications. L'auteur de l'ouvrage sur le Schematron indique que ces vérifications peuvent arriver lors des différentes étapes de traitement~:
\quoteenfrlist {
    \item[][---] For XML, checking the ground rules of the language family is an absolute precondition for any subsequent checks. An XML document that passes this stage is called well-formed.
    \item[][---] To check grammar you can choose from several XML validation languages. The most common ones are DTD, W3C XML Schema, and RELAX NG. If this stage is passed, an XML document is called valid.
    \item[][---] Schematron is a validation language that checks business rules. A set of checks written in the Schematron language is called a Schematron schema. There’s no official name for a document that passes this stage, but let’s call it Schematron valid.
} {
    \item[][---] Pour XML, la vérification des règles fondamentales de la famille de langages est une condition préalable absolue à toute vérification ultérieure. Un document XML qui réussit cette étape est dit bien formé.
    \item[][---] Pour vérifier la grammaire, plusieurs langages de validation XML sont disponibles. Les plus courants sont DTD, W3C XML Schema et RELAX NG. Si cette étape est réussie, un document XML est dit valide. 
    \item[][---] Schematron est un langage de validation qui vérifie les règles métier. Un ensemble de vérifications écrites en Schematron est appelé un schéma Schematron. Il n'existe pas de nom officiel pour un document qui réussit cette étape, mais nous l'appellerons « schéma valide ».
}{\ccite{siegel_schematron_2022}[p.~7]}
Ces étapes rappellent ceux d'un compilateur où différents analyseurs vont vérifier et traiter des informations spécifiques.
}{}

\enonce{AD27}{Assertion}{Les langages de schéma XML ne permettent pas d'exprimer totalement la sémantique.}{}{P038}
\explication{AD27}{de l'assertion}{Les langages de schéma XML ne permettent pas d'exprimer totalement la sémantique.}{
Rick Jelliffe créateur du Schematron commente sur son blogue les propos de l'auteur d'un autre langage de schéma XML (RELAX NG) Makoto Murata sur le fait que ces langages ne concernent pas la sémantique. Selon lui, cela demeure plus nuancé~:
\quoteenfr
{So I think Murata-san has a good rule-of-thumb there, that schemas don’t imply any semantics. Certainly schemas never imply all semantics, and often not even all the  structural pr data type or linking rules.  But that is not to say that schemas don’t reflect semantic considerations through and through. The implying (ultimately, by humans) is the problem: the idea that any schema completely expresses everything about a document type.}
{Je pense donc que Murata-san a raison sur ce point : les schémas n'impliquent aucune sémantique. Certes, ils n'impliquent jamais toute la sémantique, et souvent, même pas toutes les règles structurelles, de type de données ou de liaison. Mais cela ne signifie pas pour autant que les schémas ne reflètent pas pleinement les considérations sémantiques. Le problème réside dans l'implication (finalement, par les humains) : l'idée qu'un schéma puisse exprimer complètement tout ce qui concerne un type de document.}
{\footnote{Rick Jelliffe, site web, «~Schemas do not imply any semantics of documents~», \url{https://schematron.com/opinion/_schemas_do_not_imply_any_semantics_of_documents_.html} (consulté le 2026-03-10).}}
Ces vérifications ont des limites sur la sémantique, comme indiqué dans un article de 2022 qui propose une solution semi-automatique pour passer de XSD à OWL, puisque le XSD ne gère pas la sémantique~:
\quotefr
{Bien que le passage de données XML ou de schémas XSD à une représentation sémantique (RDF, RDFS ou OWL) soit une question étudiée de longue date dans le domaine du Web sémantique, une transformation automatique simple est rarement correcte. Ce processus se heurte à la difficulté de gérer des nœuds anonymes, de traiter la représentation de nœuds complexes, de capturer la sémantique des balises purement structurelles, ou encore de produire des constructeurs liés à la structuration [12, 10, 4, 22].}
{\cite{aussenac_ecise_2022}}

}{}

\enonce{AD28}{Assertion}{Le XML, le DTD, le XSD sont utilisés par des instruments de documentation.}{}{P038}
\explication{AD28}{de l'assertion}{Le XML, le DTD, le XSD sont utilisés par des instruments de documentation.}{
En 2017, un article sur le \textit{Model-based documentation} décrit des instruments servant à spécifier la documentation~:
\quoteenfrlist {
    \item[][---] Darwin Information Type Architecture (DITA) is a documentation architecture, based on XML, created to provide an end-to-end support, from semi-automatic authoring all the way through to automatic delivery of technical documentation.
    \item[][---] DocBook is a documentation system developed to facilitate the production of computer hardware and software related publications such as books and articles. \elide DocBook DTDs and schema languages such as XML Schema are employed to provide structure to the resulting documentation.
    \item[][---] S1000D is a specification originally designed for aerospace and defense industry for the preparation and production of technical documentation. \elide At its core, there is the principle of information reuse that is facilitated by Data Modules. The Data Modules in version 4.0 of S1000D are represented in XML. S1000D is suitable for publishing in both printed and electronic form. The S1000D document generation architecture includes a data storage component, which is called Common Source Database (CSDB).
} {
    \item[][---] Darwin Information Type Architecture (DITA) est une architecture de documentation, basée sur XML, conçue pour fournir une prise en charge complète, de la création semi-automatique à la diffusion automatique de la documentation technique.
    \item[][---] DocBook est un système de documentation développé pour faciliter la production de publications relatives au matériel et aux logiciels, telles que des livres et des articles. \elide Les DTD DocBook et les langages de schémas, comme XML Schema, sont utilisés pour structurer la documentation résultante.
    \item[][---] S1000D est une spécification initialement conçue pour l'industrie aérospatiale et de défense, pour la préparation et la production de documentation technique. \elide Son principe fondamental repose sur la réutilisation de l'information, facilitée par les modules de données. Dans la version 4.0 de S1000D, ces modules sont représentés en XML. S1000D convient à la publication sous forme imprimée et électronique. L'architecture de génération de documents S1000D comprend un composant de stockage de données appelé base de données source commune (CSDB).
}{\cite{farazi_model-based_2017}}
Le guide \textit{The DITA style guide: best practices for authors} expose le contexte technologique du DITA dans la figure~B.16 l'illustre. Le DITA repose sur le XML et fournit les moyens d'exprimer des informations sur le logiciel, le matériel, etc.
}{}

\figannexe{
\figmediumen
{\textit{DITA in relation to other standards and technologies}}
{figure/chap1/documentation/dita.png}
{(source~: \textit{figure} 46 \cite{self_dita_2011})}
{}{}{H}{Technologies sur lesquelles le DITA repose}

}

\enonce{AD29}{Assertion}{Le DITA permet une adaptabilité pour les documents.}{}{P038}
\explication{AD29}{de l'assertion}{Le DITA permet une adaptabilité pour les documents.}{
Voici une liste des fonctions offertes par DITA \ccite{self_dita_2011}[p.~222]~: 
\enfr{modularity, structured authoring (reduces authoring time, increases analysis time), information typing, minimalism, inheritance, specialization, single-source, topic-based, metadata, conditional processing, component publishing DITA authoring concepts, task-orientation, content re-use, translation-friendly structure}{modularité, rédaction structurée (réduction du temps de rédaction, augmentation du temps d'analyse), typage de l'information, minimalisme, héritage, spécialisation, source unique, approche thématique, métadonnées, traitement conditionnel, publication par composants : concepts de rédaction DITA, orientation tâche, réutilisation du contenu, structure facilitant la traduction}
}{}

\paragraphe{P039}{\fl{L'utilisation de langages de descriptions d'architecture est une technique, mais plusieurs utilisateurs préfèrent les langages non formels.} Un article répertorie 177 différents langages pour architecture \cite{malavolta_what_2013}. Plusieurs utilisateurs des langages d'architecture voient un intérêt à bien définir la sémantique, il s'agit de la troisième fonctionnalité la plus utile, selon eux, après les multiples vues architecturales ou un langage exprimé graphiquement. 
\sep
Néanmoins, plusieurs articles affirment le désintérêt des utilisateurs dans l'utilisation de langages formels pouvant aller jusqu'à exprimer la sémantique \cite{malavolta_what_2013, ozkaya_are_2013, rost_software_2013}. Selon plusieurs sondages, l'Unified Modeling Language (UML) est plus souvent utilisé que les langages de descriptions d'architecture (ADL, Architecture Description Language en anglais) \cite{malavolta_what_2013, rost_software_2013}. 
}{AD30, AD31, AD32}{AD30}

\enonce{AD30}{Assertion}{Les langages d'architecture sont nombreux.}{}{P039}
\explication{AD30}{de l'assertion}{Les langages d'architecture sont nombreux.}{
En 2013, un article inventorie avec des recherches et un sondage 177 langages d'architecture~:
\quoteenfr
{\elide we discovered that another community uses the term ADL to refer to languages for specifying the architecture for retargetable compilation \elide.
The list of 57 ALs produced in step 1 is enriched by the 120 ALs found in step 2.7 This output is used in Phase 2 to build the community of practitioners target of our survey.}
{\elide nous avons découvert qu'une autre communauté utilise le terme ADL pour désigner les langages permettant de spécifier l'architecture pour la compilation adaptative \elide. La liste des 57 langages d'architecture (LA) produite à l'étape 1 est enrichie par les 120 LA trouvés à l'étape 2.7. Ces résultats sont utilisés dans la phase 2 pour constituer la communauté de praticiens ciblée par notre enquête.}
{\cite{malavolta_what_2013}}

}{}

\enonce{AD31}{Assertion}{La possibilité d'exprimer la sémantique est une fonctionnalité appréciée et recherchée par plusieurs professionnels informatiques utilisant des langages d'architecture.}{}{P039}
\explication{AD31}{l'assertion}{La possibilité d'exprimer la sémantique est une fonctionnalité appréciée et recherchée par plusieurs professionnels informatiques utilisant des langages d'architecture.}{
Un sondage comprenant 48~professionnels informatiques, utilisant des langages d'architecture, est réalisé \cite{malavolta_what_2013}. Les participants ont été invités à donner leur avis sur certaines fonctionnalités, ainsi que sur leurs utilités passées et leurs utilités potentielles. La figure~B.17 présente un résumé des résultats. Ces résultats mettent en évidence les observations suivantes~:
\begin{itemize}
\item les trois fonctionnalités les plus utiles sont~: support de plusieurs vues architecturales, syntaxe graphique, sémantique bien définie ; 
\item les trois fonctionnalités les plus utiles sont~: support d'outils, support de plusieurs vues architecturales, sémantique bien définie.
\end{itemize}
Après avoir exposé différents ADL, un autre article exprime trois problèmes qui restent non résolus \cite{ozkaya_are_2013}~:
\figmsen
{\textit{Usefulness of AL features in past projects and for future projects}}
{figure/chap1/documentation/adl.png}
{(source~: \textit{Figure}~9 \cite{malavolta_what_2013})}
{}{}{H}{Fonctionnalités utiles pour les langages d'architecture selon des professionnels}
\quoteenfrlist {
    \item[][---] \elide lack of support for formal analysis of architectures \elide.
    \item[][---] \elide notations that sometimes make specifying large and complex system architectures harder than it should be \elide.
    \item[][---] \elide potential un-realizability of system architectures \elide.
} {
    \item[][---] \elide le manque de support pour l'analyse formelle des architectures \elide.
    \item[][---] \elide les notations qui ne sont pas adaptées aux systèmes complexes \elide.
    \item[][---] \elide la non-faisabilité des systèmes \elide.
}{\cite{farazi_model-based_2017}}

}{}

\enonce{AD32}{Assertion}{Plusieurs utilisateurs des langages d'architecture n'utilisent pas de langages formels, du moins pas ceux permettant d'exprimer la sémantique.}{}{P039}
\explication{AD32}{de l'assertion}{Plusieurs utilisateurs des langages d'architecture n'utilisent pas de langages formels, du moins pas ceux permettant d'exprimer la sémantique.}{
Dans l'article \textit{What Industry Needs from Architectural Languages: A Survey}, Ivano Malavolta, Patricia Lago \textit{et al.} affirment la sous-utilisation des langages formels~:
\quoteenfr
{AL tools from industry can cater to the needs of practitioners, but they lack the rigor of formal ALs. Conversely, formal ALs are not currently adopted by the industry, mainly for usability deficiencies. A bridge needs to be built to close this gap.}
{Les outils de langages d'architecture issus de l'industrie peuvent répondre aux besoins des praticiens, mais ils manquent de la rigueur des langages d'architecture formels. Inversement, les langages d'architecture formels ne sont pas actuellement adoptés par l'industrie, principalement en raison de problèmes d'utilisabilité. Il est nécessaire de créer un pont pour combler cet écart.}
{\cite{malavolta_what_2013}}
Aussi, cet article présente les résultats d'un sondage $(n=48)$ indiquant la préférence à l'égard de l'UML comparativement aux ADL. 
\quoteenfr
{What emerges from the collected data is that 42 percent of the respondents make exclusive use of UML and UML profiles, 35 percent use ADLs and UML together, while only 17 percent make and exclusively use ADLs when describing software architectures.}
{Il ressort des données recueillies que 42 \% des répondants utilisent exclusivement UML et les profils UML, 35 \% utilisent conjointement les ADL et UML, tandis que seulement 17 \% créent et utilisent exclusivement des ADL pour décrire les architectures logicielles.}
{\cite{malavolta_what_2013}}
Le même article démontre le désintérêt envers la fonctionnalité «~générer du code source~»~. Cette fonction découle directement de l'utilisation d'un langage formel~:
\quoteenfr
{While code generation is considered an interesting feature, it is rarely used. It is also not considered a very useful feature for future projects. The reasons provided by respondents are very different \elide.}
{Bien que la génération de code soit considérée comme une fonctionnalité intéressante, elle est rarement utilisée. Elle n'est pas non plus perçue comme très utile pour les projets futurs. Les raisons invoquées par les personnes interrogées sont très diverses \elide.}
{\cite{malavolta_what_2013}}

Des raisons y sont invoquées~:
\quoteenfrlist {
    \item[][---] \elide architecture description [is] on a too high level \elide.
    \item[][---] \elide[code generation] Requires too much details in the descriptions \elide.
    \item[][---] \elide[code generation] was not the intended goal of the architecture definition ! Architecture did have other focus \elide.
} {
    \item[][---] \elide la description de l'architecture est trop complexe \elide.
    \item[][---] \elide [la génération de code] exige trop de détails dans les descriptions \elide.
    \item[][---] \elide[La génération de code] n'était pas l'objectif visé par la définition de l'architecture ! L'architecture avait d'autres objectifs \elide.
}{\cite{malavolta_what_2013}}

Un autre article explique que les utilisateurs ne sont pas enclins à utiliser des langages formels qui permettent d'exprimer la sémantique~:
\quoteenfr
{The bulk of the current ADLs (e.g., Wright, LEDA, Darwin, PiLar, PRISMA, CONNECT, and SOFA) adopt a formal notation for specifying the behaviours of architectural elements. The notation is usually some process algebra (e.g.,  FSP [28], CSP [18], CCS [35], or $\pi$-calculus [36]), even though other formalisms are also used (e.g., Z [47]). Although it is important that they provide a formal means of specifying behaviours of architectures, process algebras, etc., are unfortunately not viewed favourably by practitioners [30].}
{La plupart des langages de description d'architectures (ADL) actuels (par exemple, Wright, LEDA, Darwin, PiLar, PRISMA, CONNECT et SOFA) adoptent une notation formelle pour spécifier le comportement des éléments architecturaux. Cette notation repose généralement sur une algèbre de processus (par exemple, FSP [28], CSP [18], CCS [35] ou le $\pi$-calcul [36]), même si d'autres formalismes sont également utilisés (par exemple, Z [47]). Bien qu'il soit important qu'ils fournissent un moyen formel de spécifier le comportement des architectures, les algèbres de processus, etc., sont malheureusement mal perçues par les praticiens [30].}
{\cite{ozkaya_are_2013}}
\figsmallen
{\textit{Notation of architecture documentation}}
{figure/chap1/documentation/umladl.png}
{(source~: \textit{Figure} 7 \textit{part} 1 \cite{rost_software_2013})}
{}{}{tb}{Langages d'architecture utilisés selon des professionnels}
Un sondage portant sur l'utilisation de la documentation d'architecture logicielle a été réalisé auprès de 147~professionnels informatiques provenant de divers pays et de multiples secteurs \cite{rost_software_2013}. Comme l'illustre la figure~B.18, les ADL formels ne sont utilisés qu'à environ 5~\%.
}{}


\paragraphe{P040}{\fl{Utiliser une base de données relationnelle pour maintenir la cohérence de la documentation\footnote{Le modèle est le résultat de la modélisation qui est une méthode. La modélisation en informatique regroupe un ensemble de techniques (diagramme entité-association, diagramme de flux de données et etc.).}}. D'abord, Scott R. Tilley, dans un article publié en 1993, souligne que la documentation a peu évoluée en trente ans et identifie ces trois problèmes \cite{tilley_documenting_1993}~: \enfr{in-the-small, usually out-of-date, provides just one perspective}{\elide est conçu pour de petites entités informatiques, est généralement obsolète et fournis une seule perspective \elide}. Il modélise une solution qui consiste à maintenir cohérents le code source et la documentation d'architecture. Il n'explicite pas l'utilisation d'une base de données, mais y explique son modèle relationnel et la forme des tuples $({objects}, {relationships})$. 
\sep
Ensuite, un article portant sur l'utilisation du S1000D dans le secteur critique des trains à haute vitesse indique utiliser Common Resource Data Base, une base de données relationnelle pour maintenir la cohérence \cite{weijiao_s1000d_2015}. À l'intérieur de la spécification du S1000D, il y est décrit que le Common Resource Data Base contient que la structure de données (objets d'informations), c'est-à-dire le modèle sans fournir les fonctions nécessaires aux applications l'utilisant \ccite{s1000d_spec_2024}[p.~1942].
}{AD33}{AD33}

\enonce{AD33}{Assertion}{Il y a des exemples où un modèle relationnel est utilisé pour maintenir la documentation cohérente.}{}{P040}
\explication{AD33}{de l'assertion}{Il y a des exemples où un modèle relationnel est utilisé pour maintenir la documentation cohérente.}{
Scott R. Tilley dans un article de 1993 décrit la situation de la documentation~:
\quoteenfr
{Documentation techniques have not really changed very much in thirty years ; most software documentation is still targeted towards in-the-small activities.}
{Les techniques de documentation n'ont pas vraiment beaucoup changé en trente ans ; la plupart des documentations logicielles sont encore destinées à des activités de petite envergure.}
{\cite{tilley_documenting_1993}}
Il y spécifie trois problèmes~: \enfr{in-the-small}{pour de petites entités informatiques}, \enfr{usually out-of-date}{généralement obsolète} et \enfr{provides just one perspective}{fournis une seule perspective}.
Voici la description qu'il fait du terme «~\textit{in-the-small}~»~:
\quoteenfr
{Most software documentation is documenting-in-the-small, since it typically describes the program at the algorithm and data structure level. For large legacy systems, an understanding of the structural aspects of the system's architecture is more important than any single algorithmic component.}
{La plupart des documentations logicielles sont de nature descriptive, car elles se limitent généralement à la description des algorithmes et des structures de données. Pour les grands systèmes existants, la compréhension des aspects structurels de l'architecture du système est plus importante que celle de tout composant algorithmique pris individuellement.}
{\cite{tilley_documenting_1993}}
L'article propose également d'établir des relations entre les différents modèles~:
\quoteenfr
{Our approach is to make the documentation a «~live~» representation of the source code, not separate text. If the documentation is the structure, and vice versa, no discrepancy exists. The rest of this section describes the basis for documenting software structure~: virtual subsystem stratifications (VSS's) and software interconnection models (IM's).}{Notre approche consiste à faire de la documentation une représentation «~vivante~» du code source, et non un texte séparé. Si la documentation reflète la structure, et, inversement, il n'y a pas de contradiction. La suite de cette section décrit les bases de la documentation de la structure logicielle : les stratifications de sous-systèmes virtuels et les modèles d'interconnexion logicielle.}

Un autre article décrit l'utilisation du S1000D pour le secteur des trains à haute vitesse~: 
\quoteenfr{Data module (DM) and common resource data base (CSDB) are two essential concepts of S1000D, a standardized IETM organizes its information with DM, while manages it in CSDB. \elide After compilation, the technical information (manuals) is stored in the relational CSDB with a structured and nonredundancy manner.
}{
Le module de données (MD) et la base de données de ressources communes (CSDB) sont deux concepts essentiels de la norme S1000D. Cette norme, qui définit un manuel d'utilisation de l'information (MUI), organise ses informations dans le MD et les gère dans la CSDB. \elide Après la compilation, les informations techniques (manuels) sont stockées dans la CSDB relationnelle de manière structurée et non redondante.}{\cite{weijiao_s1000d_2015}}

Voici une description de la base de données \textit{Common Resource Data Base} (CSDB) utilisée dans le S1000D~:
\quoteenfr{S1000D does not specify the design and implementation rules for a CSDB. It only specifies the data structure of the information objects, which is independent from any software implementation.}{La norme S1000D ne précise pas les règles de conception et de mise en œuvre d'une CSDB. Elle spécifie uniquement la structure des données des objets d'information, indépendamment de toute implémentation logicielle.}{\ccite{s1000d_spec_2024}[p.~1942]}

}{}

\subsubsection{Des techniques de traitement automatique des langues}

\paragraphe{P041}{Cette sous-section aborde les machines en mesure de traiter les documents, peu importe les langages employés. Il y aura une présentation des domaines, des techniques et des modèles de langage ainsi que des limitations existantes.}{}{} 

\paragraphe{P042}{Pour éviter les ambiguïtés, ce paragraphe regroupe des définitions, puis un diagramme des relations logiques entre des ensembles (figure~1.1) : 
\figmedium
{Diagramme de Venn sur les domaines, techniques et modèle de langage}
{figure/chap1/documentation/taxa.png}
{}{}{}{tb}
\begin{itemize}
    \item intelligence artificielle (IA ou AI, \textit{Artificial Intelligence} en anglais~: \enfr{\elide<discipline> research and development of mechanisms and applications of AI systems \elide.}{\elide<discipline> recherche et développement des mécanismes et applications des systèmes d'IA \elide.} \cite{iso_22989_2022};
    \item traitement automatique des langues (TAL ou NLP, \textit{Natural Language Processing} en anglais)~: «~Technique d'apprentissage automatique qui permet à l'ordinateur de comprendre le langage humain.~» \cite{Gdt_NLP_0} ;
    \item apprentissage automatique (AA ou ML, \textit{Machine Learning} en anglais)~: \enfr{\elide process of optimizing model parameters through computational techniques \elide.}{\elide processus d'optimisation des paramètres d'un modèle par des techniques de calcul \elide.} \cite{iso_22989_2022};
    \item réseaux de neurones (RN ou NN, \textit{Neural Networks en anglais)}~: \enfr{\elide network of one or more layers of neurons \elide.}{\elide réseau d'une ou plusieurs couches de neurones \elide} \cite{iso_22989_2022};
    \item apprentissage profond (AP ou DL, \textit{Deep Learning} en anglais)~: «~Mode d'apprentissage automatique généralement effectué par un réseau de neurones artificiels composé de plusieurs couches de neurones hiérarchisées \elide.~» \cite{Gdt_DL_0} ;
    \item grand modèle de langage (GML ou LLM, \textit{Large Language Model} en anglais)~: «~Modèle de langage constitué par apprentissage profond à partir de mégadonnées.~» \cite{Gdt_llm_0} ;
    \item boîte noire~: «~Système fermé dont on ignore le fonctionnement interne, qui n'est connu que par son comportement extérieur en fonction des sollicitations qui lui sont appliquées.~» \cite{Gdt_noire_0}.
\end{itemize}

}{AD34, AD35, AD36, AD37, AD38, AD39}{AD34}

\enoncewithoutexp{AD34}{Définition}{Traitement automatique des langues (TAL ou NLP, \textit{Natural Language Processing} en anglais)}{
«~Technique d'apprentissage automatique qui permet à l'ordinateur de comprendre le langage humain.~» \cite{Gdt_NLP_0}
}{P042}

\enoncewithoutexp{AD35}{Définition}{Apprentissage profond (AP ou DL~: \textit{Deep Learning})}{
«~Mode d'apprentissage automatique généralement effectué par un réseau de neurones artificiels composé de plusieurs couches de neurones hiérarchisées \elide.~» \cite{Gdt_DL_0}
}{P042}

\enoncewithoutexp{AD36}{Définition}{Grand modèle de langage (GML ou LLM, \textit{Large Language Model} en anglais)}{
«~Modèle de langage constitué par apprentissage profond à partir de mégadonnées.~» \cite{Gdt_llm_0}
}{P042}

\enoncewithoutexp{AD37}{Définition}{Boîte noire}{
«~Système fermé dont on ignore le fonctionnement interne, qui n'est connu que par son comportement extérieur en fonction des sollicitations qui lui sont appliquées.~» \cite{Gdt_noire_0}
}{P042}

\enonce{AD38}{Assertion}{Plusieurs concepts de l'intelligence artificielle sont logiquement interreliés.}{}{P042}
\explication{AD38}{de l'assertion}{Plusieurs concepts de l'intelligence artificielle sont logiquement interreliés.}{
La littérature fournit ces informations qui ont servi à construire cette assertion~:
\begin{itemize}
    \item définition d'intelligence artificielle (IA)~: \enfr{\elide<discipline> research and development of mechanisms and applications of AI systems \elide.}{\elide<discipline> recherche et développement des mécanismes et applications des systèmes d'IA \elide.} \cite{iso_22989_2022};
    \item définition de système d'intelligence artificielle~: \enfr{\elide engineered system that generates outputs such as content, forecasts, recommendations or decisions for a given set of human-defined objectives \elide}{\elide système conçu pour générer des résultats tels que du contenu, des prévisions, des recommandations ou des décisions en fonction d'un ensemble d'objectifs définis par l'humain \elide} \cite{iso_22989_2022};
    \item définition d’apprentissage automatique~: \enfr{\elide process of optimizing model parameters through computational techniques \elide.}{\elide processus d'optimisation des paramètres d'un modèle par des techniques de calcul \elide.} \cite{iso_22989_2022};
    \item définition de réseau de neurones (RN)~: \enfr{\elide network of one or more layers of neurons \elide.}{\elide réseau d'une ou plusieurs couches de neurones \elide} \cite{iso_22989_2022};
    \item synonymes de réseau de neurones~: \textit{neural net} et \textit{artificial neural network} \cite{iso_22989_2022} ;
    \item définition de réseau de neurones~: \enfr{\elide approach to creating rich hierarchical representations through the training of neural networks with many hidden layers \elide}{\elide approche pour créer des représentations hiérarchiques riches grâce à l'entraînement de réseaux neuronaux avec de nombreuses couches cachées \elide}\cite{iso_22989_2022};
    \item synonyme d'apprentissage profond~: \textit{deep neural network learning} \cite{iso_22989_2022};
    \item description d'apprentissage profond (AP)~: \enfr{Deep learning is a broad family of techniques for machine learning in which hypotheses take the form of complex algebraic circuits with tunable connection strengths.}{L'apprentissage profond est une famille élargie de techniques pour l'apprentissage machine où les hypothèses prennent la forme de circuits complexes d'algèbre avec des liaisons multiples.} \ccite{norvig_artificial_2021}[p.~1378];
    \item définition 1 de grand modèle de langage (GML)~: «~Modèle de langage constitué par apprentissage profond à partir de mégadonnées.~» \cite{Gdt_llm_0};
    \item définition 2 de grand modèle de langage~: \enfr{\elide LLMs like GPT (Generative Pre-trained Transformer) use deep learning techniques \elide.}{\elide Les GML comme GPT (Generative Pre-trained Transformer) utilisent des techniques d'apprentissage profond \elide.} \footnote{University of British Columbia, site web, «~Glossary of GenAI Terms~», \url{https://ai.ctlt.ubc.ca/resources/glossary-of-genai-terms/} (consulté le 2025-10-16).}.
\end{itemize}
\figmediumen
{\textit{The relationship between the different disciplines}}
{figure/chap1/documentation/gazit.png}
{(modifié la \textit{Figure}~1.2 \cite{gazit_mastering_2024})}
{}
{
\begin{tablenotes}
\footnotesize
\item[a] {Dans l'ouvrage, le diagramme est en couleur, puisqu'il comporte des noms, une version en gris a été retenu ici.}
\end{tablenotes}
}{tb}{Diagramme de Venn sur l'IA et le TAL}
\figmediumen
{\textit{A Venn diagram showing how deep learning is a kind of representation learning \elide}}
{figure/chap1/documentation/goodfellow2.png}
{(modifié la \textit{Figure}~1.4 \cite{goodfellow_deep_2016})}
{}{
\begin{tablenotes}
\footnotesize
\item[a] {Diagramme est en partie.}
\end{tablenotes}
}{tb}{Diagramme de Venn sur l'IA}
Dans l'ouvrage \titleen{Mastering NLP from foundations to LLMs}{Maîtriser le TAL des fondations aux GML} inclut un diagramme de Venn (figure~B.19) représentant les relations entre les différentes disciplines d'IA \ccite{gazit_mastering_2024}[p.~25]. Tandis que l'ouvrage  \titleen{Deep learning}{Apprentissage profond} d'Ian Goodfellow, Yoshua Bengio et Aaron Courville, présentent notamment ce type de diagramme de Venn \ccite{goodfellow_deep_2016}[p.~9]. Il s'agit de la figure~B.20. 
}{}

\figannexe{

}

\enonce{AD39}{Assertion}{Les domaines du traitement automatique des langues et de l'intelligence artificielle se croisent sur le plan historique.}{}{P042}
\explication{AD39}{de l'assertion}{Les domaines du traitement automatique des langues et de l'intelligence artificielle se croisent sur le plan historique.}{
Les ouvrages \titleen{Mastering NLP from foundations to LLMs}{} et \titleen{Artificial Intelligence A Modern Approach}{}proposent une synthèse de l'histoire du TAL \ccite{gazit_mastering_2024}[p.~20-21][norvig_artificial_2021][p.~63-79]. En voici un très bref résumé fondé sur les éléments présentés dans ces ouvrages~:
\begin{itemize}
    \item Dans les années 1950, les systèmes à base de règles de production étaient utilisés \footnote{En anglais~: \textit{rule-based systems}}. Ces systèmes sont devenus plus sophistiqués, grâce à des bases de connaissances plus grandes, et ce, jusqu'à la fin des années 1980.
    \item À la fin des années 1980, les méthodes statistiques se sont imposées sous l'appellation d'apprentissage automatique (ML)\cite{Gdt_ml_0}.
    \item Les années 2010 marquent le début de l'apprentissage profond (DL, \textit{Deep Learning} en anglais), ce qui a relancé la popularité des réseaux de neurones (NN, \textit{Neural Network} en anglais). Ces derniers existent depuis les années 1950 et ont déjà été populaires dans les années 1980.
\end{itemize}
  
Dans l'ouvrage \textit{Deep learning}, l'auteur affirme que l'apprentissage profond (AP) est simplement le dernier nom donné, voici la citation complète ;
\quoteenfr
{Deep learning only appears to be new, because it was relatively unpopular for several years preceding its current popularity, and because it has gone through many different names, only recently being called "deep learning." The field has been rebranded many times, reflecting the influence of different researchers and different perspectives.}
{L'apprentissage profond n'est une nouveauté qu'en apparence, car il a été relativement impopulaire pendant plusieurs années avant sa popularité actuelle, et parce qu'il a connu de nombreux noms, jusqu'à être récemment appelé "~apprentissage profond~". Ce domaine a connu de nombreuses refontes, reflétant l'influence de différents chercheurs et perspectives.}
{\ccite{goodfellow_deep_2016}[p.~12]}

L'ensemble de ces techniques tente de produire des modèles. Dans le contexte du TAL, le résultat est appelé un modèle de langage. Parmi ces modèles de langage apparaissent les grands modèles de langage (GML) à la fin des années 2010.
}{}

\paragraphe{P043}{\fl{Toute la documentation écrite en langage non formel peut alimenter le développement de techniques de traitement automatique des langues puisque ces techniques traitent ce type d'information.} D'abord, plusieurs professionnels informatiques affirment que l'architecture et les exigences s'expriment couramment en langage naturel. Selon deux sondages, cette proportion atteint 90~\% $(n=109)$ pour l'architecture \cite{rost_software_2013} et 69~\% $(n=117)$ pour les exigences \cite{kassab_current_2022}. 
\sep
Ensuite, les documents de spécification des exigences sont ceux qui emploient le plus les techniques de TAL. C'est ce qui ressort d'une revue de littérature portant sur des articles publiés entre 1983 et 2019 \cite{zhao_natural_2022}. 
\sep
Pour finir, les techniques de TAL sont de plus en plus exploitées dans la littérature scientifique. Une tentative de taxinomie est réalisée sur le TAL à partir de \mbox{quatre-vingt mille} articles provenant de l'Association for Computational Linguistics et le nombre d'articles suit une courbe exponentielle entre 1952 et 2023 \cite{schopf_exploring_2023}.
}{AD40, AD41, AD42}{AD40}

\enonce{AD40}{Assertion}{Plusieurs professionnels informatiques affirment que l'architecture et les exigences s'expriment couramment en langage naturel.}{}{P043}
\explication{AD40}{de l'assertion}{Plusieurs professionnels informatiques affirment que l'architecture et les exigences s'expriment couramment en langage naturel.}{
Un article de D. Rost  \cite{rost_software_2013} sonde l'utilisation de la documentation d'architecture logicielle. Selon son sondage, auquel 109~des 147~professionnels informatiques ont répondu, 90~\% utilisent un langage non formel.
(voir AD28 p.\href{AD28}{\pageref{explication:AD28}} figure~B.8 «~Comment les informations architecturales sont-elles décrites ?~»)

Un article publié en 2022 décrit la pratique des professionnels informatiques en ce qui concerne les exigences \cite{kassab_current_2022}. Il analyse quatre sondages qui ont été réalisés en 2003 $(n=?)$, en 2008 $(n=?)$, en 2013 $(n=?)$ et en 2020 $(n=117)$. Au cours de ces sondages, le taux des professionnels informatiques jugeant que leurs exigences deviennent de plus en plus formelles est passé de 45~\% à 69~\%.
}{}

\enonce{AD41}{Assertion}{Au moins jusqu'en 2020, les documents de spécification des exigences est la catégorie de documents utilisant le plus les techniques de TAL.}{}{P043}
\explication{AD41}{l'assertion}{Au moins jusqu'en 2020, les documents de spécification des exigences est la catégorie de documents utilisant le plus les techniques de TAL.}{
Une revue de littérature publiée en 2022 portant sur l'utilisation du TAL pour les exigences \cite{zhao_natural_2022} analyse un corpus de 404~articles publiés entre 1983 et 2019 (dont plus de la moitié [214 articles] a été publiée après 2012. Parmi l'ensemble des articles recensés, 271~articles portent sur une proposition de solution, tandis que 239~articles traitent de spécification des exigences système ou logicielles comme type de documents, les autres types sont des \textit{use cases}, des \textit{user stories}, etc. D'ailleurs, les fonctions et les techniques les plus utilisées sont~: 
étiquetage grammatical (\textit{POS Tagging}) ($\approx$185), analyse lexicale (\textit{Tokenization}) ($\approx$80), analyse syntaxique (\textit{Parsing}) ($\approx$75), suppression des mots vides (\textit{Stop-Words Removal}) ($\approx$75), extraction de termes (\textit{Term Extraction}) ($\approx$70), racinisation (\textit{Stemming}) ($\approx$70), lemmatisation (\textit{Lemmatizaation}) ($\approx$50), mesures de similarité (\textit{Similarity Measures}) ($\approx$50), TF-IDF ($\approx$40), etc.
}{}

\enonce{AD42}{Énoncé}{Les techniques de TAL sont de plus en plus utilisées dans la littérature scientifique.}{}{P043}
\explication{AD42}{de l'énoncé}{Les techniques de TAL sont de plus en plus utilisées dans la littérature scientifique.}{
Les utilisations des techniques de TAL croissent, il y a de plus en plus de propositions et d'expérimentations. Des articles publiés entre 2014 et 2023 démontrent cette progression~:  
\begin{itemize}
    \item un article publié en 2014 présente une solution pour générer un sommaire depuis le code source \cite{mcburney_automatic_2015} ;
    \item un article publié en 2016 décrit qu'il est possible d'extraire les concepts des documents, tout en soulignant le caractère mitigé des résultats obtenus \cite{menard_concept_2016} ;
    \item un article publié en 2018 démontre qu'il est possible de se servir des techniques TAL comme instruments de suggestions ou de complétions pour la documentation \cite{terrasa_using_2018} ;
    \item un article publié en 2023 énumère tous les domaines d'application du TAL \cite{schopf_exploring_2023}. Le TAL couvre 12~domaines d'application, chacun divisé. Il s'agit d'une tentative de taxonomie réalisée à partir de 80 000 articles provenant de la bibliothèque \textit{Anthology} (qui est consacrée au TAL) de l'Association for Computational Linguistics (ACL). Le nombre d'articles contenus dans cette bibliothèque suit une courbe exponentielle entre 1952 et 2023.
\end{itemize}
\vspace{-1\baselineskip}
}{}

\paragraphe{P044}{\fl{Plusieurs techniques de TAL\footnote{Pour rappel, le traitement automatique des langues comprend les techniques d'apprentissage automatique, de réseaux de neurones et d'apprentissage profond.} ont des lacunes scientifiques.} D'abord, il n'y a pas de règles exactes pour sélectionner adéquatement la chaîne ou sélectionner la classe de modèle à entraîner \ccite{gazit_mastering_2024}[p.~176]. En fait, une chaîne (\textit{pipeline}) prépare, développe et exécute un modèle de langage qui servira pour le TAL. La pratique actuelle se contente de fournir des exemples de chaînes et certains experts appellent à utiliser l'intuition \ccite{norvig_artificial_2021}[p.~1310]. À la fin, les résultats obtenus de plusieurs de ces techniques de TAL sont en fait des modèles de type «~boîte noire~» qui sont plus difficilement explicables. \cite{hsieh_comprehensive_2024}.
\sep
Néanmoins, il y a de plus en plus d'articles qui s'intéressent à comprendre les modèles \cite{barredo_explainable_2020} et certains peuvent fournir des explications demeurant incertaines, voire le plus souvent douteuses. Cynthia Rudin souligne des problèmes existants avec plusieurs de ces explications~:
\quoteenfrlist{
\item Explainable ML methods provide explanations that are not faithful to what the original model computes. 
\item Explanations often do not make sense or do not provide enough detail to understand what the black box is doing.
}{
\item Les méthodes d'apprentissage automatique explicables fournissent des explications qui ne reflètent pas fidèlement les calculs du modèle original.
\item Ces explications sont souvent incohérentes ou insuffisamment détaillées pour comprendre le fonctionnement des modèles de type « boîte noire~».
}{\cite{rudin_stop_2019}}
Pour finir, il y a des manquements importants pour assurer la reproductibilité. 
\begin{itemize}
    \item Un article publié en 2025 énonce cinq items essentiels à la reproductibilité~: versionner le code, avoir accès aux données, versionner les données, avoir les journaux de l'expérimentation ainsi que la chaîne qui a servi à créer le modèle \cite{semmelrock_reproducibility_2025}.
    \item Une analyse sur 24~expérimentations importantes sur des modèles conclut que \cite{reuel_betterbench_2024}~:
    \begin{itemize}
        \item 14 n'ont pas indiqué de signification statistique;
        \item 17 ne comportaient pas de scripts pour la réplication des résultats;
        \item la plupart souffraient d'une documentation insuffisante.
    \end{itemize}
    \item Plusieurs experts vont jusqu'à affirmer qu'il n'y a pas d'accumulation fiable des connaissances \ccite{silberztein_linguistic_2024}[p.~22][weidinger_toward_2025][].
\end{itemize}
\vspace{-1\baselineskip}
}{AD43, AD44, AD45, AD46, AD47}{AD43}

\enonce{AD43}{Assertion}{Il n'y a pas de règles exactes pour sélectionner une chaîne ou une classe de modèle à entraîner.}{}{P044}
\explication{AD43}{de l'assertion}{Il n'y a pas de règles exactes pour sélectionner une chaîne ou une classe de modèle à entraîner.}{
L'ouvrage \textit{Mastering NLP from foundations to LLMs} montre une chaîne typique (figure~B.21) pour l'entraînement d'un modèle avec de l'apprentissage automatique \ccite{gazit_mastering_2024}[p.~176].
\figen
{\textit{The structure of a typical ML pipeline}}
{figure/chap1/documentation/chaine.png}
{(source~: Figure~5.2 \ccite{gazit_mastering_2024}[p.~176])}
{}{}{tb}{Procédé typique pour la création d'un modèle de langage}
Dans l'ouvrage \textit{Artificial Intelligence A Modern Approach}, il est écrit qu'il n'existe pas de procédé parfaitement établi~: 
\quoteenfr
{With cleaned data in hand and an intuitive feel for it, it is time to build a model. \elide There is no guaranteed way to pick the best model class, but there are some rough guidelines.}
{Avec des données nettoyées et une compréhension intuitive de celles-ci, il est temps de construire un modèle. \elide Il n'existe aucune méthode infaillible pour choisir la meilleure classe de modèle, mais voici quelques lignes directrices générales.}
{\ccite{norvig_artificial_2021}[p.~1310]}
\vspace{-1\baselineskip}
}{}


\enonce{AD44}{Assertion}{Plusieurs techniques de TAL conduisent à la création de modèles de type boîte noire.}{}{P044}
\explication{AD44}{de l'assertion}{Plusieurs techniques de TAL conduisent à la création de modèles de type boîte noire.}{
En 2024, un guide consacré à l'explication des modèles d'IA est publié \cite{hsieh_comprehensive_2024}. Rédigé par un collectif de 27~auteurs issus, majoritairement, du milieu universitaire, il énumère des modèles de types boîte blanche et boîte noire \footnote{La boîte est blanche est le contraire de la boîte noire. }~:  
\begin{itemize}
    \item Les modèles «~boîtes blanches~» sont~: 
    \enfr{Linear Regression, Logistic Regression, Decision Trees, Rule-based Systems, K-Nearest Neighbors (KNN), Naive Bayes Classifiers, Generalized Additive Models (GAMs)}{Régression linéaire, régression logistique, arbres de décision, systèmes à base de règles, K-plus proches voisins, classificateurs bayésiens naïfs, modèles additifs généralisés}. 
    \item Les modèles «~boîtes noires~» sont~: 
    \enfr {Neural Networks (e.g., Deep Learning Models), Support Vector Machines (SVMs), Ensemble Methods (e.g., Random Forests, Gradient Boosting Machines), Transformer Models (e.g., BERT, GPT), Graph Neural Networks (GNNs)}{ Réseaux de neurones (par exemple, modèles d'apprentissage profond), machines à vecteurs de support (SVM), méthodes d'ensemble (par exemple, forêts aléatoires, machines à gradient boosting), modèles de transformateurs (par exemple, BERT, GPT), réseaux de neurones graphiques}. 
\end{itemize}
\vspace{-1\baselineskip}
}{}

\figannexe{%
\figen
{\textit{Evolution of the number of total publications whose title, abstract and/or keywords refer to the field of XAI during the last years.}}
{figure/chap1/documentation/xai.png}
{(source~: \textit{Figure}~1 \cite{barredo_explainable_2020})}
{}{}{tb}{Nombre d'articles voulant expliquer l'IA entre 2012 et 2019}}
\enonce{AD45}{Assertion}{Les tentatives pour expliquer ces modèles sont de plus en plus nombreuses, cela montre l'insuffisance des explications.}{}{P044}
\explication{AD45}{de l'assertion}{Les tentatives pour expliquer ces modèles sont de plus en plus nombreuses, cela montre l'insuffisance des explications.}{
Un article publié en 2020 décrit un domaine récent de l'AI, l'\textit{Explainable Artificial Intelligence} (XAI) \cite{barredo_explainable_2020}. L'article réalise des recherches sur Scopus, lesquelles mettent en évidence une  croissance du nombre de publications (figure~B.22).
Dans un article publié en 2019, Cynthia Rudin expose des problèmes d'interprétabilité, d'explicabilité, de compréhensibilité et de transparence avec l'apprentissage machine~:
\quoteenfrlist {
    \item[][---] It is a myth that there is necessarily a trade-off between accuracy and interpretability.
    \item[][---] Explainable ML methods provide explanations that are not faithful to what the original model computes.
    \item[][---] Explanations often do not make sense or do not provide enough detail to understand what the black box is doing.
    \item[][---] Black box models are often not compatible with situations where information outside the database needs to be combined with a risk assessment.
    \item[][---] Black box models with explanations can lead to an overly complicated decision pathway that is ripe for human error.
} {
    \item[][---] Il est faux de croire qu'il existe nécessairement un compromis entre précision et interprétabilité.
    \item[][---] Les méthodes d'apprentissage automatique explicables fournissent des explications qui ne reflètent pas fidèlement les calculs du modèle original.
    \item[][---] Ces explications sont souvent incohérentes ou insuffisamment détaillées pour comprendre le fonctionnement de la boîte noire.
    \item[][---] Les modèles de type «~boîte noire~»sont souvent incompatibles avec les situations où des informations externes à la base de données doivent être intégrées à une évaluation des risques.
    \item[][---] Les modèles de type «~boîte noire~», même avec explications, peuvent engendrer un processus de décision excessivement complexe, propice aux erreurs humaines.
}{\cite{rudin_stop_2019}}
\vspace{-1\baselineskip}
}{}

\enonce{AD46}{Assertion}{L'absence de culture sur la reproductibilité limite l'amélioration des connaissances.}{}{P044}
\explication{AD46}{l'assertion}{L'absence de culture sur la reproductibilité limite l'amélioration des connaissances.}{
En 2025, un article décrit les problèmes de reproductibilité des modèles en AA \cite{semmelrock_reproducibility_2025} et énonce cinq items essentiels à la reproductibilité ; 
\begin{itemize}
\item \sfren{versionner le code}{code versioning};
\item \sfren{avoir accès aux données}{data access};
\item \sfren{versionner les données}{data versioning};
\item \sfren{avoir les journaux de l'expérimentation}{experiment logging};
\item \sfren{la chaîne qui a servi à créer le modèle}{pipeline creation}.
\end{itemize}

En 2025, le populaire \textit{Artificial Intelligence Index Report} de l'Université de Stanford cite un article décrivant les problèmes de reproductibilité et d'évaluation sur les expérimentations \ccite{maslej_artificial_2025}[p.~101]. Le nom de l'article cité est \textit{BetterBench} \cite{reuel_betterbench_2024}. Au total, 24~expérimentations majeures ont été analysées et il en ressort que~:
\begin{itemize}
\item 14~n'ont pas indiqué de signification statistique (\textit{failed to report statistical significance});
\item 17~ne comportaient pas de scripts pour la réplication des résultats (\textit{lacked scripts for result replication});
\item la plupart souffraient d'une documentation insuffisante (\textit{most suffered from inadequate documentation}).
\end{itemize}
\vspace{-1\baselineskip}
}{}
%\footnote{Il est l'auteur de l'instrument NooJ, un environnement de développement linguistique.}
\enonce{AD47}{Assertion}{Plusieurs experts affirment qu'il n'y a pas d'accumulation fiable des connaissances dans le domaine du TAL.}{}{P044}
\explication{AD47}{de l'assertion}{Plusieurs experts affirment qu'il n'y a pas d'accumulation fiable des connaissances dans le domaine du TAL.}{
En 2024, Max Silberztei, auteur de l'ouvrage \textit{Linguistic Resources for Natural Language Processing: On the Necessity of Using Linguistic Methods to Develop NLP Software}, met en garde contre les modèles de type \emphase{boîte noire} \ccite{silberztein_linguistic_2024}[p.~22]~:
\quoteenfr
{Using training-corpus-based "black box" methods that produce satisfactory results, but without any power of explanation nor generalization, is a viable approach to non-critical NLP software engineering problems. However, it is time for all scientists—both linguists and computer scientists—to realize that this approach is the opposite of what scientific approaches require~: reliable accumulation of knowledge.}
{L'utilisation de méthodes de type «~boîte noire~» basées sur des corpus d'entraînement, produisant des résultats satisfaisants, mais dépourvus de toute capacité d'explication ou de généralisation, constitue une approche viable pour résoudre les problèmes d'ingénierie logicielle non critiques en TAL. Cependant, il est temps que tous les scientifiques, linguistes et informaticiens, prennent conscience que cette approche est à l'opposé de ce que les approches scientifiques exigent~: une accumulation fiable de connaissances.}
{\ccite{silberztein_linguistic_2024}[p.~22]}
En 2025, dans l'article  \textit{Toward an evaluation science for generative AI systems} \cite{weidinger_toward_2025}, Laura Weidinger \textit{et al.} affirment que la manière actuelle d'expérimenter n'est tout simplement pas suffisante pour améliorer nos connaissances~:
\quoteenfr
{For AI evaluation to mature into a proper "science," it must meet certain criteria. Sciences are marked by having theories about their targets of study, which can be expressed as testable hypotheses. Measurement instruments to test these hypotheses must provide experimental consistency (i.e., reliability, internal validity) and generalizability (i.e., external validity). Finally, sciences are marked by iteration: Over time, measurement approaches and instruments are refined and new insights are uncovered. Collectively, these properties of sciences contrast sharply with the practice of rapidly developing static benchmarks for evaluating generative AI systems, while anticipating that within a few months such benchmarks will become much less useful or obsolete.}
{Pour que l'évaluation de l'IA devienne une véritable «~science~», elle doit répondre à certains critères. Les sciences se caractérisent par des théories sur leurs objets d'étude, lesquelles peuvent être exprimées sous forme d'hypothèses testables. Les instruments de mesure utilisés pour tester ces hypothèses doivent garantir la cohérence expérimentale (fiabilité, validité interne) et la généralisabilité (validité externe). Enfin, les sciences se caractérisent par l'itération~: au fil du temps, les approches et les instruments de mesure sont affinés et de nouvelles connaissances sont mises au jour. Collectivement, ces propriétés des sciences contrastent fortement avec la pratique consistant à développer rapidement des référentiels statiques pour évaluer les systèmes d'IA générative, tout en anticipant que, d'ici quelques mois, ces référentiels deviendront beaucoup moins utiles, voire obsolètes.}
{\cite{weidinger_toward_2025}}
\vspace{-1\baselineskip}
}{}

\subsubsection{Des techniques d'apprentissage profond}

\paragraphe{P045}{Certaines techniques de TAL sont si répandues qu'elles nécessitent leurs propres sous-sections, il s'agit des techniques d'apprentissage profond. Il est expliqué que, même avec leurs limitations, les professionnels informatiques sont prêts à les utiliser et à les offrir. Il est présenté qu'il existe une absence d'expérimentations démontrant sans équivoque les avantages pour la documentation et une absence d'analyses approfondies des risques. La sous-section se termine en envisageant l'intégration ou la coexistence de ces techniques.}{}{}

\paragraphe{P046}{\fl{Malgré certaines limites en matière de reproductibilité, l'apprentissage profond demeure utile pour de nombreux individus.} Les techniques d'apprentissage profond comptent parmi les techniques les plus utilisées en TAL \ccite{norvig_artificial_2021}[p.~1378]. Elles fournissent des résultats jugés satisfaisants. Les experts du domaine vont jusqu'à changer les bancs d'essai lorsqu'ils leur semblent suffisamment réussis \ccite{maslej_artificial_2025}[p.~137]. Un sondage mené auprès de quatre-vingt-dix mille professionnels informatiques indique qu'une majorité utilise les instruments basés sur les grands modèles de langage (GML), donc produits par l'apprentissage profond\ccite{maslej_artificial_2024}[p.~269]. Les réponses indiquent notamment que~:
\begin{itemize}
    \item 82~\% s'en servent pour écrire du code source ;
    \item 48~\% s'en servent comme une aide envers du code source ;
    \item 55~\% voudraient s'en servir pour tester du code source ;
    \item 50~\% voudraient s'en servir pour documenter du code source.
\end{itemize}
\vspace{-1\baselineskip}
}{AD48, AD49, AD50}{AD48}

\enonce{AD48}{Citation}{L'apprentissage profond crée des grands modèles de langage qui améliorent les résultats en traitement automatique des langues.}{}{P046}
\explication{AD48}{la citation}{L'apprentissage profond crée des grands modèles de langage qui améliorent les résultats en traitement automatique des langues.}{
L'AP est l'approche la plus utilisée en TAL et est celle qui a eu le plus d'impact, selon l'ouvrage d’\textit{Artificial Intelligence A Modern Approach}, en voici deux extraits~:
\quoteenfr
{Deep learning is currently the most widely used approach for applications such as visual object recognition,  machine translation, speech recognition, speech synthesis, and image synthesis; it also plays  a significant role in reinforcement learning applications. \elide Deep learning has also had a huge impact on natural language processing (NLP)  applications such as machine translation and speech recognition. Some advantages of deep  learning for these applications include the possibility of end-to-end learning, the automatic  generation of internal representations for the meanings of words, and the interchangeability  of learned encoders and decoders.}
{L'apprentissage profond est actuellement l'approche la plus utilisée pour des applications, telles que la reconnaissance visuelle d'objets, la traduction automatique, la reconnaissance vocale, la synthèse vocale et la synthèse d'images ; il joue également un rôle important dans les applications d'apprentissage par renforcement. \elide L'apprentissage profond a également eu un impact considérable sur les applications de traitement automatique du langage naturel (TAL), telles que la traduction automatique et la reconnaissance vocale. Parmi ses avantages pour ces applications, on peut citer la possibilité d'un apprentissage de bout en bout, la génération automatique de représentations internes du sens des mots et l'interchangeabilité des encodeurs et décodeurs appris.}
{\ccite{norvig_artificial_2021}[p.~1378,1439]}
\vspace{-1\baselineskip}
}{}

\enonce{AD49}{Assertion}{Les bancs d'essai pour les GML sont davantage réussis, ce qui rend nécessaire le développement de nouveaux bancs d'essai.}{}{P046}
\explication{AD49}{de l'assertion}{Les bancs d'essai pour les GML sont davantage réussis, ce qui rend nécessaire le développement de nouveaux bancs d'essai.}{
L'\titleen{Artificial Intelligence Index Report}{} de 2025 contient des informations sur la performance des techniques \cite{maslej_artificial_2025}. Il y est discuté majoritairement de GML, comme l'indique le tableau \titleen{Timeline: Significant Model and Dataset Release}{} \ccite{maslej_artificial_2025}[p.~87-92]. Voici un sommaire du tableau (tableau B.10), GML y apparaît 14 fois sur un total de 38.
Dans le même rapport se trouve un comparatif de différents bancs d'essai à la \textit{Figure}~2.1.33 \ccite{maslej_artificial_2025}[p.~93]. Les résultats sont disposés de 0\% à 120\% où 100\% équivaudraient aux compétences humaines. Il est observé qu'il y a un banc d'essai qui disparaît en 2021 et deux en 2022. Tandis que d'autres apparaissent, dont deux en 2019, un en 2021 puis deux en 2023. Les auteurs expliquent ces changements de la façon suivante~:
\quoteenfr
{In recent years, the reasoning abilities of AI systems have  advanced so much that older benchmarks like SQuAD (for  textual reasoning) and VQA (for visual reasoning) have become  saturated, indicating a need for more challenging reasoning tests.}
{Ces dernières années, les capacités de raisonnement des systèmes d'IA ont tellement progressé que les anciens \textit{benchmarks}, comme SQuAD (pour le raisonnement textuel) et VQA (pour le raisonnement visuel), sont devenus saturés, ce qui indique un besoin de tests de raisonnement plus difficiles.}
{\ccite{maslej_artificial_2025}[p.~137]}
\vspace{-1\baselineskip}
}{}

\figannexe{
\tabtaben
{}
{
\begin{tabular}{p{4.8cm}p{4.8cm}p{4.8cm}}
    \begin{itemize}
        \item \textit{LLM} (14)
        \item \textit{dataset} (2)
        \item \textit{model/dataset} (1)
        \item \textit{multimodal} (3)
        \item \textit{language} (2)
        \item \textit{vision} (1)
    \end{itemize}
    &
    \begin{itemize}
        \item \textit{text-to-text} (1)
        \item \textit{text-to-image} (3)
        \item \textit{text-to-video} (2)
        \item \textit{image-to-video} (1)
        \item \textit{agentic capability} (1)
        \item \textit{math} (1)
    \end{itemize}
    &
    \begin{itemize}
        \item \textit{text-to-song} (1)
        \item \textit{song-to-song} (1)
        \item \textit{text-to-podcast} (1)
        \item \textit{tool} (1)
        \item \textit{iPhone feature} (1)
        \item \textit{biology} (1)
    \end{itemize}
\end{tabular}
}
{(modifié de~: \ccite{maslej_artificial_2025}[p.~87-92])}
{}
{
\begin{tablenotes}
\footnotesize
\item[a] {Il s'agit d'un sommaire de la colonne \emphase{Category}}
\item[b] {Le rapport parle de \emphase{catégories}, mais cela ressemble davantage à des étiquettes.}
\end{tablenotes}
}{tb}{Sommaire des étiquettes sur les techniques utilisées}
}

\enonce{AD50}{Assertion}{Plusieurs professionnels informatiques utilisent les GML.}{}{P046}
\explication{AD50}{de l'assertion}{Plusieurs professionnels informatiques utilisent les GML.}{
Un article publié en 2025 a analysé l'utilisation de ChatGPT par des professionnels informatiques \cite{li_unveiling_2025}. Pour ce faire, les auteurs analysent les données provenant du jeu de données DevChat et conservent uniquement les messages qui possèdent des liens URL vers GitHub. On en dénombre 2~547, répartis comme suit~: \ccite{li_unveiling_2025}[p.~18]~; 1~105 liens vers des fichiers de code source, 823 vers des \textit{Commits}, 295 vers des \textit{Issues}, 266 vers des \textit{Pull Requests}, et 58 vers des fils de discussions. 
Les cinq tâches les plus demandées (présentées en ordre décroissant) étaient~: 
\begin{itemize}
    \item \sfren{la génération et complétion de code source}{code generation and completion} (888); 
    \item \sfren{la modification et l'optimisation de code source}{code modification and optimization} (686);
    \item \sfren{la revue de code source et le débogage}{code review and debugging} (184); 
    \item \sfren{le déploiement et la configuration}{deployment and configuration} (123); 
    \item \sfren{l'explication de code source et la compréhension}{code explanation and understanding} (111).
\end{itemize}

Dans le \titleen{Artificial Intelligence Index Report} de 2024, il est présenté le sondage de Stack Overflow \footnote {Stack Overflow est un site Web proposant des questions et réponses dans le domaine de l'informatique.} portant sur la relation entre les professionnels informatiques et les instruments d'IA \ccite{maslej_artificial_2024}[p.~269]. Réalisé en 2023, ce sondage a recueilli les réponses de plus de quatre-vingt-dix mille professionnels informatiques.

Selon ce sondage, les deux instruments les plus utilisés par catégories sont~:
\begin{itemize}
    \item outils de développement~:
    \begin{itemize}
        \item GitHub Copilot (56\%); 
        \item Tabnine (12~\%);
    \end{itemize}
    \item outils de recherche~:
        \begin{itemize}
        \item ChatGPT (83~\%);
        \item Bing AI (19~\%).
    \end{itemize}
\end{itemize}
GitHub Copilot et ChatGPT utilisent des GML. 
\quoteenfr
{Copilot enables developers to choose between multiple large language models (LLMs), including OpenAI's GPT models, Anthropic's Claude, xAI's Grok, and Google's Gemini.}
{Copilot permet aux développeurs de choisir parmi plusieurs grands modèles de langage (GML), notamment les modèles GPT d'OpenAI, Claude d'Anthropic, Grok de xAI et Gemini de Google.}
{\footnote{Wikipedia, «~Github Copilot~», \url{https://en.wikipedia.org/wiki/GitHub_Copilot} (consulté le 2025-08-15).}}

Le figure~B.23 présente les réponses des professionnels informatiques sur les usages actuels, les usages désirés et les usages non désirés quant aux instruments d'IA. Comme il est indiqué, les professionnels informatiques s'en servent pour écrire du code (82~\%). Cependant, ils souhaiteraient s'en servir pour tester le code (55~\%) et documenter le code (50~\%). 
\figen
{}
{figure/chap1/documentation/IAuse.png}
{(source~: \textit{Figure}~4.4.15 \ccite{maslej_artificial_2024}[p.~270])}
{}{}{H}{Adoption d'outils d'IA pour des tâches de développement}
}{}

% Plusieurs avancent efficience, je n'ai pas vu de preuve qui allait à l'encontre.
\paragraphe{P047}{\fl{Plusieurs caractéristiques de qualité ne sont pas respectées ou démontrées comme telles par les grands modèles de langage.} D'abord, la cohérence ainsi que les ambiguïtés provenant de la langue et du contexte demeurent même pour les GML \cite{khurana_natural_2023, yang_natural_2023}. Les ellipses\footnote{Il s'agit de l'«~omission d’un ou plusieurs mots dans une phrase qui reste cependant compréhensible.~» \cite{Robert_ellipse_0}} en sont un exemple \cite{cavar_computing_2024}. Ces ambiguïtés peuvent conduire les GML à donner deux réponses contradictoires. 
\sep
Ensuite l'efficience n'est pas démontrée pour la rédaction des documents de spécifications des exigences. Pourtant, il s'agit d'un sujet plusieurs fois abordé dans l'utilisation des GML. Un revue de littérature trouve 22 articles qui abordent le sujet \cite{marques_using_2024}. Parmi ces articles, 15 abordent l'efficience, mais seulement quatre établissent un comparatif. L'accès à trois des quatre comparatifs confirme l'utilisation d'un échantillon de dix personnes ou moins \cite{kutzner_supporting_2023, waseem_chatgpt_2024, ronanki_investigating_2023}. 
\sep
Puis ce qui est de la réfutabilité, un article soulève les risques de contamination des jeux de données servant à l'entraînement. Cette situation aurait comme effet d'invalider des comparatifs entre les GML et de former des hypothèses trompeuses lors de l'étude des techniques de TAL \cite{sainz_nlp_2023}. 
\sep
Enfin, en ce qui concerne l'éthique, Laura Weidinger publie en 2022, avec 22 autres auteurs, l'article \titleen{Taxonomy of Risks posed by Language Models}{Taxonomie des risques posés par les modèles de langage}. Plusieurs des risques soulevés sont des enjeux éthiques, et les techniques de mitigation réduisent le risque, mais ne l'élimine pas entièrement \cite{weidinger_taxonomy_2022}.
Par ailleurs, plusieurs de ces risques se concrétiseront inévitablement à une fréquence difficile, voire impossible à évaluer (attendu la faible documentation des méthodes et des instruments). Dès lors il faudra avoir recours à des techniques de palliation qui sont très rarement abordées suite à l’exposé de la mitigation. Une analyse judicieuse des coûts-bénéfices des instruments d’IA devrait contenir les coûts de mitigation et de palliation, comme cela le serait pour tout projet \ccite{pmi_guidefr_2017}[p.~450].
}{AD51, AD52, AD53, AD54}{AD51}

\enonce{AD51}{Assertion}{Les ambiguïtés provenant de la langue et du contexte demeurent même pour les GML.}{}{P047}
\explication{AD51}{de l'assertion}{Les ambiguïtés provenant de la langue et du contexte demeurent même pour les GML.}{
Un article publié en 2022 fait un état de l'art sur le TAL \cite{khurana_natural_2023}. Les auteurs dressent quelques constats~:
\quoteenfrlist {
    \item[][---] Contextual words and phrases in the language where same words and phrases can have different meanings in a sentence which are easy for the humans to understand but makes a challenging task.
    \item[][---] Language containing informal phrases, expressions, idioms, and culture-specific lingo make difficult to design models intended for the broad use, \elide.
    \item[][---] \elide the models for NLP may be working good for an individual domain, geographic area but for a broad use such challenges need to be tackled.
} {
\item[][---] Dans une langue, les mots et expressions contextuels peuvent avoir des significations différentes au sein d'une même phrase. Si ces mots et expressions sont faciles à comprendre pour les humains, leur utilisation représente néanmoins un défi.
\item[][---] Les langues contenant des expressions familières, des tournures de phrase, des idiomes et un jargon culturel spécifique compliquent la conception de modèles destinés à un usage généralisé, \elide.
\item[][---] \elide les modèles de TAL peuvent être performants dans un domaine ou une zone géographique spécifique, mais pour une utilisation à grande échelle, ces difficultés doivent être surmontées
}{\cite{khurana_natural_2023}}

Un article publié en 2023 effectue une revue de littérature sur l'utilisation du TAL appliqué à la sécurité de l'aviation. \cite{yang_natural_2023}.  
\begin{itemize}
    \item L'article résume vingt articles dont la majorité avait comme objectifs d'identifier les incidents.
    \item Les auteurs des différents articles utilisent des rapports ou des bases de données standardisés, treize utilisent la base de données \textit{Aviation Safety Reporting System} (ASRS) qui provient de la NASA.
    \item Le langage utilisé est l'anglais dans dix-huit des articles.
    \item Les techniques de TAL utilisées peuvent être des combinaisons de~; \textit{SVM, LSA, LDA, TF-IDF, Long short-term memory (LSTM), Bag of Word (BoW), Word2Vec, Diffusion Maps (DM), Depth Neural Network (DNN)}, etc.
    \item Les techniques s'appliquent à la reconnaissance de texte et de la parole.
\end{itemize}
Le contexte revient comme un défi important~: 
\quoteenfr
{NLP struggles to understand the context and meaning of words and phrases, especially in aviation-specific jargon. For example, the term «~runway~» can refer to both the physical runway and the takeoff/landing procedures associated with it. \cite{yang_natural_2023}
}
{Le traitement automatique du langage naturel peine à comprendre le contexte et le sens des mots et des expressions, notamment dans le jargon spécifique à l'aviation. Par exemple, le terme « piste » peut désigner à la fois la piste physique et les procédures de décollage et d'atterrissage qui y sont associées.}

Un article publié en 2024 démontre que les ellipses sont particulièrement difficiles pour le TAL et que les GML n'y échappent pas~:
\quoteenfrlist {
    \item[][---] Ellipsis constructions are obviously still challenging for all the common SOTA [State-of-the-art] NLP pipelines and rule-based systems.
    \item[][---] The fact that specific models trained on the prediction of ellipses in sentences outperform LLMs seems to indicate that the lack of explicit data and pure self-supervised machine learning is not sufficient to handle opaque elements in language, either. Training LLMs on purely overt data ignores significant properties of language. Ellipsis phenomena are grammatical and systematic, and it seems problematic for current LLMs to guess covert continuations.
} {
\item[][---] Les constructions d'ellipses restent évidemment un défi pour tous les pipelines TAL SOTA courants et les systèmes basés sur des règles.
\item[][---] Le fait que des modèles spécifiques entraînés à prédire les ellipses dans les phrases surpassent les modèles de langage semble indiquer que l'absence de données explicites et l'apprentissage automatique autosupervisé pur ne suffisent pas non plus à traiter les éléments opaques du langage. L'entraînement des GML sur des données purement explicites ignore des propriétés importantes du langage. Les phénomènes d'ellipses sont grammaticaux et systématiques, et il semble problématique pour les GML actuels de deviner les suites implicites.
}{\cite{cavar_computing_2024}}
\vspace{-1\baselineskip}
}{}

\enonce{AD52}{Assertion}{Il n'a pas été encore démontré, dans un contexte réel, que les GML peuvent apporter une aide pour les documents de spécification des exigences.}{}{P047}
\explication{AD52}{de l'assertion}{Il n'a pas été encore démontré, dans un contexte réel, que les GML peuvent apporter une aide pour les documents de spécification des exigences.}{
Il y a plusieurs articles qui s'intéressent aux performances des GML pour les exigences en logiciel. En 2024, une revue de littérature obtient à 22 articles à analyser sur l'utilisation de ChatGPT et les exigences logicielles \cite{marques_using_2024}. Plusieurs articles s'intéressent aux performances des GML vis-à-vis de la spécification des exigences logicielles. En 2024, une revue de littérature recense 22 articles portant sur l'utilisation de ChatGPT dans ce domaine. Les thèmes abordés par l'analyse de ces articles sont les suivants~:
\begin{itemize}
\item \sfren{efficacité dans la génération des exigences}{Effectiveness in Generating Requirements} (15);
\item \sfren{rôle de l'intervention humaine}{Role of Human Input} (3);
\item \sfren{exploration de l'ingénierie rapide}{Exploration of Prompt Engineering} (6);
\item \sfren{défis et limites de l'IA}{AI Challenges and Limitations} (6);
\item \sfren{orientations futures de la recherche}{Future Research Directions} (15);
\item \sfren{expertise humaine comparée}{Comparative Human Expertise}  (4).
\end{itemize}
Parmi les articles faisant l'objet de la revue, seuls quatre présentaient un comparatif. L'accès à trois d'entre eux ne permet pas de conclure qu'il y a définitivement un bénéfice à utiliser ChatGPT, puisque les échantillons sont trop faibles  \cite{kutzner_supporting_2023, waseem_chatgpt_2024, ronanki_investigating_2023}~:
\begin{itemize}
    \item Les trois articles arrivent à la conclusion qu'il y aurait un bénéfice à utiliser ChatGPT pour documenter les exigences.
    \item Les échantillons pour chaque article respectivement, sont~: sept étudiants en informatique, dix étudiants en informatique, puis dix professionnels informatiques experts en exigence logiciel.
\end{itemize}
\vspace{-1\baselineskip}
}{}

\enonce{AD53}{Assertion}{Les contaminations des jeux de données servant à l'entraînement d'un modèle peuvent invalider les essais de celui-ci.}{}{P047}
\explication{AD53}{l'assertion}{Les contaminations des jeux de données servant à l'entraînement d'un modèle peuvent invalider les essais de celui-ci.}{
Les données d’entraînement peuvent contaminer. En 2023, un article met en garde sur la contamination \cite{sainz_nlp_2023} et indique les trois moments où la contamination peut avoir lieu~: 
\begin{itemize}
    \item contamination lors du pré-entraînement (\textit{Contamination during pre-training});
    \item contamination lors du réglage fin supervisé (\textit{Contamination on supervised fine-tuning});
    \item contamination après déploiement(\textit{Contamination after deployment}).
\end{itemize}
Les deux conséquences de la contamination sont également soulevées dans l'article~:
\quoteenfr
{The first one is that the performance of an LLM when evaluated on a benchmark it already processed during pre-training will be overestimated, causing it to be preferred with respect to other LLMs. This affects the comparative assessment of the quality of LLMs.
The second is that papers proposing scientific hypotheses on certain NLP tasks could be using contaminated LLMs, and thus make wrong claims about their hypotheses, and invalidate alternative hypotheses that could be true. This second consequence has an enormous negative impact on our field and is our main focus.}
{Premièrement, les performances d'un GML évalué sur un ensemble de données de référence qu'il a déjà traité lors du préentraînement seront surestimées, ce qui le favorisera par rapport à d'autres GML. Cela affecte l'évaluation comparative de la qualité des GML.
Deuxièmement, les articles proposant des hypothèses scientifiques sur certaines tâches de traitement automatique du langage naturel pourraient utiliser des GML contaminés, et ainsi formuler des affirmations erronées concernant leurs hypothèses et invalider des hypothèses alternatives qui pourraient être valides. Cette seconde conséquence a un impact négatif considérable sur notre domaine et constitue notre principal axe de recherche.}
{\cite{sainz_nlp_2023}}
\vspace{-1\baselineskip}
}{}

\enonce{AD54}{Assertion}{Les répercussions des techniques de TAL commencent à être formalisées et plusieurs sont néfastes.}{}{P047}
\explication{AD54}{de l'assertion}{Les répercussions des techniques de TAL commencent à être formalisées et plusieurs sont néfastes.}{
En 2020, une revue de littérature portant sur les biais dans le TAL est publiée \cite{blodgett_language_2020}. Au total, 146~articles y sont analysés. L'auteur affirme que plusieurs articles manquent d'arguments et de connaissances pour étayer leurs propos. Il formule alors trois recommandations visant à améliorer la recherche sur ce sujet~:
\quoteenfrlist {
    \item[](R1) Ground work analyzing "bias" in NLP systems in the relevant literature outside of NLP that explores the relationships between language and social hierarchies. Treat representational harms as harmful in their own right.
    \item[](R2) Provide explicit statements of why the system behaviors that are described as "bias" are harmful, in what ways, and to whom. Be forthright about the normative reasoning (Green, 2019) underlying these statements.
    \item[](R3) Examine language use in practice by engaging with the lived experiences of members of communities affected by NLP systems. Interrogate and reimagine the power relations between technologists and such communities.
} {
\item[](R1) Effectuer un travail de fond analysant les "biais"» dans les systèmes de TAL à partir de la littérature pertinente, en dehors du domaine du TAL, qui explore les relations entre le langage et les hiérarchies sociales. Considérer les préjudices représentationnels comme nuisibles en soi.
\item[](R2) Fournir des énoncés explicites expliquant pourquoi les comportements du système décrits comme des "biais"» sont nuisibles, de quelle manière et pour qui. Assumer clairement le raisonnement normatif (Green, 2019) sous-jacent à ces énoncés.
\item[](R3) Examiner l’usage du langage en pratique en s’appuyant sur les expériences vécues des membres des communautés affectées par les systèmes de TAL. Interroger et repenser les rapports de pouvoir entre les technologues et ces communautés.
}{\cite{blodgett_language_2020}}

En 2021, Laura Weidinger et 22 coauteurs publient un article sur les risques associés aux modèles de langage \cite{weidinger_ethical_2021}. Ils y présentent 21 risques préjudiciables, répartis selon des domaines de risque. Un article ultérieur approfondit encore cette analyse.

En 2022, Laura Weidinger et ses collègues de DeepMind proposent un regroupement plus structuré des risques dans un article intitulé \textit{Taxonomy of Risks posed by Language Models} \cite{weidinger_taxonomy_2022}. Leurs regroupements s'appuient sur 35 références, et les auteurs présentent également plusieurs techniques d'atténuation. Voici une partie du tableau sommaire qui se retrouve à l'intérieur de l'article (tableau~B.11). 
Les zones de risques, les mécanismes et les techniques de mitigation demeurent parfois diffus, mais il s'agit néanmoins d'un pas en avant dans la compréhension des risques. Certaines zones de risques ne sont d'ailleurs pas propres aux modèles de langage et peuvent s'appliquer à bien d'autres disciplines.
\begin{itemize}
    \item risques technologiques, en général~: la récupération d'une technologie à des fins malveillantes (4) ou les conséquences de son utilisation pour l'environnement et la société (6) ;
    \item risques en informatique~: la fuite d'information (2) et la désinformation préjudiciable (3) ;
    \item risques à simuler l'humain~: la discrimination, les discours haineux et l'exclusion (1) et les interactions personne-machine (5) ;
\end{itemize}
\tab
{Risques avec les modèles de langage.}
{figure/chap1/documentation/risk.png}
{[Traduction libre] (modifié de \cite{weidinger_taxonomy_2022})}
{}{
\begin{tablenotes}
\footnotesize
\item[a] {Il a été retiré les colonnes ; \textit{Type of Harm} et \textit{Evidence of this risk in large-scale LMs}}
\end{tablenotes}
}{H}
\vspace{-1\baselineskip}
}{}

\paragraphe{P048}{\fl{Les techniques d'IA ne remplacent pas les techniques informatiques.}
Un sondage mené en 2020 auprès de cent trente experts en méthodes formelles \ccite{ter_formal_2020}[p.~36] révèle que 93~\% d'entre eux estiment que l'IA ne remplacera probablement pas, ou ne remplacera définitivement pas, les méthodes formelles. Parmi les explications fournies, il y a~:
\quoteenfrlist {
\item[][---] \elide only formal methods can provide guarantees about correctnes \elide
\item[][---] \elide the crucial role of formal methods in requirements specification: "~at the end of the day, both the ambiguous setting and the mapping to the unambiguous setting are characteristic of human activities. I have a hard time imagining that these creative aspects can be fully automated~".
\item[][---]\elide we need approaches that combine model-driven and data-driven techniques \elide.
} {
\item[][---] \elide seules les méthodes formelles peuvent garantir l'exactitude \elide
\item[][---] \elide le rôle crucial des méthodes formelles dans la spécification des exigences : «~En fin de compte, tant le contexte ambigu que la traduction vers un contexte non ambigu sont caractéristiques des activités humaines. J'ai de la difficulté à imaginer que ces aspects créatifs puissent être entièrement automatisés.~»
\item[][---] \elide nous avons besoin d'approches combinant les techniques basées sur les modèles et les techniques basées sur les données \elide
}{\ccite{ter_formal_2020}[p.~36]}

De plus, les auteurs de l’ouvrage de référence \titleen{Artificial Intelligence A Modern Approach}{L'intelligence artificielle : une approche moderne} expliquent pourquoi, avec un titre semblable, un grand nombre des techniques informatiques y sont précisément décrites et que les modèles transformateurs \footnote{Le terme \emphase{modèle transformateur} signifie «~Architecture de réseau de neurones artificiels fondée sur le mécanisme d'attention \elide~.».\cite{Gdt_mt_0}} utilisés pour les GML ne sont pas suffisant~:
\quoteenfr
{It may be that transformer models and their relatives are learning latent representations that capture the same basic ideas as grammars and semantic information, or it may be that something entirely different is happening within these enormous models ; we simply don’t know.}
{Il se peut que les modèles transformateurs et leurs dérivés apprennent des représentations latentes qui capturent les mêmes idées fondamentales que les grammaires et l'information sémantique, ou bien il se peut que quelque chose de totalement différent se produise au sein de ces modèles gigantesques ; nous l'ignorons tout simplement.}
{\ccite{norvig_artificial_2021}[p.~1609]}
Durant toute la revue de littérature, le terme \emphase{modèle} a été utilisé pour décrire les résultats obtenus des différentes techniques de TAL. 
\sep
Néanmoins, l'usage de ce terme en IA désigne tout autre chose et cela est expliqué dans la prochaine section.
}{AD55}{AD55}

\enonce{AD55}{Assertion}{Les techniques d'IA ne remplacent pas les techniques informatiques.}{}{P048}
\explication{AD55}{de l'assertion}{Les techniques d'IA ne remplacent pas les techniques informatiques.}{
Dans \textit{Artificial Intelligence A Modern Approach} plusieurs chapitres sont consacrés aux différentes techniques permettant d'obtenir un modèle exprimé depuis un langage formel ou un langage non formel. Les auteurs prennent soin d'expliquer les raisons pour lesquelles toutes ces techniques sont réunies dans le même ouvrage~:
\quoteenfr
{The curious reader may wonder why we learned about grammars, parsing, and semantic interpretation in the previous chapter, only to discard those notions in favor of purely datadriven models in this chapter ? At present, the answer is simply that the data-driven models are easier to develop and maintain, and score better on standard benchmarks, compared to the hand-built systems that can be constructed using a reasonable amount of human effort with the approaches described in Chapter 23. It may be that transformer models and their relatives are learning latent representations that capture the same basic ideas as grammars and semantic information, or it may be that something entirely different is happening within these enormous models ; we simply don’t know. We do know that a system that is trained  with textual data is easier to maintain and to adapt to new domains and new natural languages than a system that relies on hand-crafted features.}
{Le lecteur curieux pourrait se demander pourquoi nous avons étudié les grammaires, l'analyse syntaxique et l'interprétation sémantique dans le chapitre précédent, pour ensuite les abandonner au profit de modèles purement axés sur les données dans ce chapitre. Pour l'instant, la réponse est simple~: les modèles axés sur les données sont plus faciles à développer et à maintenir, et obtiennent de meilleurs résultats aux tests de performance standard, que les systèmes construits manuellement, avec un effort humain raisonnable, grâce aux approches décrites au chapitre 23. Il se peut que les modèles transformateurs et leurs équivalents apprennent des représentations latentes qui capturent les mêmes idées fondamentales que les grammaires et les informations sémantiques, ou qu'il se passe quelque chose de complètement différent au sein de ces modèles gigantesques ; nous l'ignorons. Nous savons qu'un système entraîné avec des données textuelles est plus facile à maintenir et à adapter à de nouveaux domaines et à de nouvelles langues naturelles qu'un système qui repose sur des caractéristiques conçues manuellement.}
{\cite{norvig_artificial_2021}}

En 2020, un article ayant pour titre \titleen{The 2020 Expert Survey on Formal Methods}{Enquête auprès d'experts de 2020 sur les méthodes formelles}, interroge 130 experts en méthodes formelles \ccite{ter_formal_2020}[p.~36]. Le tableau~B.12 présente les résultats obtenus à la question portant sur l'éventualité que les méthodes formelles soient remplacées~:
\taben
{\textit{Whether other rising approaches}}
{figure/chap1/documentation/Qf.png}
{(source~: \cite{ter_formal_2020})}
{}{}{tb}{Avis d'experts en méthodes formelles sur un éventuel remplacement par des méthodes non formelles}

Le sondage révèle également que la majorité des experts ne voit pas l'IA remplacer les méthodes formelles. Quelques-uns des commentaires de ces experts sont décrits ci-dessous~:
\quoteenfrlist
{
\item [][---] Many of them argue that only formal methods can provide guarantees about correctness, e.g. «~artificial intelligence is wrong in 10 to 25\% of the cases and must be hand-tuned. What is formal there?~» and «~artificial intelligence will need to be certified. What methods will be used to certify it?~». 
\item [][---] \elide two comments that praise the crucial role of formal methods in requirements specification: «~at the end of the day, both the ambiguous setting and the mapping to the unambiguous setting are characteristic of human activities. I have a hard time imagining that these creative aspects can be fully automated~». 
\item [][---] Other comments see fruitful interactions between both fields. Another comment adds: «~we need approaches that combine model-driven and data-driven techniques}
{
\item [][---] Nombre d'entre eux soutiennent que seules les méthodes formelles peuvent garantir l'exactitude des résultats, par exemple~: «~L'intelligence artificielle se trompe dans 10 à 25~\% des cas et doit être paramétrée manuellement. En quoi les méthodes formelles sont-elles pertinentes dans ce cas ?~» et «~L'intelligence artificielle devra être certifiée. Quelles méthodes seront utilisées pour cette certification ?~». 
\item [][---] \elide
Deux commentaires, soulignant le rôle crucial des méthodes formelles dans la spécification des exigences, avancent une raison bien différente~: «~En fin de compte, l’ambiguïté du contexte et sa traduction en un contexte non ambigu sont caractéristiques de l’activité humaine. J’ai du mal à imaginer que ces aspects créatifs puissent être entièrement automatisés.~» 
\item [][---] D’autres commentaires entrevoient des interactions fructueuses entre les deux domaines. \elide Un autre commentaire ajoute~: "Nous avons besoin d’approches combinant techniques basées sur les modèles et techniques basées sur les données.}
{\cite{ter_formal_2020}}
\vspace{-1\baselineskip}
}{}
